%============================================================
%  Bd. 30 - Halbverbände und Verbände
%============================================================

\documentclass[main.tex]{subfiles}

\ifSubfilesClassLoaded{
  \BandSetupStandalone{B30}
}{
}

\title{Bd. 30 - Halbverbände und Verbände}
\author{Martin Kunze}
\date{}
\setcounter{file}{30}

\hypersetup{
  pdftitle={Bd. 30 - Halbverbände und Verbände},
  pdfauthor={Martin Kunze}
}

\begin{document}

\title{Bd. 30 - Halbverbände und Verbände}
\author{Martin Kunze}
\date{}
\hypersetup{pageanchor=false}
\maketitle
\hypersetup{pageanchor=true}
\setcounter{tocdepth}{3}
\tableofcontents
\setcounter{file}{30}
\renewcommand{\FormulaBandID}{30}
\newcommand{\ProofRefStack}[1]{%
  \ensuremath{\begin{gathered}#1\end{gathered}}}

\chapter{Einleitung}

Ausgehend von den kommutativen idempotenten Halbgruppen aus Band~21 und der allgemeinen Ordnungstheorie
in Band~11 werden Halbverbände und Verbände algebraisch aufgebaut. Ein
Halbverband ist dabei eine kommutative idempotente Halbgruppe. Ein eigenes
Kapitel trennt
die allgemeine Theorie der Supremums-Halbverbände mit Null und der
beschränkten Supremums-Halbverbände vollständig von den Anwendungen auf
Frankls Vermutung.

\chapter{Algebraische Halbverbände}

Band~11 entwickelt partielle Ordnungen sowie obere und untere Schranken,
Infima und Suprema. Der vorliegende Band geht den umgekehrten Weg: Aus den
Gesetzen einer binären Operation wird zunächst eine partielle Ordnung
gewonnen.

Die Bezeichnung \emph{Halbgitter} wird synonym zu \emph{Halbverband}
verwendet. In der allgemeinen Axiomatik schreiben wir \(\Semi{A}\), solange
die zugrunde liegende Operation aus dem Kontext mitgeführt und in den
abstrakten Termen durch Juxtaposition dargestellt wird. Die Juxtaposition
\(xy\) bezeichnet dann das Ergebnis dieser Operation; sie ist keine
gewöhnliche Multiplikation. Sobald der Operator benannt wird, gehört er
obligatorisch zur Strukturbezeichnung. Für den Operator \(\NeutralSemiOp\)
schreiben wir also
\[
  \Semi{A,\NeutralSemiOp}
  \qquad\text{und}\qquad
  x\NeutralSemiOp{}y.
\]
Allgemein steht der Träger an erster und der Operator an zweiter Stelle.

\section{Allgemeine axiomatische Definition}

Wie in den Bänden~03 und~04 werden die Strukturaxiome zunächst einzeln
registriert und anschließend in einer Definition gesammelt. In allen vier
Axiomen bezeichnet die Juxtaposition dieselbe, kontextuell festgelegte
Operation.

\FormulaAxiomDeltaK[Abgeschlossenheit einer Halbverbandsoperation]{%
  x,y\in A\vdash x\OmittedSemiOp{}y\in A%
}{SemilatticeClosureLaw}{%
  \DeltaRow{Mengen}{A\dsep x\dsep y}
}
\par\smallskip

\FormulaAxiomDeltaK[Assoziativität einer Halbverbandsoperation]{%
  x,y,z\in A\vdash
  \bigl(x\OmittedSemiOp{}y\bigr)\OmittedSemiOp{}z
    =x\OmittedSemiOp\bigl(y\OmittedSemiOp{}z\bigr)%
}{SemilatticeAssociativityLaw}{%
  \DeltaRow{Mengen}{A\dsep x\dsep y\dsep z}
}
\par\smallskip

\FormulaAxiomDeltaK[Kommutativität einer Halbverbandsoperation]{%
  x,y\in A\vdash x\OmittedSemiOp{}y=y\OmittedSemiOp{}x%
}{SemilatticeCommutativityLaw}{%
  \DeltaRow{Mengen}{A\dsep x\dsep y}
}
\par\smallskip

\FormulaAxiomDeltaK[Idempotenz einer Halbverbandsoperation]{%
  x\in A\vdash x\OmittedSemiOp{}x=x%
}{SemilatticeIdempotenceLaw}{%
  \DeltaRow{Mengen}{A\dsep x}
}

\FormulaDefDeltaK[Begriff des Halbverbandes]{%
  \Semi{A}%
}{Semilattice}{%
  \DeltaRow{Mengen}{A\dsep x\dsep y\dsep z}
  \DeltaRow{\textbf{Axiome}}{}
  \DeltaRow{Abgeschlossenheit}{%
    x,y\in A\vdash x\OmittedSemiOp{}y\in A}
    [\FormulaRefAuto{SemilatticeClosureLaw}[ax]]
  \DeltaRow{Assoziativität}{%
    \begin{aligned}[t]
      &x,y,z\in A\vdash{}\\
      &\bigl(x\OmittedSemiOp{}y\bigr)\OmittedSemiOp{}z
        =x\OmittedSemiOp\bigl(y\OmittedSemiOp{}z\bigr)
    \end{aligned}}
    [\FormulaRefAuto{SemilatticeAssociativityLaw}[ax]]
  \DeltaRow{Kommutativität}{%
    x,y\in A\vdash x\OmittedSemiOp{}y=y\OmittedSemiOp{}x}
    [\FormulaRefAuto{SemilatticeCommutativityLaw}[ax]]
  \DeltaRow{Idempotenz}{%
    x\in A\vdash x\OmittedSemiOp{}x=x}
    [\FormulaRefAuto{SemilatticeIdempotenceLaw}[ax]]
  \DeltaRow{\textbf{Neues Symbol}}{}
  \DeltaPrem{Kontextuelle Schreibweise}{\Semi{A}}
}

Die folgenden Zugriffsaxiome verbinden das Strukturprädikat mit den vier
soeben registrierten Gesetzen. Wird die Operation später benannt, wird die
Juxtaposition überall einheitlich durch den gewählten Operator ersetzt.

\FormulaAxiomDeltaK[Zugriff auf die Abgeschlossenheit eines Halbverbandes]{%
  \Semi{A}\dsep x,y\in A
  \vdash x\OmittedSemiOp{}y\in A%
}{SemilatticeClosureAxiom}{%
  \DeltaRow{Mengen}{A\dsep x\dsep y}
}
\par\smallskip

\FormulaAxiomDeltaK[Zugriff auf die Assoziativität eines Halbverbandes]{%
  \Semi{A}\dsep x,y,z\in A\vdash
  \bigl(x\OmittedSemiOp{}y\bigr)\OmittedSemiOp{}z
    =x\OmittedSemiOp\bigl(y\OmittedSemiOp{}z\bigr)%
}{SemilatticeAssociativityAxiom}{%
  \DeltaRow{Mengen}{A\dsep x\dsep y\dsep z}
}
\par\smallskip

\FormulaAxiomDeltaK[Zugriff auf die Kommutativität eines Halbverbandes]{%
  \Semi{A}\dsep x,y\in A
  \vdash x\OmittedSemiOp{}y=y\OmittedSemiOp{}x%
}{SemilatticeCommutativityAxiom}{%
  \DeltaRow{Mengen}{A\dsep x\dsep y}
}
\par\smallskip

\FormulaAxiomDeltaK[Zugriff auf die Idempotenz eines Halbverbandes]{%
  \Semi{A}\dsep x\in A
  \vdash x\OmittedSemiOp{}x=x%
}{SemilatticeIdempotenceAxiom}{%
  \DeltaRow{Mengen}{A\dsep x}
}

Die Abgeschlossenheit ist die auf den Träger eingeschränkte Form der in
Band~05 eingeführten Funktionsauswertung. Beim Nachweis eines konkreten
Halbverbandes werden alle vier Eigenschaften einzeln gezeigt; nur die letzte
Beweiszeile verwendet die zusammenfassende Definition.

\section{Homomorphismen und Isomorphismen von Halbverbänden}

Ein Homomorphismus muss die jeweils benannten Operationen erhalten. Deshalb
werden sowohl Quell- als auch Zieloperator in der Bezeichnung mitgeführt.

\FormulaDefDeltaK[Begriff des Halbverbandshomomorphismus]{%
  \SemilatticeHom{f;A,\NeutralSemiOp;B,\mathbin{\diamond}}%
}{SemilatticeHomomorphism}{%
  \DeltaRow{Mengen}{A\dsep B\dsep x\dsep y}
  \DeltaRow{Funktionen}{f\dsep\NeutralSemiOp\dsep\mathbin{\diamond}}
  \DeltaPrem{Quellhalbverband}{\Semi{A,\NeutralSemiOp}}
  \DeltaPrem{Zielhalbverband}{\Semi{B,\mathbin{\diamond}}}
  \DeltaPrem{Abbildung}{f\colon A\to B}
  \DeltaRow{Operationserhaltung}{%
    x,y\in A\vdash
    f(x\NeutralSemiOp{}y)=f(x)\mathbin{\diamond}f(y)}
  \DeltaRow{\textbf{Neues Symbol}}{}
  \DeltaPrem{Halbverbandshomomorphismen}{%
    \SemilatticeHom{f;A,\NeutralSemiOp;B,\mathbin{\diamond}}}
}
\par\smallskip

\FormulaAxiomDeltaK[Zugriff auf den Quellhalbverband eines Halbverbandshomomorphismus]{%
  \SemilatticeHom{f;A,\NeutralSemiOp;B,\mathbin{\diamond}}
  \vdash\Semi{A,\NeutralSemiOp}%
}{SemilatticeHomSourceAxiom}{%
  \DeltaRow{Mengen}{A\dsep B}
  \DeltaRow{Funktionen}{f\dsep\NeutralSemiOp\dsep\mathbin{\diamond}}
}
\par\smallskip

\FormulaAxiomDeltaK[Zugriff auf den Zielhalbverband eines Halbverbandshomomorphismus]{%
  \SemilatticeHom{f;A,\NeutralSemiOp;B,\mathbin{\diamond}}
  \vdash\Semi{B,\mathbin{\diamond}}%
}{SemilatticeHomTargetAxiom}{%
  \DeltaRow{Mengen}{A\dsep B}
  \DeltaRow{Funktionen}{f\dsep\NeutralSemiOp\dsep\mathbin{\diamond}}
}
\par\smallskip

\FormulaAxiomDeltaK[Zugriff auf die Abbildung eines Halbverbandshomomorphismus]{%
  \SemilatticeHom{f;A,\NeutralSemiOp;B,\mathbin{\diamond}}
  \vdash f\colon A\to B%
}{SemilatticeHomMapAxiom}{%
  \DeltaRow{Mengen}{A\dsep B}
  \DeltaRow{Funktionen}{f\dsep\NeutralSemiOp\dsep\mathbin{\diamond}}
}
\par\smallskip

\FormulaAxiomDeltaK[Auswertung eines Halbverbandshomomorphismus]{%
  \begin{aligned}[t]
    &\SemilatticeHom{f;A,\NeutralSemiOp;B,\mathbin{\diamond}}
      \dsep x,y\in A\vdash{}\\
    &f(x\NeutralSemiOp{}y)=f(x)\mathbin{\diamond}f(y)
  \end{aligned}
}{SemilatticeHomEvaluationAxiom}{%
  \DeltaRow{Mengen}{A\dsep B\dsep x\dsep y}
  \DeltaRow{Funktionen}{f\dsep\NeutralSemiOp\dsep\mathbin{\diamond}}
}

\FormulaDefDeltaK[Begriff des Halbverbandsisomorphismus]{%
  \SemilatticeIso{f;A,\NeutralSemiOp;B,\mathbin{\diamond}}%
}{SemilatticeIsomorphism}{%
  \DeltaRow{Mengen}{A\dsep B}
  \DeltaRow{Funktionen}{f\dsep\NeutralSemiOp\dsep\mathbin{\diamond}}
  \DeltaPrem{Halbverbandshomomorphismus}{%
    \SemilatticeHom{f;A,\NeutralSemiOp;B,\mathbin{\diamond}}}
  \DeltaPrem{Bijektivität}{f\colon A\bij B}
  \DeltaRow{\textbf{Neues Symbol}}{}
  \DeltaPrem{Halbverbandsisomorphismen}{%
    \SemilatticeIso{f;A,\NeutralSemiOp;B,\mathbin{\diamond}}}
}
\par\smallskip

\FormulaAxiomDeltaK[Ein Halbverbandsisomorphismus ist ein Homomorphismus]{%
  \SemilatticeIso{f;A,\NeutralSemiOp;B,\mathbin{\diamond}}
  \vdash
  \SemilatticeHom{f;A,\NeutralSemiOp;B,\mathbin{\diamond}}%
}{SemilatticeIsoHomAxiom}{%
  \DeltaRow{Mengen}{A\dsep B}
  \DeltaRow{Funktionen}{f\dsep\NeutralSemiOp\dsep\mathbin{\diamond}}
}
\par\smallskip

\FormulaAxiomDeltaK[Ein Halbverbandsisomorphismus ist bijektiv]{%
  \SemilatticeIso{f;A,\NeutralSemiOp;B,\mathbin{\diamond}}
  \vdash f\colon A\bij B%
}{SemilatticeIsoBijectiveAxiom}{%
  \DeltaRow{Mengen}{A\dsep B}
  \DeltaRow{Funktionen}{f\dsep\NeutralSemiOp\dsep\mathbin{\diamond}}
}

\FormulaThmDeltaK[Die Identität erhält die Halbverbandsoperation]{%
  \Semi{A,\NeutralSemiOp}\dsep x,y\in A\vdash
  \Id_A(x\NeutralSemiOp{}y)
    =\Id_A(x)\NeutralSemiOp\Id_A(y)
}{SemilatticeIdentityPreservesOperation}{%
  \DeltaRow{Mengen}{A\dsep x\dsep y}
  \DeltaRow{Funktionen}{\Id_A\dsep\NeutralSemiOp}
}
\begin{tabproofwide}
  \proofstepwidestar[1]{\Semi{A,\NeutralSemiOp}}{\rA}
  \proofstepwidestar[2]{x\in A}{\rA}
  \proofstepwidestar[3]{y\in A}{\rA}
  \proofstepwidestar[1,2,3]{x\NeutralSemiOp{}y\in A}{%
    \FormulaRefAuto{SemilatticeClosureAxiom}[ax]{1,2,3}}
  \proofstepwidestar[1,2,3]{%
    \Id_A(x\NeutralSemiOp{}y)=x\NeutralSemiOp{}y}{%
    \FormulaRefAuto{x\in A\vdash \Id_A(x)=x}{4}}
  \proofstepwidestar[2]{\Id_A(x)=x}{%
    \FormulaRefAuto{x\in A\vdash \Id_A(x)=x}{2}}
  \proofstepwidestar[3]{\Id_A(y)=y}{%
    \FormulaRefAuto{x\in A\vdash \Id_A(x)=x}{3}}
  \proofstepwidestar[2,3]{%
    \Id_A(x)\NeutralSemiOp\Id_A(y)=x\NeutralSemiOp{}y}{%
    \rIE{6,\rIE{7,\rII}}}
  \proofstepwidestar[1,2,3]{%
    \Id_A(x\NeutralSemiOp{}y)
      =\Id_A(x)\NeutralSemiOp\Id_A(y)}{%
    \FormulaRefAuto{a=b,\,c=b\vdash a=c}{5,8}}
\end{tabproofwide}

\FormulaThmDeltaK[Die Identität ist ein Halbverbandshomomorphismus]{%
  \Semi{A,\NeutralSemiOp}\vdash
  \SemilatticeHom{\Id_A;A,\NeutralSemiOp;A,\NeutralSemiOp}%
}{SemilatticeIdentityHomomorphism}{%
  \DeltaRow{Mengen}{A\dsep x\dsep y}
  \DeltaRow{Funktionen}{\Id_A\dsep\NeutralSemiOp}
}
\begin{tabproof}
  \proofstep{1}{\Semi{A,\NeutralSemiOp}}{\rA}
  \proofstep{}{\Id_A\colon A\to A}{\FormulaRefAuto{IdA}}
  \proofstep{3}{x\in A}{\rA}
  \proofstep{4}{y\in A}{\rA}
  \proofstep{1,3,4}{%
    \Id_A(x\NeutralSemiOp{}y)
      =\Id_A(x)\NeutralSemiOp\Id_A(y)}{%
    \FormulaRefAuto{SemilatticeIdentityPreservesOperation}[thm]{1,3,4}}
  \proofstep{1}{%
    \SemilatticeHom{\Id_A;A,\NeutralSemiOp;A,\NeutralSemiOp}}{%
    \FormulaRefAuto{SemilatticeHomomorphism}[def]{1,1,2,5}}
\end{tabproof}

\FormulaThmDeltaK[Komposition von Halbverbandshomomorphismen]{%
  \begin{aligned}[t]
    &\SemilatticeHom{f;A,\NeutralSemiOp;B,\mathbin{\diamond}}
      \dsep
      \SemilatticeHom{g;B,\mathbin{\diamond};C,\mathbin{\bullet}}
      \vdash{}\\
    &\SemilatticeHom{g\circ f;A,\NeutralSemiOp;C,\mathbin{\bullet}}
  \end{aligned}
}{SemilatticeHomComposition}{%
  \DeltaRow{Mengen}{A\dsep B\dsep C\dsep x\dsep y}
  \DeltaRow{Funktionen}{%
    f\dsep g\dsep\NeutralSemiOp\dsep\mathbin{\diamond}\dsep\mathbin{\bullet}}
}
\begin{tabproofwide}
  \proofstepwidestar[1]{%
    \SemilatticeHom{f;A,\NeutralSemiOp;B,\mathbin{\diamond}}}{\rA}
  \proofstepwidestar[2]{%
    \SemilatticeHom{g;B,\mathbin{\diamond};C,\mathbin{\bullet}}}{\rA}
  \proofstepwidestar[1]{\Semi{A,\NeutralSemiOp}}{%
    \FormulaRefAuto{SemilatticeHomSourceAxiom}[ax]{1}}
  \proofstepwidestar[2]{\Semi{C,\mathbin{\bullet}}}{%
    \FormulaRefAuto{SemilatticeHomTargetAxiom}[ax]{2}}
  \proofstepwidestar[1]{f\colon A\to B}{%
    \FormulaRefAuto{SemilatticeHomMapAxiom}[ax]{1}}
  \proofstepwidestar[2]{g\colon B\to C}{%
    \FormulaRefAuto{SemilatticeHomMapAxiom}[ax]{2}}
  \proofstepwidestar[1,2]{g\circ f\colon A\to C}{%
    \FormulaRefAuto{G\circ F\colon A\to C}[thm]{5,6}}
  \proofstepwidestar[8]{x\in A}{\rA}
  \proofstepwidestar[9]{y\in A}{\rA}
  \proofstepwidestar[1,8,9]{x\NeutralSemiOp{}y\in A}{%
    \FormulaRefAuto{SemilatticeClosureAxiom}[ax]{3,8,9}}
  \proofstepwidestar[1,8]{f(x)\in B}{%
    \FormulaRefAuto{F\colon A\to B\dsep x\in A\vdash F(x)\in B}[thm]{5,8}}
  \proofstepwidestar[1,9]{f(y)\in B}{%
    \FormulaRefAuto{F\colon A\to B\dsep x\in A\vdash F(x)\in B}[thm]{5,9}}
  \proofstepwidestar[1,8,9]{%
    f(x\NeutralSemiOp{}y)=f(x)\mathbin{\diamond}f(y)}{%
    \FormulaRefAuto{SemilatticeHomEvaluationAxiom}[ax]{1,8,9}}
  \proofstepwidestar[1,2,8]{g(f(x))=(g\circ f)(x)}{%
    \FormulaRefAuto{x\in A\vdash F(G(x))=(F\circ G)(x)}[thm]{5,6,8}}
  \proofstepwidestar[1,2,9]{g(f(y))=(g\circ f)(y)}{%
    \FormulaRefAuto{x\in A\vdash F(G(x))=(F\circ G)(x)}[thm]{5,6,9}}
  \proofstepwide[1,2,8,9]{(g\circ f)(x\NeutralSemiOp{}y)}{=}{%
    g\bigl(f(x\NeutralSemiOp{}y)\bigr)}{%
    \FormulaRefAuto{x\in A\vdash(F\circ G)(x)=F(G(x))}[thm]{5,6,10}}
  \proofstepwide[1,2,8,9]{}{=}{%
    g\bigl(f(x)\mathbin{\diamond}f(y)\bigr)}{\rIE{13,\rII}}
  \proofstepwide[1,2,8,9]{}{=}{%
    g(f(x))\mathbin{\bullet}g(f(y))}{%
    \FormulaRefAuto{SemilatticeHomEvaluationAxiom}[ax]{2,11,12}}
  \proofstepwide[1,2,8,9]{}{=}{%
    (g\circ f)(x)\mathbin{\bullet}(g\circ f)(y)}{%
    \rIE{14,\rIE{15,\rII}}}
  \proofstepwide[1,2,8,9]{(g\circ f)(x\NeutralSemiOp{}y)}{=}{%
    (g\circ f)(x)\mathbin{\bullet}(g\circ f)(y)}{%
    \rChain{16,17,18,19}}
  \proofstepwidestar[1,2]{%
    \SemilatticeHom{g\circ f;A,\NeutralSemiOp;C,\mathbin{\bullet}}}{%
    \FormulaRefAuto{SemilatticeHomomorphism}[def]{3,4,7,20}}
\end{tabproofwide}

\FormulaThmDeltaK[Die Identität ist ein Halbverbandsisomorphismus]{%
  \Semi{A,\NeutralSemiOp}\vdash
  \SemilatticeIso{\Id_A;A,\NeutralSemiOp;A,\NeutralSemiOp}%
}{SemilatticeIdentityIsomorphism}{%
  \DeltaRow{Mengen}{A}
  \DeltaRow{Funktionen}{\Id_A\dsep\NeutralSemiOp}
}
\begin{tabproof}
  \proofstep{1}{\Semi{A,\NeutralSemiOp}}{\rA}
  \proofstep{1}{%
    \SemilatticeHom{\Id_A;A,\NeutralSemiOp;A,\NeutralSemiOp}}{%
    \FormulaRefAuto{SemilatticeIdentityHomomorphism}[thm]{1}}
  \proofstep{}{\Id_A\colon A\bij A}{%
    \FormulaRefAuto{\Id_A\colon A\bij A}}
  \proofstep{1}{%
    \SemilatticeIso{\Id_A;A,\NeutralSemiOp;A,\NeutralSemiOp}}{%
    \FormulaRefAuto{SemilatticeIsomorphism}[def]{2,3}}
\end{tabproof}

\FormulaThmDeltaK[Komposition von Halbverbandsisomorphismen]{%
  \begin{aligned}[t]
    &\SemilatticeIso{f;A,\NeutralSemiOp;B,\mathbin{\diamond}}
      \dsep
      \SemilatticeIso{g;B,\mathbin{\diamond};C,\mathbin{\bullet}}
      \vdash{}\\
    &\SemilatticeIso{g\circ f;A,\NeutralSemiOp;C,\mathbin{\bullet}}
  \end{aligned}
}{SemilatticeIsoComposition}{%
  \DeltaRow{Mengen}{A\dsep B\dsep C}
  \DeltaRow{Funktionen}{%
    f\dsep g\dsep\NeutralSemiOp\dsep\mathbin{\diamond}\dsep\mathbin{\bullet}}
}
\begin{tabproofwide}
  \proofstepwidestar[1]{%
    \SemilatticeIso{f;A,\NeutralSemiOp;B,\mathbin{\diamond}}}{\rA}
  \proofstepwidestar[2]{%
    \SemilatticeIso{g;B,\mathbin{\diamond};C,\mathbin{\bullet}}}{\rA}
  \proofstepwidestar[1]{%
    \SemilatticeHom{f;A,\NeutralSemiOp;B,\mathbin{\diamond}}}{%
    \FormulaRefAuto{SemilatticeIsoHomAxiom}[ax]{1}}
  \proofstepwidestar[2]{%
    \SemilatticeHom{g;B,\mathbin{\diamond};C,\mathbin{\bullet}}}{%
    \FormulaRefAuto{SemilatticeIsoHomAxiom}[ax]{2}}
  \proofstepwidestar[1,2]{%
    \SemilatticeHom{g\circ f;A,\NeutralSemiOp;C,\mathbin{\bullet}}}{%
    \FormulaRefAuto{SemilatticeHomComposition}[thm]{3,4}}
  \proofstepwidestar[1]{f\colon A\bij B}{%
    \FormulaRefAuto{SemilatticeIsoBijectiveAxiom}[ax]{1}}
  \proofstepwidestar[2]{g\colon B\bij C}{%
    \FormulaRefAuto{SemilatticeIsoBijectiveAxiom}[ax]{2}}
  \proofstepwidestar[1,2]{g\circ f\colon A\bij C}{%
    \FormulaRefAuto{BijectiveFunctionComposition}[thm]{6,7}}
  \proofstepwidestar[1,2]{%
    \SemilatticeIso{g\circ f;A,\NeutralSemiOp;C,\mathbin{\bullet}}}{%
    \FormulaRefAuto{SemilatticeIsomorphism}[def]{5,8}}
\end{tabproofwide}

\section{Der supremal induzierte Relationsoperator}

\subsection{Axiomatische Definition}

Auch in diesem Abschnitt bleibt die Halbverbandsoperation unbenannt. Ihre
Abhängigkeit vom Träger wird durch den Kontext \(\Semi{A}\) festgehalten; für
den induzierten Relationsoperator verwenden wir unmittelbar die natürliche
Schreibweise \(\leq\).

\FormulaAxiomDeltaK[Trägerbedingung des supremal induzierten Relationsoperators]{%
  \Semi{A}\dsep x\InducedLeq y\vdash x,y\in A
}{JoinOrderCarrierAxiom}{%
  \DeltaRow{Mengen}{A\dsep x\dsep y}
  \DeltaPrem{\makecell[l]{Supremal aus\\
                          \(\Semi{A}\) induzierter\\
                          Relationsoperator}}{\InducedLeq}
}

\FormulaAxiomDeltaK[Auswertung des supremal induzierten Relationsoperators]{%
  \Semi{A}\dsep x,y\in A\vdash
  \bigl(x\InducedLeq y\leftrightarrow x\OmittedSemiOp{}y=y\bigr)%
}{JoinOrderEvaluationAxiom}{%
  \DeltaRow{Mengen}{A\dsep x\dsep y}
  \DeltaPrem{Induzierter Relationsoperator}{\InducedLeq}
}

\FormulaDefDeltaK[Supremal induzierter Relationsoperator]{%
  x\InducedLeq y%
}{JoinOrder}{%
  \DeltaRow{Mengen}{A\dsep x\dsep y}
  \DeltaRow{\textbf{Axiome}}{}
  \DeltaRow{Trägerbedingung}{%
    \Semi{A}\dsep x\InducedLeq y\vdash x,y\in A}
    [\FormulaRefAuto{JoinOrderCarrierAxiom}[ax]]
  \DeltaRow{Auswertung}{%
    \begin{aligned}[t]
      &\Semi{A}\dsep x,y\in A\vdash{}\\[-2pt]
      &\bigl(x\InducedLeq y\leftrightarrow x\OmittedSemiOp{}y=y\bigr)
    \end{aligned}}
    [\FormulaRefAuto{JoinOrderEvaluationAxiom}[ax]]
  \DeltaRow{\textbf{Neues Symbol}}{}
  \DeltaPrem{Induzierter Relationsoperator}{\InducedLeq}
}

\begin{remark}
Die induzierte Ordnung wird ausschließlich in der natürlichen Infixform
\(x\InducedLeq y\) verwendet.
\end{remark}

\subsection{Träger und Auswertung}

\FormulaThmDeltaK[Erste Komponente der supremal induzierten Relation]{%
  \Semi{A}\dsep x\InducedLeq y\vdash x\in A
}{JoinOrderFirstCarrier}{%
  \DeltaRow{Mengen}{A\dsep x\dsep y}
  \DeltaPrem{Von \(\Semi{A}\) supremal induzierter Relationsoperator}{\InducedLeq}
}
\begin{tabproof}
  \proofstep{1}{\Semi{A}}{\rA}
  \proofstep{2}{x\InducedLeq y}{\rA}
  \proofstep{1,2}{x\in A\land y\in A}{%
    \FormulaRefAuto{JoinOrderCarrierAxiom}[ax]{1,2}}
  \proofstep{1,2}{x\in A}{\rAEa{3}}
\end{tabproof}

\FormulaThmDeltaK[Zweite Komponente der supremal induzierten Relation]{%
  \Semi{A}\dsep x\InducedLeq y\vdash y\in A
}{JoinOrderSecondCarrier}{%
  \DeltaRow{Mengen}{A\dsep x\dsep y}
  \DeltaPrem{Von \(\Semi{A}\) supremal induzierter Relationsoperator}{\InducedLeq}
}
\begin{tabproof}
  \proofstep{1}{\Semi{A}}{\rA}
  \proofstep{2}{x\InducedLeq y}{\rA}
  \proofstep{1,2}{x\in A\land y\in A}{%
    \FormulaRefAuto{JoinOrderCarrierAxiom}[ax]{1,2}}
  \proofstep{1,2}{y\in A}{\rAEb{3}}
\end{tabproof}

\FormulaThmDeltaK[Auswertung der supremal induzierten Relation]{%
  \Semi{A}\dsep x\InducedLeq y
  \vdash x\OmittedSemiOp{}y=y
}{JoinOrderEvaluation}{%
  \DeltaRow{Mengen}{A\dsep x\dsep y}
  \DeltaPrem{Von \(\Semi{A}\) supremal induzierter Relationsoperator}{\InducedLeq}
}
\begin{tabproof}
  \proofstep{1}{\Semi{A}}{\rA}
  \proofstep{2}{x\InducedLeq y}{\rA}
  \proofstep{1,2}{x\in A}{\FormulaRefAuto{JoinOrderFirstCarrier}[thm]{1,2}}
  \proofstep{1,2}{y\in A}{\FormulaRefAuto{JoinOrderSecondCarrier}[thm]{1,2}}
  \proofstep{1,2}{x\InducedLeq y\leftrightarrow x\OmittedSemiOp{}y=y}{%
    \FormulaRefAuto{JoinOrderEvaluationAxiom}[ax]{1,3,4}}
  \proofstep{1,2}{x\OmittedSemiOp{}y=y}{%
    \FormulaRefAuto{P\leftrightarrow Q, P\vdash Q}{5,2}}
\end{tabproof}

\subsection{Die induzierte partielle Ordnung}

\FormulaThmDeltaK[Reflexivität der supremal induzierten Relation]{%
  \Semi{A}\dsep x\in A\vdash x\InducedLeq x
}{JoinOrderReflexive}{%
  \DeltaRow{Mengen}{A\dsep x}
  \DeltaPrem{Von \(\Semi{A}\) supremal induzierter Relationsoperator}{\InducedLeq}
}
\begin{tabproof}
  \proofstep{1}{\Semi{A}}{\rA}
  \proofstep{2}{x\in A}{\rA}
  \proofstep{1,2}{x\OmittedSemiOp{}x=x}{%
    \FormulaRefAuto{SemilatticeIdempotenceAxiom}[ax]{1,2}}
  \proofstep{1,2}{x\InducedLeq x\leftrightarrow x\OmittedSemiOp{}x=x}{%
    \FormulaRefAuto{JoinOrderEvaluationAxiom}[ax]{1,2,2}}
  \proofstep{1,2}{x\InducedLeq x}{%
    \FormulaRefAuto{P\leftrightarrow Q, Q\vdash P}{4,3}}
\end{tabproof}

\FormulaThmDeltaK[Transitivität der supremal induzierten Relation]{%
  \Semi{A}\dsep x\InducedLeq y
  \dsep y\InducedLeq z\vdash x\InducedLeq z
}{JoinOrderTransitive}{%
  \DeltaRow{Mengen}{A\dsep x\dsep y\dsep z}
  \DeltaPrem{Von \(\Semi{A}\) supremal induzierter Relationsoperator}{\InducedLeq}
}
\begin{tabproofwide}
  \proofstepwidestar[1]{\Semi{A}}{\rA}
  \proofstepwidestar[2]{x\InducedLeq y}{\rA}
  \proofstepwidestar[3]{y\InducedLeq z}{\rA}
  \proofstepwidestar[1,2]{x\in A}{%
    \FormulaRefAuto{JoinOrderFirstCarrier}[thm]{1,2}}
  \proofstepwidestar[1,2]{y\in A}{%
    \FormulaRefAuto{JoinOrderSecondCarrier}[thm]{1,2}}
  \proofstepwidestar[1,3]{z\in A}{%
    \FormulaRefAuto{JoinOrderSecondCarrier}[thm]{1,3}}
  \proofstepwidestar[1,2]{x\OmittedSemiOp{}y=y}{%
    \FormulaRefAuto{JoinOrderEvaluation}[thm]{1,2}}
  \proofstepwidestar[1,3]{y\OmittedSemiOp{}z=z}{%
    \FormulaRefAuto{JoinOrderEvaluation}[thm]{1,3}}
  \proofstepwidestar[1,2,3]{%
    \bigl(x\OmittedSemiOp{}y\bigr)\OmittedSemiOp{}z
      =x\OmittedSemiOp\bigl(y\OmittedSemiOp{}z\bigr)}{%
    \FormulaRefAuto{SemilatticeAssociativityAxiom}[ax]{1,4,5,6}}

  \proofstepwide[1,2,3]{x\OmittedSemiOp{}z}{=}{%
    x\OmittedSemiOp\bigl(y\OmittedSemiOp{}z\bigr)}{%
    \FormulaRefAuto{a=b\vdash b=a}{\rIE{8,\rII}}}
  \proofstepwide[1,2,3]{}{=}{%
    \bigl(x\OmittedSemiOp{}y\bigr)\OmittedSemiOp{}z}{%
    \FormulaRefAuto{a=b\vdash b=a}{9}}
  \proofstepwide[1,2,3]{}{=}{y\OmittedSemiOp{}z}{\rIE{7}}
  \proofstepwide[1,2,3]{}{=}{z}{\rIE{8}}
  \proofstepwidestar[1,2,3]{x\OmittedSemiOp{}z=z}{%
    \rChain{10,11,12,13}}

  \proofstepwidestar[1,2,3]{%
    x\InducedLeq z\leftrightarrow x\OmittedSemiOp{}z=z}{%
    \FormulaRefAuto{JoinOrderEvaluationAxiom}[ax]{1,4,6}}
  \proofstepwidestar[1,2,3]{x\InducedLeq z}{%
    \FormulaRefAuto{P\leftrightarrow Q, Q\vdash P}{15,14}}
\end{tabproofwide}

\FormulaThmDeltaK[Antisymmetrie der supremal induzierten Relation]{%
  \Semi{A}\dsep x\InducedLeq y
  \dsep y\InducedLeq x\vdash x=y
}{JoinOrderAntisymmetric}{%
  \DeltaRow{Mengen}{A\dsep x\dsep y}
  \DeltaPrem{Von \(\Semi{A}\) supremal induzierter Relationsoperator}{\InducedLeq}
}
\begin{tabproof}
  \proofstep{1}{\Semi{A}}{\rA}
  \proofstep{2}{x\InducedLeq y}{\rA}
  \proofstep{3}{y\InducedLeq x}{\rA}
  \proofstep{1,2}{x\in A}{\FormulaRefAuto{JoinOrderFirstCarrier}[thm]{1,2}}
  \proofstep{1,2}{y\in A}{\FormulaRefAuto{JoinOrderSecondCarrier}[thm]{1,2}}
  \proofstep{1,2}{x\OmittedSemiOp{}y=y}{\FormulaRefAuto{JoinOrderEvaluation}[thm]{1,2}}
  \proofstep{1,3}{y\OmittedSemiOp{}x=x}{\FormulaRefAuto{JoinOrderEvaluation}[thm]{1,3}}
  \proofstep{1,2}{x\OmittedSemiOp{}y=y\OmittedSemiOp{}x}{%
    \FormulaRefAuto{SemilatticeCommutativityAxiom}[ax]{1,4,5}}
  \proofstep{1,2,3}{x\OmittedSemiOp{}y=x}{\rIE{8,7}}
  \proofstep{1,2,3}{y=x}{\FormulaRefAuto{a=b,\,a=c\vdash b=c}{6,9}}
  \proofstep{1,2,3}{x=y}{\FormulaRefAuto{a=b\vdash b=a}{10}}
\end{tabproof}

\FormulaThmDeltaK[Ein supremal gelesener Halbverband induziert eine partielle Ordnung]{%
  \Semi{A}\vdash \PartOrd{A,\InducedLeq}
}{JoinSemilatticeInducesOrder}{%
  \DeltaRow{Mengen}{A\dsep x\dsep y\dsep z}
  \DeltaPrem{Von \(\Semi{A}\) supremal induzierter Relationsoperator}{\InducedLeq}
}
\begin{tabproofwide}
  \proofstepwidestar[1]{\Semi{A}}{\rA}
  \proofstepwidestar[2]{x\in A}{\rA}
  \proofstepwidestar[1,2]{x\InducedLeq x}{%
    \FormulaRefAuto{JoinOrderReflexive}[thm]{1,2}}
  \proofstepwidestar[4]{x\InducedLeq y}{\rA}
  \proofstepwidestar[5]{y\InducedLeq z}{\rA}
  \proofstepwidestar[1,4,5]{x\InducedLeq z}{%
    \FormulaRefAuto{JoinOrderTransitive}[thm]{1,4,5}}
  \proofstepwidestar[7]{y\InducedLeq x}{\rA}
  \proofstepwidestar[1,4,7]{x=y}{%
    \FormulaRefAuto{JoinOrderAntisymmetric}[thm]{1,4,7}}
  \proofstepwidestar[1]{\PartOrd{A,\InducedLeq}}{%
    \FormulaRefAuto{Partielle Ordnung}[def]{3,6,8}}
\end{tabproofwide}

\FormulaThmDeltaK[Charakterisierung des Hauptfilters im Halbverband]{%
  \begin{aligned}[t]
    &\Semi{A,\JoinOp}\dsep j,x\in A\vdash{}\\
    &x\in\PrincipalFilter{A,\InducedLeq}{j}
      \leftrightarrow j\JoinOp{}x=x
  \end{aligned}
}{JoinPrincipalFilterChar}{%
  \DeltaRow{Mengen}{A\dsep j\dsep x}
  \DeltaRow{Funktionen}{\JoinOp}
}
\begin{tabproofwide}
  \proofstepwidestar[1]{\Semi{A,\JoinOp}}{\rA}
  \proofstepwidestar[2]{j\in A}{\rA}
  \proofstepwidestar[3]{x\in A}{\rA}
  \proofstepwide[2,3]{x\in\PrincipalFilter{A,\InducedLeq}{j}}{\leftrightarrow}{%
    j\InducedLeq x}{\FormulaRefAuto{PrincipalFilter}[def]}
  \proofstepwide[2,3]{}{\leftrightarrow}{j\JoinOp{}x=x}{%
    \FormulaRefAuto{JoinOrderEvaluationAxiom}[ax]{1,2,3}}
\end{tabproofwide}

\subsection{Universelle Eigenschaft der supremalen Lesart}

Band~11 definiert das Supremum bereits für jede Teilmenge \(T\subseteq A\).
Für ein Elementpaar verwenden wir daher ohne neue Definition den Spezialfall
\(\sup_{A,\leq}(\{x,y\})\). Die folgenden algebraischen Sätze zeigen, dass
dieser Term durch \(x\OmittedSemiOp{}y\) dargestellt wird.

\FormulaThmDeltaK[Der erste Operand liegt unter dem algebraischen Supremum]{%
  \Semi{A}\dsep x,y\in A
  \vdash x\InducedLeq x\OmittedSemiOp{}y
}{JoinUpperLeft}{%
  \DeltaRow{Mengen}{A\dsep x\dsep y}
  \DeltaPrem{Von \(\Semi{A}\) supremal induzierter Relationsoperator}{\InducedLeq}
}
\begin{tabproof}
  \proofstep{1}{\Semi{A}}{\rA}
  \proofstep{2}{x\in A}{\rA}
  \proofstep{3}{y\in A}{\rA}
  \proofstep{1,2,3}{x\OmittedSemiOp{}y\in A}{%
    \FormulaRefAuto{SemilatticeClosureAxiom}[ax]{1,2,3}}
  \proofstep{1,2}{x\OmittedSemiOp{}x=x}{%
    \FormulaRefAuto{SemilatticeIdempotenceAxiom}[ax]{1,2}}
  \proofstep{1,2,3}{%
    \bigl(x\OmittedSemiOp{}x\bigr)\OmittedSemiOp{}y=x\OmittedSemiOp{}y}{\rIE{5,\rII}}
  \proofstep{1,2,3}{%
    \bigl(x\OmittedSemiOp{}x\bigr)\OmittedSemiOp{}y
      =x\OmittedSemiOp\bigl(x\OmittedSemiOp{}y\bigr)}{%
    \FormulaRefAuto{SemilatticeAssociativityAxiom}[ax]{1,2,2,3}}
  \proofstep{1,2,3}{%
    x\OmittedSemiOp\bigl(x\OmittedSemiOp{}y\bigr)=x\OmittedSemiOp{}y}{%
    \FormulaRefAuto{a=b,\,a=c\vdash b=c}{7,6}}
  \proofstep{1,2,3}{x\InducedLeq x\OmittedSemiOp{}y
    \leftrightarrow
      x\OmittedSemiOp\bigl(x\OmittedSemiOp{}y\bigr)=x\OmittedSemiOp{}y}{%
    \FormulaRefAuto{JoinOrderEvaluationAxiom}[ax]{1,2,4}}
  \proofstep{1,2,3}{x\InducedLeq x\OmittedSemiOp{}y}{%
    \FormulaRefAuto{P\leftrightarrow Q, Q\vdash P}{9,8}}
\end{tabproof}

\FormulaThmDeltaK[Der zweite Operand liegt unter dem algebraischen Supremum]{%
  \Semi{A}\dsep x,y\in A
  \vdash y\InducedLeq x\OmittedSemiOp{}y
}{JoinUpperRight}{%
  \DeltaRow{Mengen}{A\dsep x\dsep y}
  \DeltaPrem{Von \(\Semi{A}\) supremal induzierter Relationsoperator}{\InducedLeq}
}
\begin{tabproof}
  \proofstep{1}{\Semi{A}}{\rA}
  \proofstep{2}{x\in A}{\rA}
  \proofstep{3}{y\in A}{\rA}
  \proofstep{1,2,3}{y\InducedLeq y\OmittedSemiOp{}x}{%
    \FormulaRefAuto{JoinUpperLeft}[thm]{1,3,2}}
  \proofstep{1,2,3}{y\OmittedSemiOp{}x=x\OmittedSemiOp{}y}{%
    \FormulaRefAuto{SemilatticeCommutativityAxiom}[ax]{1,3,2}}
  \proofstep{1,2,3}{y\InducedLeq x\OmittedSemiOp{}y}{\rIE{5,4}}
\end{tabproof}

\FormulaThmDeltaK[Das algebraische Supremum ist die kleinste obere Schranke]{%
  \begin{aligned}[t]
    &\Semi{A}\dsep x,y,u\in A\dsep
      x\InducedLeq u\dsep y\InducedLeq u\vdash{}\\
    &x\OmittedSemiOp{}y\InducedLeq u
  \end{aligned}
}{JoinLeastUpper}{%
  \DeltaRow{Mengen}{A\dsep x\dsep y\dsep u}
  \DeltaPrem{Von \(\Semi{A}\) supremal induzierter Relationsoperator}{\InducedLeq}
}
\begin{tabproofwide}
  \proofstepwidestar[1]{\Semi{A}}{\rA}
  \proofstepwidestar[2]{x\in A}{\rA}
  \proofstepwidestar[3]{y\in A}{\rA}
  \proofstepwidestar[4]{u\in A}{\rA}
  \proofstepwidestar[5]{x\InducedLeq u}{\rA}
  \proofstepwidestar[6]{y\InducedLeq u}{\rA}
  \proofstepwidestar[1,5]{x\OmittedSemiOp{}u=u}{%
    \FormulaRefAuto{JoinOrderEvaluation}[thm]{1,5}}
  \proofstepwidestar[1,6]{y\OmittedSemiOp{}u=u}{%
    \FormulaRefAuto{JoinOrderEvaluation}[thm]{1,6}}
  \proofstepwidestar[1,2,3]{x\OmittedSemiOp{}y\in A}{%
    \FormulaRefAuto{SemilatticeClosureAxiom}[ax]{1,2,3}}

  \proofstepwide[1,2,3,4]{%
    \bigl(x\OmittedSemiOp{}y\bigr)\OmittedSemiOp{}u}{=}{%
    x\OmittedSemiOp\bigl(y\OmittedSemiOp{}u\bigr)}{%
    \FormulaRefAuto{SemilatticeAssociativityAxiom}[ax]{1,2,3,4}}
  \proofstepwide[1,2,3,4,6]{}{=}{x\OmittedSemiOp{}u}{\rIE{8}}
  \proofstepwide[1,2,3,4,5,6]{}{=}{u}{\rIE{7}}
  \proofstepwidestar[1,2,3,4,5,6]{%
    \bigl(x\OmittedSemiOp{}y\bigr)\OmittedSemiOp{}u=u}{%
    \rChain{10,11,12}}

  \proofstepwidestar[1,2,3,4]{x\OmittedSemiOp{}y\InducedLeq u
    \leftrightarrow
      \bigl(x\OmittedSemiOp{}y\bigr)\OmittedSemiOp{}u=u}{%
    \FormulaRefAuto{JoinOrderEvaluationAxiom}[ax]{1,9,4}}
  \proofstepwidestar[1,2,3,4,5,6]{x\OmittedSemiOp{}y\InducedLeq u}{%
    \FormulaRefAuto{P\leftrightarrow Q, Q\vdash P}{14,13}}
\end{tabproofwide}

\FormulaThmDeltaK[Die Halbverbandsoperation liefert das Paarsupremum]{%
  \Semi{A}\dsep x,y\in A
  \vdash x\OmittedSemiOp{}y=\sup_{A,\InducedLeq}(\{x,y\})
}{JoinIsPairSupremum}{%
  \DeltaRow{Mengen}{A\dsep x\dsep y\dsep u\dsep z}
  \DeltaPrem{Von \(\Semi{A}\) supremal induzierter Relationsoperator}{\InducedLeq}
}
\begin{tabproofwide}
  \proofstepwidestar[1]{\Semi{A}}{\rA}
  \proofstepwidestar[2]{x\in A}{\rA}
  \proofstepwidestar[3]{y\in A}{\rA}
  \proofstepwidestar[1]{\PartOrd{A,\InducedLeq}}{%
    \FormulaRefAuto{JoinSemilatticeInducesOrder}[thm]{1}}
  \proofstepwidestar[2,3]{\{x,y\}\subseteq A}{%
    \FormulaRefAuto{a \in A,\, b \in A \vdash \{a,b\} \subseteq A}{2,3}}
  \proofstepwidestar[1,2,3]{x\OmittedSemiOp{}y\in A}{%
    \FormulaRefAuto{SemilatticeClosureAxiom}[ax]{1,2,3}}
  \proofstepwidestar[1,2,3]{x\InducedLeq x\OmittedSemiOp{}y}{%
    \FormulaRefAuto{JoinUpperLeft}[thm]{1,2,3}}
  \proofstepwidestar[1,2,3]{y\InducedLeq x\OmittedSemiOp{}y}{%
    \FormulaRefAuto{JoinUpperRight}[thm]{1,2,3}}
  \proofstepwidestar[1,2,3]{%
    \begin{aligned}[t]
      &x\OmittedSemiOp{}y\in\UB_{A,\InducedLeq}(\{x,y\})\\[-2pt]
      &\quad\leftrightarrow
        x\OmittedSemiOp{}y\in A\land
        x\InducedLeq x\OmittedSemiOp{}y\land
        y\InducedLeq x\OmittedSemiOp{}y
    \end{aligned}}{%
    \FormulaRefAuto{PairUpperBoundSetCharacterization}[thm]{4,2,3}}
  \proofstepwidestar[1,2,3]{%
    x\OmittedSemiOp{}y\in\UB_{A,\InducedLeq}(\{x,y\})}{%
    \FormulaRefAuto{P\leftrightarrow Q, Q\vdash P}{%
      9,\rAI{6,\rAI{7,8}}}}
  \proofstepwidestar[11]{u\in\UB_{A,\InducedLeq}(\{x,y\})}{\rA}
  \proofstepwidestar[2,3]{%
    \begin{aligned}[t]
      u\in\UB_{A,\InducedLeq}(\{x,y\})
      \leftrightarrow
      \bigl(u\in A\land x\InducedLeq u\land y\InducedLeq u\bigr)
    \end{aligned}}{%
    \FormulaRefAuto{PairUpperBoundSetCharacterization}[thm]{4,2,3}}
  \proofstepwidestar[2,3,11]{%
    u\in A\land x\InducedLeq u\land y\InducedLeq u}{%
    \FormulaRefAuto{P\leftrightarrow Q, P\vdash Q}{12,11}}
  \proofstepwidestar[2,3,11]{u\in A}{\rAEa{13}}
  \proofstepwidestar[2,3,11]{x\InducedLeq u}{\rAEa{\rAEb{13}}}
  \proofstepwidestar[2,3,11]{y\InducedLeq u}{\rAEb{\rAEb{13}}}
  \proofstepwidestar[1,2,3,11]{x\OmittedSemiOp{}y\InducedLeq u}{%
    \FormulaRefAuto{JoinLeastUpper}[thm]{1,2,3,14,15,16}}
  \proofstepwidestar[1,2,3]{%
    \forall u\in\UB_{A,\InducedLeq}(\{x,y\})\,
      (x\OmittedSemiOp{}y\InducedLeq u)}{%
    \rUI{\rRI{11,17}}}
  \proofstepwidestar[1,2,3]{%
    x\OmittedSemiOp{}y=\sup_{A,\InducedLeq}(\{x,y\})}{%
    \FormulaRefAuto{SupremumCandidateEqualsSupremum}[thm]{4,5,10,18}}
\end{tabproofwide}

\section{Die ordnungstheoretische Charakterisierung}

Ziel dieses Abschnitts ist die Umkehrung der bisherigen Konstruktion: Aus
einer partiellen Ordnung, in der jedes Elementpaar ein Supremum besitzt, wird
die dadurch eindeutig bestimmte Halbverbandsoperation zurückgewonnen. Die
partielle Ordnung auf \(A\) bezeichnen wir mit \(\leq\), die aus den
Paarsuprema gewonnene Operation mit \(\NeutralSemiOp\). Da der Operator hier
ausdrücklich benannt ist, schreiben wir
\(\PairJoinOp{A,\NeutralSemiOp}\), also
\(\mathsf{SupOp}(A,\NeutralSemiOp)\), und verwenden ihn in allen Termen in
Infixnotation.

\FormulaAxiomDeltaK[Existenzaxiom einer Paarsupremumsoperation]{%
  \begin{aligned}[t]
    &\PairJoinOp{A,\NeutralSemiOp}\dsep x,y\in A\vdash{}\\
    &\exists m\in\UB_{A,\leq}(\{x,y\})\,
      \forall u\in\UB_{A,\leq}(\{x,y\})\,(m\leq u)
  \end{aligned}
}{PairJoinExistenceAxiom}{%
  \DeltaRow{Mengen}{A\dsep x\dsep y\dsep m\dsep u}
  \DeltaRow{Relationsoperatoren}{\leq}
  \DeltaRow{Funktionen}{\NeutralSemiOp}
}
\par\smallskip

\FormulaAxiomDeltaK[Abgeschlossenheitsaxiom einer Paarsupremumsoperation]{%
  \PairJoinOp{A,\NeutralSemiOp}\dsep x,y\in A
  \vdash x\NeutralSemiOp{}y\in A
}{PairJoinClosureAxiom}{%
  \DeltaRow{Mengen}{A\dsep x\dsep y}
  \DeltaRow{Relationsoperatoren}{\leq}
  \DeltaRow{Funktionen}{\NeutralSemiOp}
}
\par\smallskip

\FormulaAxiomDeltaK[Auswertungsaxiom einer Paarsupremumsoperation]{%
  \PairJoinOp{A,\NeutralSemiOp}\dsep x,y\in A\vdash
  x\NeutralSemiOp{}y=\sup_{A,\leq}(\{x,y\})
}{PairJoinEvaluationAxiom}{%
  \DeltaRow{Mengen}{A\dsep x\dsep y}
  \DeltaRow{Relationsoperatoren}{\leq}
  \DeltaRow{Funktionen}{\NeutralSemiOp}
}

\FormulaDefDeltaK[Durch Paarsuprema bestimmte Operation]{%
  \PairJoinOp{A,\NeutralSemiOp}%
}{PairJoinOperation}{%
  \DeltaRow{Mengen}{A\dsep x\dsep y\dsep m\dsep u}
  \DeltaRow{Relationsoperatoren}{\leq}
  \DeltaRow{Funktionen}{\NeutralSemiOp}
  \DeltaRow{\textbf{Axiome}}{}
  \DeltaRow{\makecell[l]{Existenz der\\Paarsuprema}}{%
    \begin{aligned}[t]
      &x,y\in A\vdash{}\\[-2pt]
      &\exists m\in\UB_{A,\leq}(\{x,y\})\,\\[-2pt]
      &\quad\forall u\in\UB_{A,\leq}(\{x,y\})\,(m\leq u)
    \end{aligned}}
    [\FormulaRefAuto{PairJoinExistenceAxiom}[ax]]
  \DeltaRow{Abgeschlossenheit}{%
    x,y\in A\vdash x\NeutralSemiOp{}y\in A}
    [\FormulaRefAuto{PairJoinClosureAxiom}[ax]]
  \DeltaRow{Auswertung}{%
    x,y\in A\vdash
      x\NeutralSemiOp{}y=\sup_{A,\leq}(\{x,y\})}
    [\FormulaRefAuto{PairJoinEvaluationAxiom}[ax]]
  \DeltaRow{\textbf{Neues Symbol}}{}
  \DeltaPrem{Paarsupremumsoperator}{%
    \PairJoinOp{A,\NeutralSemiOp}}
}

\FormulaThmDeltaK[Erster Operand einer Paarsupremumsoperation]{%
  \begin{aligned}[t]
    &\PartOrd{A,\leq}\dsep\PairJoinOp{A,\NeutralSemiOp}
      \dsep x,y\in A\vdash{}\\
    &x\leq x\NeutralSemiOp{}y
  \end{aligned}
}{PairJoinUpperLeft}{%
  \DeltaRow{Mengen}{A\dsep x\dsep y\dsep m\dsep u}
  \DeltaRow{Relationsoperatoren}{\leq}
  \DeltaRow{Funktionen}{\NeutralSemiOp}
}
\begin{tabproof}
  \proofstep{1}{\PartOrd{A,\leq}}{\rA}
  \proofstep{2}{\PairJoinOp{A,\NeutralSemiOp}}{\rA}
  \proofstep{3}{x\in A}{\rA}
  \proofstep{4}{y\in A}{\rA}
  \proofstep{2,3,4}{%
    \exists m\in\UB_{A,\leq}(\{x,y\})\,
      \forall u\in\UB_{A,\leq}(\{x,y\})\,(m\leq u)}{%
    \FormulaRefAuto{PairJoinExistenceAxiom}[ax]{2,3,4}}
  \proofstep{1,2,3,4}{x\leq\sup_{A,\leq}(\{x,y\})}{%
    \FormulaRefAuto{PairSupremumFirstOperand}[thm]{1,3,4,5}}
  \proofstep{2,3,4}{x\NeutralSemiOp{}y=\sup_{A,\leq}(\{x,y\})}{%
    \FormulaRefAuto{PairJoinEvaluationAxiom}[ax]{2,3,4}}
  \proofstep{1,2,3,4}{x\leq x\NeutralSemiOp{}y}{\rIE{7,6}}
\end{tabproof}

\FormulaThmDeltaK[Zweiter Operand einer Paarsupremumsoperation]{%
  \PartOrd{A,\leq}\dsep\PairJoinOp{A,\NeutralSemiOp}\dsep x,y\in A
  \vdash y\leq x\NeutralSemiOp{}y
}{PairJoinUpperRight}{%
  \DeltaRow{Mengen}{A\dsep x\dsep y}
  \DeltaRow{Relationsoperatoren}{\leq}
  \DeltaRow{Funktionen}{\NeutralSemiOp}
}
\begin{tabproof}
  \proofstep{1}{\PartOrd{A,\leq}}{\rA}
  \proofstep{2}{\PairJoinOp{A,\NeutralSemiOp}}{\rA}
  \proofstep{3}{x\in A}{\rA}
  \proofstep{4}{y\in A}{\rA}
  \proofstep{2,3,4}{%
    \exists m\in\UB_{A,\leq}(\{x,y\})\,
      \forall u\in\UB_{A,\leq}(\{x,y\})\,(m\leq u)}{%
    \FormulaRefAuto{PairJoinExistenceAxiom}[ax]{2,3,4}}
  \proofstep{1,2,3,4}{y\leq\sup_{A,\leq}(\{x,y\})}{%
    \FormulaRefAuto{PairSupremumSecondOperand}[thm]{1,3,4,5}}
  \proofstep{2,3,4}{x\NeutralSemiOp{}y=\sup_{A,\leq}(\{x,y\})}{%
    \FormulaRefAuto{PairJoinEvaluationAxiom}[ax]{2,3,4}}
  \proofstep{1,2,3,4}{y\leq x\NeutralSemiOp{}y}{\rIE{7,6}}
\end{tabproof}

\FormulaThmDeltaK[Kleinste obere Schranke einer Paarsupremumsoperation]{%
  \begin{aligned}[t]
    &\PartOrd{A,\leq}\dsep\PairJoinOp{A,\NeutralSemiOp}
      \dsep x,y,u\in A\dsep x\leq u\dsep y\leq u\vdash{}\\
    &x\NeutralSemiOp{}y\leq u
  \end{aligned}
}{PairJoinLeastUpper}{%
  \DeltaRow{Mengen}{A\dsep x\dsep y\dsep u\dsep z\dsep m}
  \DeltaRow{Relationsoperatoren}{\leq}
  \DeltaRow{Funktionen}{\NeutralSemiOp}
}
\begin{tabproofwide}
  \proofstepwidestar[1]{\PartOrd{A,\leq}}{\rA}
  \proofstepwidestar[2]{\PairJoinOp{A,\NeutralSemiOp}}{\rA}
  \proofstepwidestar[3]{x\in A}{\rA}
  \proofstepwidestar[4]{y\in A}{\rA}
  \proofstepwidestar[5]{u\in A}{\rA}
  \proofstepwidestar[6]{x\leq u}{\rA}
  \proofstepwidestar[7]{y\leq u}{\rA}
  \proofstepwidestar[2,3,4]{%
    \begin{aligned}[t]
      &\exists s\in\UB_{A,\leq}(\{x,y\})\,{}\\[-2pt]
      &\quad\forall v\in\UB_{A,\leq}(\{x,y\})\,(s\leq v)
    \end{aligned}}{%
    \FormulaRefAuto{PairJoinExistenceAxiom}[ax]{2,3,4}}
  \proofstepwidestar[1,2,3,4,5,6,7]{%
    \sup_{A,\leq}(\{x,y\})\leq u}{%
    \FormulaRefAuto{PairSupremumUniversal}[thm]{%
      1,3,4,5,8,6,7}}
  \proofstepwidestar[2,3,4]{%
    x\NeutralSemiOp{}y=\sup_{A,\leq}(\{x,y\})}{%
    \FormulaRefAuto{PairJoinEvaluationAxiom}[ax]{2,3,4}}
  \proofstepwidestar[1,2,3,4,5,6,7]{%
    x\NeutralSemiOp{}y\leq u}{\rIE{10,9}}
\end{tabproofwide}

\FormulaThmDeltaK[Idempotenz einer Paarsupremumsoperation]{%
  \PartOrd{A,\leq}\dsep \PairJoinOp{A,\NeutralSemiOp}\dsep x\in A
  \vdash x\NeutralSemiOp{}x=x
}{PairJoinIdempotent}{%
  \DeltaRow{Mengen}{A\dsep x}
  \DeltaRow{Relationsoperatoren}{\leq}
  \DeltaRow{Funktionen}{\NeutralSemiOp}
}
\begin{tabproof}
  \proofstep{1}{\PartOrd{A,\leq}}{\rA}
  \proofstep{2}{\PairJoinOp{A,\NeutralSemiOp}}{\rA}
  \proofstep{3}{x\in A}{\rA}
  \proofstep{2,3}{%
    x\NeutralSemiOp{}x=\sup_{A,\leq}(\{x,x\})}{%
    \FormulaRefAuto{PairJoinEvaluationAxiom}[ax]{2,3,3}}
  \proofstep{1,3}{\sup_{A,\leq}(\{x,x\})=x}{%
    \FormulaRefAuto{PairSupremumIdempotent}[thm]{1,3}}
  \proofstep{1,2,3}{x\NeutralSemiOp{}x=x}{%
    \rChain{4,5}}
\end{tabproof}

\FormulaThmDeltaK[Kommutativität einer Paarsupremumsoperation]{%
  \PartOrd{A,\leq}\dsep \PairJoinOp{A,\NeutralSemiOp}\dsep x,y\in A
  \vdash x\NeutralSemiOp{}y=y\NeutralSemiOp{}x
}{PairJoinCommutative}{%
  \DeltaRow{Mengen}{A\dsep x\dsep y}
  \DeltaRow{Relationsoperatoren}{\leq}
  \DeltaRow{Funktionen}{\NeutralSemiOp}
}
\begin{tabproof}
  \proofstep{1}{\PartOrd{A,\leq}}{\rA}
  \proofstep{2}{\PairJoinOp{A,\NeutralSemiOp}}{\rA}
  \proofstep{3}{x\in A}{\rA}
  \proofstep{4}{y\in A}{\rA}
  \proofstep{2,3,4}{%
    \exists m\in\UB_{A,\leq}(\{x,y\})\,
      \forall u\in\UB_{A,\leq}(\{x,y\})\,(m\leq u)}{%
    \FormulaRefAuto{PairJoinExistenceAxiom}[ax]{2,3,4}}
  \proofstep{2,3,4}{x\NeutralSemiOp{}y=\sup_{A,\leq}(\{x,y\})}{%
    \FormulaRefAuto{PairJoinEvaluationAxiom}[ax]{2,3,4}}
  \proofstep{2,3,4}{y\NeutralSemiOp{}x=\sup_{A,\leq}(\{y,x\})}{%
    \FormulaRefAuto{PairJoinEvaluationAxiom}[ax]{2,4,3}}
  \proofstep{1,2,3,4}{%
    \sup_{A,\leq}(\{x,y\})=\sup_{A,\leq}(\{y,x\})}{%
    \FormulaRefAuto{PairSupremumCommutative}[thm]{1,3,4,5}}
  \proofstep{1,2,3,4}{x\NeutralSemiOp{}y=y\NeutralSemiOp{}x}{%
    \FormulaRefAuto{a=b,\,c=b\vdash a=c}{6,\rIE{8,7}}}
\end{tabproof}

\FormulaThmDeltaK[Assoziativität einer Paarsupremumsoperation]{%
  \begin{aligned}[t]
    &\PartOrd{A,\leq}\dsep \PairJoinOp{A,\NeutralSemiOp}
      \dsep x,y,z\in A\vdash{}\\
    &\bigl(x\NeutralSemiOp{}y\bigr)\NeutralSemiOp{}z
      =x\NeutralSemiOp\bigl(y\NeutralSemiOp{}z\bigr)
  \end{aligned}
}{PairJoinAssociative}{%
  \DeltaRow{Mengen}{A\dsep x\dsep y\dsep z}
  \DeltaRow{Relationsoperatoren}{\leq}
  \DeltaRow{Funktionen}{\NeutralSemiOp}
}
\begin{tabproofwide}
  \proofstepwidestar[1]{\PartOrd{A,\leq}}{\rA}
  \proofstepwidestar[2]{\PairJoinOp{A,\NeutralSemiOp}}{\rA}
  \proofstepwidestar[3]{x\in A}{\rA}
  \proofstepwidestar[4]{y\in A}{\rA}
  \proofstepwidestar[5]{z\in A}{\rA}
  \proofstepwidestar[2,3,4]{x\NeutralSemiOp{}y\in A}{%
    \FormulaRefAuto{PairJoinClosureAxiom}[ax]{2,3,4}}
  \proofstepwidestar[2,4,5]{y\NeutralSemiOp{}z\in A}{%
    \FormulaRefAuto{PairJoinClosureAxiom}[ax]{2,4,5}}
  \proofstepwidestar[2,3,4]{%
    \begin{aligned}[t]
      &\exists s\in\UB_{A,\leq}(\{x,y\})\,{}\\[-2pt]
      &\quad\forall u\in\UB_{A,\leq}(\{x,y\})\,(s\leq u)
    \end{aligned}}{%
    \ProofRefStack{%
      \FormulaRefAuto{PairJoinExistenceAxiom}[ax]\\
      (2,3,4)}}
  \proofstepwidestar[1,2,3,4]{\operatorname{SupEx}_{A,\leq}(x,y)}{%
    \ProofRefStack{%
      \FormulaRefAuto{PairSupremumExistence}[def]\\
      \bigl(1,\rAI{\rAI{3,4},8}\bigr)}}
  \proofstepwidestar[2,4,5]{%
    \begin{aligned}[t]
      &\exists s\in\UB_{A,\leq}(\{y,z\})\,{}\\[-2pt]
      &\quad\forall u\in\UB_{A,\leq}(\{y,z\})\,(s\leq u)
    \end{aligned}}{%
    \ProofRefStack{%
      \FormulaRefAuto{PairJoinExistenceAxiom}[ax]\\
      (2,4,5)}}
  \proofstepwidestar[1,2,4,5]{\operatorname{SupEx}_{A,\leq}(y,z)}{%
    \ProofRefStack{%
      \FormulaRefAuto{PairSupremumExistence}[def]\\
      \bigl(1,\rAI{\rAI{4,5},10}\bigr)}}
  \proofstepwidestar[2,3,4,5]{%
    \begin{aligned}[t]
      &\exists s\in\UB_{A,\leq}(\{x\NeutralSemiOp{}y,z\})\,{}\\[-2pt]
      &\quad\forall u\in\UB_{A,\leq}(\{x\NeutralSemiOp{}y,z\})\,(s\leq u)
    \end{aligned}}{%
    \ProofRefStack{%
      \FormulaRefAuto{PairJoinExistenceAxiom}[ax]\\
      (2,6,5)}}
  \proofstepwidestar[1,2,3,4,5]{%
    \operatorname{SupEx}_{A,\leq}(x\NeutralSemiOp{}y,z)}{%
    \ProofRefStack{%
      \FormulaRefAuto{PairSupremumExistence}[def]\\
      \bigl(1,\rAI{\rAI{6,5},12}\bigr)}}
  \proofstepwidestar[2,3,4,5]{%
    \begin{aligned}[t]
      &\exists s\in\UB_{A,\leq}(\{x,y\NeutralSemiOp{}z\})\,{}\\[-2pt]
      &\quad\forall u\in\UB_{A,\leq}(\{x,y\NeutralSemiOp{}z\})\,(s\leq u)
    \end{aligned}}{%
    \ProofRefStack{%
      \FormulaRefAuto{PairJoinExistenceAxiom}[ax]\\
      (2,3,7)}}
  \proofstepwidestar[1,2,3,4,5]{%
    \operatorname{SupEx}_{A,\leq}(x,y\NeutralSemiOp{}z)}{%
    \ProofRefStack{%
      \FormulaRefAuto{PairSupremumExistence}[def]\\
      \bigl(1,\rAI{\rAI{3,7},14}\bigr)}}
  \proofstepwidestar[2,3,4]{%
    x\NeutralSemiOp{}y=\sup_{A,\leq}(\{x,y\})}{%
    \FormulaRefAuto{PairJoinEvaluationAxiom}[ax]{2,3,4}}
  \proofstepwidestar[2,4,5]{%
    y\NeutralSemiOp{}z=\sup_{A,\leq}(\{y,z\})}{%
    \FormulaRefAuto{PairJoinEvaluationAxiom}[ax]{2,4,5}}
  \proofstepwidestar[2,3,4,5]{%
    \operatorname{SupEx}_{A,\leq}(\sup_{A,\leq}(\{x,y\}),z)}{%
    \rIE{16,13}}
  \proofstepwidestar[2,3,4,5]{%
    \operatorname{SupEx}_{A,\leq}(x,\sup_{A,\leq}(\{y,z\}))}{%
    \rIE{17,15}}
  \proofstepwidestar[1,2,3,4,5]{%
    \begin{aligned}[t]
      &\sup_{A,\leq}\bigl(\{\sup_{A,\leq}(\{x,y\}),z\}\bigr)\\[-2pt]
      &\qquad=
      \sup_{A,\leq}\bigl(\{x,\sup_{A,\leq}(\{y,z\})\}\bigr)
    \end{aligned}}{%
    \FormulaRefAuto{PairSupremumAssociative}[thm]{1,9,11,18,19}}
  \proofstepwidestar[2,3,4,5]{%
    \begin{aligned}[t]
      &\bigl(x\NeutralSemiOp{}y\bigr)\NeutralSemiOp{}z\\[-2pt]
      &\quad=\sup_{A,\leq}
        \bigl(\{\sup_{A,\leq}(\{x,y\}),z\}\bigr)
    \end{aligned}}{%
    \ProofRefStack{%
      \rIE{16,\FormulaRefAuto{PairJoinEvaluationAxiom}[ax]{2,6,5}}}}
  \proofstepwidestar[2,3,4,5]{%
    \begin{aligned}[t]
      &x\NeutralSemiOp\bigl(y\NeutralSemiOp{}z\bigr)\\[-2pt]
      &\quad=\sup_{A,\leq}
        \bigl(\{x,\sup_{A,\leq}(\{y,z\})\}\bigr)
    \end{aligned}}{%
    \ProofRefStack{%
      \rIE{17,\FormulaRefAuto{PairJoinEvaluationAxiom}[ax]{2,3,7}}}}
  \proofstepwidestar[2,3,4,5]{%
    \sup_{A,\leq}\bigl(\{x,\sup_{A,\leq}(\{y,z\})\}\bigr)
      =x\NeutralSemiOp\bigl(y\NeutralSemiOp{}z\bigr)}{%
    \FormulaRefAuto{a=b\vdash b=a}{22}}
  \proofstepwidestar[1,2,3,4,5]{%
    \bigl(x\NeutralSemiOp{}y\bigr)\NeutralSemiOp{}z
      =x\NeutralSemiOp\bigl(y\NeutralSemiOp{}z\bigr)}{%
    \rChain{21,20,23}}
\end{tabproofwide}

\FormulaThmDeltaK[Charakterisierung durch Paarsuprema]{%
  \PartOrd{A,\leq}\dsep \PairJoinOp{A,\NeutralSemiOp}
  \vdash \Semi{A,\NeutralSemiOp}
}{PairJoinCharacterization}{%
  \DeltaRow{Mengen}{A\dsep x\dsep y\dsep z}
  \DeltaRow{Relationsoperatoren}{\leq}
  \DeltaRow{Funktionen}{\NeutralSemiOp}
}
\begin{tabproof}
  \proofstep{1}{\PartOrd{A,\leq}}{\rA}
  \proofstep{2}{\PairJoinOp{A,\NeutralSemiOp}}{\rA}
  \proofstep{3}{x\in A}{\rA}
  \proofstep{4}{y\in A}{\rA}
  \proofstep{5}{z\in A}{\rA}
  \proofstep{2,3,4}{x\NeutralSemiOp{}y\in A}{%
    \FormulaRefAuto{PairJoinClosureAxiom}[ax]{2,3,4}}
  \proofstep{1,2,3,4,5}{%
    \bigl(x\NeutralSemiOp{}y\bigr)\NeutralSemiOp{}z
      =x\NeutralSemiOp\bigl(y\NeutralSemiOp{}z\bigr)}{%
    \FormulaRefAuto{PairJoinAssociative}[thm]{1,2,3,4,5}}
  \proofstep{1,2,3,4}{%
    x\NeutralSemiOp{}y=y\NeutralSemiOp{}x}{%
    \FormulaRefAuto{PairJoinCommutative}[thm]{1,2,3,4}}
  \proofstep{1,2,3}{x\NeutralSemiOp{}x=x}{%
    \FormulaRefAuto{PairJoinIdempotent}[thm]{1,2,3}}
  \proofstep{1,2}{\Semi{A,\NeutralSemiOp}}{%
    \FormulaRefAuto{Semilattice}[def]{6,7,8,9}}
\end{tabproof}

\FormulaThmDeltaK[Die Paarsupremumsoperation gewinnt die Ordnung zurück]{%
  \begin{aligned}[t]
    &\PartOrd{A,\leq}\dsep \PairJoinOp{A,\NeutralSemiOp}
      \dsep x,y\in A\vdash{}\\
    &x\leq y\leftrightarrow x\InducedLeq y
  \end{aligned}
}{PairJoinRecoversOrder}{%
  \DeltaRow{Mengen}{A\dsep x\dsep y}
  \DeltaRow{Relationsoperatoren}{\leq}
  \DeltaRow{Funktionen}{\NeutralSemiOp}
  \DeltaPrem{Von \(\Semi{A,\NeutralSemiOp}\) supremal induziert}{\InducedLeq}
}
\begin{tabproofsplitwide}
  \proofpartwide{\(\vdash\)}
    \proofstepwidestar[1]{\PartOrd{A,\leq}}{\rA}
    \proofstepwidestar[2]{\PairJoinOp{A,\NeutralSemiOp}}{\rA}
    \proofstepwidestar[3]{x\in A}{\rA}
    \proofstepwidestar[4]{y\in A}{\rA}
    \proofstepwidestar[5]{x\leq y}{\rA}
    \proofstepwidestar[1,4]{y\leq y}{%
      \FormulaRefAuto{EqRelReflexivity}[ax]{1,4}}
    \proofstepwidestar[1,2,3,4,5]{x\NeutralSemiOp{}y\leq y}{%
      \FormulaRefAuto{PairJoinLeastUpper}[thm]{1,2,3,4,4,5,6}}
    \proofstepwidestar[1,2,3,4]{y\leq x\NeutralSemiOp{}y}{%
      \FormulaRefAuto{PairJoinUpperRight}[thm]{1,2,3,4}}
    \proofstepwidestar[1,2,3,4,5]{x\NeutralSemiOp{}y=y}{%
      \FormulaRefAuto{x\leq y\dsep y\leq x\vdash x=y}[ax]{1,7,8}}
    \proofstepwidestar[1,2,3,4,5]{x\InducedLeq y}{%
      \FormulaRefAuto{P\leftrightarrow Q, Q\vdash P}{%
        \FormulaRefAuto{JoinOrderEvaluationAxiom}[ax]{%
          \FormulaRefAuto{PairJoinCharacterization}[thm]{1,2},3,4},9}}
  \closeproofpartwide

  \proofpartwide{\(\dashv\)}
    \proofstepwidestar[1]{\PartOrd{A,\leq}}{\rA}
    \proofstepwidestar[2]{\PairJoinOp{A,\NeutralSemiOp}}{\rA}
    \proofstepwidestar[3]{x\in A}{\rA}
    \proofstepwidestar[4]{y\in A}{\rA}
    \proofstepwidestar[5]{x\InducedLeq y}{\rA}
    \proofstepwidestar[1,2,3,4,5]{x\NeutralSemiOp{}y=y}{%
      \FormulaRefAuto{JoinOrderEvaluation}[thm]{%
        \FormulaRefAuto{PairJoinCharacterization}[thm]{1,2},5}}
    \proofstepwidestar[1,2,3,4]{x\leq x\NeutralSemiOp{}y}{%
      \FormulaRefAuto{PairJoinUpperLeft}[thm]{1,2,3,4}}
    \proofstepwidestar[2,3,4,5]{x\leq y}{\rIE{6,7}}
  \closeproofpartwide
\end{tabproofsplitwide}

Damit ist die zentrale Äquivalenz bewiesen:
\[
  \boxed{
    \begin{aligned}
      &\text{assoziative, kommutative und idempotente binäre Operation}\\
      &\qquad\Longleftrightarrow\\
      &\text{partielle Ordnung mit Suprema aller Elementpaare.}
    \end{aligned}}
\]

\section{Die infimale Lesart des Halbverbandes}

Die infimale Lesart benötigt keine zweite Halbverbandsdefinition und keine
Wiederholung der Ordnungsaxiome. Sei \(\leq\) die von
\(\Semi{A}\) supremal induzierten Relationsoperator. Dann verwenden wir
für ihre in Band~11 definierte duale partielle Ordnung unmittelbar die
Schreibweise \(\geq\). Insbesondere dürfen wir im Folgenden ohne erneuten
Dualitätsbeweis verwenden:
\[
  x\InducedGeq y\quad\Longleftrightarrow\quad y\InducedLeq x.
\]
Die Umkehrung der Leserichtung ist also bereits ordnungstheoretisch
abgedeckt. Hier ist nur noch zu zeigen, wie sie mit der
Halbverbandsoperation zusammenwirkt.

\FormulaThmDeltaK[Auswertung der infimalen Lesart]{%
  \Semi{A}\dsep x,y\in A\vdash
  \bigl(x\InducedGeq y\leftrightarrow x\OmittedSemiOp{}y=x\bigr)
}{MeetOrderEvaluation}{%
  \DeltaRow{Mengen}{A\dsep x\dsep y}
  \DeltaPrem{Von \(\Semi{A}\) supremal induzierter Relationsoperator}{\InducedLeq}
  \DeltaPrem{Dualer, infimal gelesener Relationsoperator}{\InducedGeq}
}
\begin{tabproofwide}
  \proofstepwidestar[1]{\Semi{A}}{\rA}
  \proofstepwidestar[2]{x\in A}{\rA}
  \proofstepwidestar[3]{y\in A}{\rA}
  \proofstepwide[2,3]{x\InducedGeq y}{\leftrightarrow}{y\InducedLeq x}{%
    \FormulaRefAuto{OrderDualRelation}[def]{2,3}}
  \proofstepwide[1,2,3]{}{\leftrightarrow}{y\OmittedSemiOp{}x=x}{%
    \FormulaRefAuto{JoinOrderEvaluationAxiom}[ax]{1,3,2}}
  \proofstepwide[1,2,3]{}{\leftrightarrow}{x\OmittedSemiOp{}y=x}{%
    \rLRS{\FormulaRefAuto{SemilatticeCommutativityAxiom}[ax]{1,2,3},5}}
\end{tabproofwide}

Band~11 definiert das Infimum bereits für jede Teilmenge \(T\subseteq A\).
Für ein Elementpaar ist daher keine weitere Definition erforderlich. Unter
der dualen Ordnung \(\geq\) ist die Halbverbandsoperation das Infimum von
\(\{x,y\}\).

\FormulaThmDeltaK[Erste Untere-Schranken-Eigenschaft der infimalen Lesart]{%
  \Semi{A}\dsep x,y\in A
  \vdash x\OmittedSemiOp{}y\InducedGeq x
}{MeetLowerLeft}{%
  \DeltaRow{Mengen}{A\dsep x\dsep y}
  \DeltaPrem{Von \(\Semi{A}\) supremal induzierter Relationsoperator}{\InducedLeq}
  \DeltaPrem{Dualer, infimal gelesener Relationsoperator}{\InducedGeq}
}
\begin{tabproof}
  \proofstep{1}{\Semi{A}}{\rA}
  \proofstep{2}{x\in A}{\rA}
  \proofstep{3}{y\in A}{\rA}
  \proofstep{1,2,3}{x\OmittedSemiOp{}y\in A}{%
    \FormulaRefAuto{SemilatticeClosureAxiom}[ax]{1,2,3}}
  \proofstep{1,2,3}{x\InducedLeq x\OmittedSemiOp{}y}{%
    \FormulaRefAuto{JoinUpperLeft}[thm]{1,2,3}}
  \proofstep{1,2,3}{x\OmittedSemiOp{}y\InducedGeq x}{%
    \FormulaRefAuto{P\leftrightarrow Q, Q\vdash P}{%
      \FormulaRefAuto{OrderDualRelation}[def]{4,2},5}}
\end{tabproof}

\FormulaThmDeltaK[Zweite Untere-Schranken-Eigenschaft der infimalen Lesart]{%
  \Semi{A}\dsep x,y\in A
  \vdash x\OmittedSemiOp{}y\InducedGeq y
}{MeetLowerRight}{%
  \DeltaRow{Mengen}{A\dsep x\dsep y}
  \DeltaPrem{Von \(\Semi{A}\) supremal induzierter Relationsoperator}{\InducedLeq}
  \DeltaPrem{Dualer, infimal gelesener Relationsoperator}{\InducedGeq}
}
\begin{tabproof}
  \proofstep{1}{\Semi{A}}{\rA}
  \proofstep{2}{x\in A}{\rA}
  \proofstep{3}{y\in A}{\rA}
  \proofstep{1,2,3}{x\OmittedSemiOp{}y\in A}{%
    \FormulaRefAuto{SemilatticeClosureAxiom}[ax]{1,2,3}}
  \proofstep{1,2,3}{y\InducedLeq x\OmittedSemiOp{}y}{%
    \FormulaRefAuto{JoinUpperRight}[thm]{1,2,3}}
  \proofstep{1,2,3}{x\OmittedSemiOp{}y\InducedGeq y}{%
    \FormulaRefAuto{P\leftrightarrow Q, Q\vdash P}{%
      \FormulaRefAuto{OrderDualRelation}[def]{4,3},5}}
\end{tabproof}

\FormulaThmDeltaK[Die Halbverbandsoperation liefert das Paarinfimum]{%
  \Semi{A}\dsep x,y\in A
  \vdash x\OmittedSemiOp{}y=\inf_{A,\InducedGeq}(\{x,y\})
}{MeetIsPairInfimum}{%
  \DeltaRow{Mengen}{A\dsep x\dsep y\dsep u\dsep z}
  \DeltaPrem{Von \(\Semi{A}\) supremal induzierter Relationsoperator}{\InducedLeq}
  \DeltaPrem{Dualer, infimal gelesener Relationsoperator}{\InducedGeq}
}
\begin{tabproofwide}
  \proofstepwidestar[1]{\Semi{A}}{\rA}
  \proofstepwidestar[2]{x\in A}{\rA}
  \proofstepwidestar[3]{y\in A}{\rA}
  \proofstepwidestar[1]{\PartOrd{A,\InducedLeq}}{%
    \FormulaRefAuto{JoinSemilatticeInducesOrder}[thm]{1}}
  \proofstepwidestar[2,3]{\{x,y\}\subseteq A}{%
    \FormulaRefAuto{a \in A,\, b \in A \vdash \{a,b\} \subseteq A}{2,3}}
  \proofstepwidestar[1,2,3]{x\OmittedSemiOp{}y\in A}{%
    \FormulaRefAuto{SemilatticeClosureAxiom}[ax]{1,2,3}}
  \proofstepwidestar[1,2,3]{x\InducedLeq x\OmittedSemiOp{}y}{%
    \FormulaRefAuto{JoinUpperLeft}[thm]{1,2,3}}
  \proofstepwidestar[1,2,3]{y\InducedLeq x\OmittedSemiOp{}y}{%
    \FormulaRefAuto{JoinUpperRight}[thm]{1,2,3}}
  \proofstepwidestar[1,2,3]{%
    \begin{aligned}[t]
      &x\OmittedSemiOp{}y\in\UB_{A,\InducedLeq}(\{x,y\})\\[-2pt]
      &\quad\leftrightarrow
        x\OmittedSemiOp{}y\in A\land
        x\InducedLeq x\OmittedSemiOp{}y\land
        y\InducedLeq x\OmittedSemiOp{}y
    \end{aligned}}{%
    \FormulaRefAuto{PairUpperBoundSetCharacterization}[thm]{4,2,3}}
  \proofstepwidestar[1,2,3]{%
    x\OmittedSemiOp{}y\in\UB_{A,\InducedLeq}(\{x,y\})}{%
    \FormulaRefAuto{P\leftrightarrow Q, Q\vdash P}{%
      9,\rAI{6,\rAI{7,8}}}}
  \proofstepwidestar[11]{u\in\UB_{A,\InducedLeq}(\{x,y\})}{\rA}
  \proofstepwidestar[2,3]{%
    \begin{aligned}[t]
      u\in\UB_{A,\InducedLeq}(\{x,y\})
      \leftrightarrow
      \bigl(u\in A\land x\InducedLeq u\land y\InducedLeq u\bigr)
    \end{aligned}}{%
    \FormulaRefAuto{PairUpperBoundSetCharacterization}[thm]{4,2,3}}
  \proofstepwidestar[2,3,11]{%
    u\in A\land x\InducedLeq u\land y\InducedLeq u}{%
    \FormulaRefAuto{P\leftrightarrow Q, P\vdash Q}{12,11}}
  \proofstepwidestar[2,3,11]{u\in A}{\rAEa{13}}
  \proofstepwidestar[2,3,11]{x\InducedLeq u}{\rAEa{\rAEb{13}}}
  \proofstepwidestar[2,3,11]{y\InducedLeq u}{\rAEb{\rAEb{13}}}
  \proofstepwidestar[1,2,3,11]{x\OmittedSemiOp{}y\InducedLeq u}{%
    \FormulaRefAuto{JoinLeastUpper}[thm]{1,2,3,14,15,16}}
  \proofstepwidestar[1,2,3]{%
    \forall u\in\UB_{A,\InducedLeq}(\{x,y\})\,
      (x\OmittedSemiOp{}y\InducedLeq u)}{%
    \rUI{\rRI{11,17}}}
  \proofstepwidestar[1,2,3]{%
    \exists s\in\UB_{A,\InducedLeq}(\{x,y\})\,
      \forall u\in\UB_{A,\InducedLeq}(\{x,y\})\,(s\InducedLeq u)}{%
    \rEI{\rAI{10,18}}}
  \proofstepwidestar[1,2,3]{%
    \sup_{A,\InducedLeq}(\{x,y\})
      =\inf_{A,\InducedGeq}(\{x,y\})}{%
    \FormulaRefAuto{DualSupremumEqualsInfimum}[thm]{4,5,19}}
  \proofstepwidestar[1,2,3]{%
    x\OmittedSemiOp{}y=\sup_{A,\InducedLeq}(\{x,y\})}{%
    \FormulaRefAuto{JoinIsPairSupremum}[thm]{1,2,3}}
  \proofstepwidestar[1,2,3]{%
    x\OmittedSemiOp{}y=\inf_{A,\InducedGeq}(\{x,y\})}{%
    \rChain{21,20}}
\end{tabproofwide}

\section{Beispiele von Halbverbänden}

\subsection{Vereinigung und Durchschnitt auf Potenzmengen}

Die beiden grundlegenden Mengenoperationen aus Band~03 liefern die ersten
konkreten Beispiele der allgemeinen Axiomatik. Wir zeigen zunächst jede
Halbverbandsstruktur für sich und bestimmen anschließend ihre induzierte
Relation.

\FormulaThmDeltaK[Abgeschlossenheit der Vereinigung auf einer Potenzmenge]{%
  X,Y\in\powerset(S)\vdash X\cup Y\in\powerset(S)
}{PowerSetUnionClosure}{%
  \DeltaRow{Mengen}{S\dsep X\dsep Y}
}
\begin{tabproof}
  \proofstep{1}{X\in\powerset(S)}{\rA}
  \proofstep{2}{Y\in\powerset(S)}{\rA}
  \proofstep{1}{X\subseteq S}{%
    \FormulaRefAuto{x\in\powerset(A)\eqvdash x\subseteq A}{1}}
  \proofstep{2}{Y\subseteq S}{%
    \FormulaRefAuto{x\in\powerset(A)\eqvdash x\subseteq A}{2}}
  \proofstep{1,2}{X\cup Y\subseteq S}{%
    \FormulaRefAuto{A\subseteq M\dsep B\subseteq M \vdash A\cup B\subseteq M}{3,4}}
  \proofstep{1,2}{X\cup Y\in\powerset(S)}{%
    \FormulaRefAuto{x\in\powerset(A)\eqvdash x\subseteq A}{5}}
\end{tabproof}

\FormulaThmDeltaK[Vereinigung ist eine Halbverbandsoperation]{%
  \Semi{\powerset(S),\cup}
}{PowerSetUnionSemilattice}{%
  \DeltaRow{Mengen}{S\dsep X\dsep Y\dsep Z}
  \DeltaRow{Funktionen}{\cup}
}
\begin{tabproofwide}
  \proofstepwidestar[1]{X\in\powerset(S)}{\rA}
  \proofstepwidestar[2]{Y\in\powerset(S)}{\rA}
  \proofstepwidestar[3]{Z\in\powerset(S)}{\rA}
  \proofstepwidestar[1,2]{X\cup Y\in\powerset(S)}{%
    \FormulaRefAuto{PowerSetUnionClosure}[thm]{1,2}}
  \proofstepwidestar[1,2,3]{(X\cup Y)\cup Z=X\cup(Y\cup Z)}{%
    \FormulaRefAuto{(A \cup B) \cup C = A \cup (B \cup C)}}
  \proofstepwidestar[1,2]{X\cup Y=Y\cup X}{%
    \FormulaRefAuto{A \cup B = B \cup A}}
  \proofstepwidestar[1]{X\cup X=X}{%
    \FormulaRefAuto{A \cup A = A}}
  \proofstepwidestar[]{\Semi{\powerset(S),\cup}}{%
    \FormulaRefAuto{Semilattice}[def]{4,5,6,7}}
\end{tabproofwide}

\FormulaThmDeltaK[Abgeschlossenheit des Durchschnitts auf einer Potenzmenge]{%
  X,Y\in\powerset(S)\vdash X\cap Y\in\powerset(S)
}{PowerSetIntersectionClosure}{%
  \DeltaRow{Mengen}{S\dsep X\dsep Y}
}
\begin{tabproof}
  \proofstep{1}{X\in\powerset(S)}{\rA}
  \proofstep{2}{Y\in\powerset(S)}{\rA}
  \proofstep{1}{X\subseteq S}{%
    \FormulaRefAuto{x\in\powerset(A)\eqvdash x\subseteq A}{1}}
  \proofstep{1}{X\cap Y\subseteq S}{%
    \FormulaRefAuto{A\subseteq M \vdash A\cap B\subseteq M}{3}}
  \proofstep{1,2}{X\cap Y\in\powerset(S)}{%
    \FormulaRefAuto{x\in\powerset(A)\eqvdash x\subseteq A}{4}}
\end{tabproof}

\FormulaThmDeltaK[Durchschnitt ist eine Halbverbandsoperation]{%
  \Semi{\powerset(S),\cap}
}{PowerSetIntersectionSemilattice}{%
  \DeltaRow{Mengen}{S\dsep X\dsep Y\dsep Z}
  \DeltaRow{Funktionen}{\cap}
}
\begin{tabproofwide}
  \proofstepwidestar[1]{X\in\powerset(S)}{\rA}
  \proofstepwidestar[2]{Y\in\powerset(S)}{\rA}
  \proofstepwidestar[3]{Z\in\powerset(S)}{\rA}
  \proofstepwidestar[1,2]{X\cap Y\in\powerset(S)}{%
    \FormulaRefAuto{PowerSetIntersectionClosure}[thm]{1,2}}
  \proofstepwidestar[1,2,3]{(X\cap Y)\cap Z=X\cap(Y\cap Z)}{%
    \FormulaRefAuto{(A \cap B) \cap C = A \cap (B \cap C)}}
  \proofstepwidestar[1,2]{X\cap Y=Y\cap X}{%
    \FormulaRefAuto{A \cap B = B \cap A}}
  \proofstepwidestar[1]{X\cap X=X}{%
    \FormulaRefAuto{a=b\vdash b=a}{\FormulaRefAuto{A = A \cap A}}}
  \proofstepwidestar[]{\Semi{\powerset(S),\cap}}{%
    \FormulaRefAuto{Semilattice}[def]{4,5,6,7}}
\end{tabproofwide}

\FormulaThmDeltaK[Inklusion durch Vereinigung]{%
  X,Y\in\powerset(S)\vdash
  \bigl(X\subseteq Y\leftrightarrow X\cup Y=Y\bigr)
}{SubsetIffUnionEqualsRight}{%
  \DeltaRow{Mengen}{S\dsep X\dsep Y}
}
\begin{tabproofsplitwide}
  \proofpartwide{\(\vdash\)}
    \proofstepwidestar[1]{X\in\powerset(S)}{\rA}
    \proofstepwidestar[2]{Y\in\powerset(S)}{\rA}
    \proofstepwidestar[3]{X\subseteq Y}{\rA}
    \proofstepwidestar[]{Y\subseteq Y}{\FormulaRefAuto{A\subseteq A}}
    \proofstepwidestar[3]{X\cup Y\subseteq Y}{%
      \FormulaRefAuto{A\subseteq C,\, B\subseteq C \vdash A\cup B\subseteq C}{3,4}}
    \proofstepwidestar[]{Y\subseteq Y\cup X}{%
      \FormulaRefAuto{A \subseteq A \cup B}}
    \proofstepwidestar[]{Y\cup X=X\cup Y}{%
      \FormulaRefAuto{A \cup B = B \cup A}}
    \proofstepwidestar[]{Y\subseteq X\cup Y}{\rIE{7,6}}
    \proofstepwidestar[3]{X\cup Y=Y}{%
      \FormulaRefAuto{A \subseteq B \land B \subseteq A \eqvdash A = B}{5,8}}
  \closeproofpartwide

  \proofpartwide{\(\dashv\)}
    \proofstepwidestar[1]{X\in\powerset(S)}{\rA}
    \proofstepwidestar[2]{Y\in\powerset(S)}{\rA}
    \proofstepwidestar[3]{X\cup Y=Y}{\rA}
    \proofstepwidestar[]{X\subseteq X\cup Y}{%
      \FormulaRefAuto{A \subseteq A \cup B}}
    \proofstepwidestar[3]{X\subseteq Y}{\rIE{3,4}}
  \closeproofpartwide
\end{tabproofsplitwide}

\FormulaThmDeltaK[Die Vereinigungsrelation ist die Inklusion]{%
  X,Y\in\powerset(S)\vdash
  \bigl(X\InducedLeq Y\leftrightarrow X\subseteq Y\bigr)
}{PowerSetUnionOrderIsInclusion}{%
  \DeltaRow{Mengen}{S\dsep X\dsep Y}
  \DeltaPrem{\makecell[l]{Supremal durch
    \(\Semi{\powerset(S),\cup}\)\\induzierter Relationsoperator}}{\InducedLeq}
}
\begin{tabproofwide}
  \proofstepwidestar[1]{X\in\powerset(S)}{\rA}
  \proofstepwidestar[2]{Y\in\powerset(S)}{\rA}
  \proofstepwidestar[]{\Semi{\powerset(S),\cup}}{%
    \FormulaRefAuto{PowerSetUnionSemilattice}[thm]}
  \proofstepwide[1,2]{X\InducedLeq Y}{\leftrightarrow}{X\cup Y=Y}{%
    \FormulaRefAuto{JoinOrderEvaluationAxiom}[ax]{3,1,2}}
  \proofstepwide[1,2]{}{\leftrightarrow}{X\subseteq Y}{%
    \FormulaRefAuto{SubsetIffUnionEqualsRight}[thm]{1,2}}
\end{tabproofwide}

\FormulaThmDeltaK[Die duale Durchschnittsrelation ist die Inklusion]{%
  X,Y\in\powerset(S)\vdash
  \bigl(X\InducedGeq Y\leftrightarrow X\subseteq Y\bigr)
}{PowerSetIntersectionOrderIsInclusion}{%
  \DeltaRow{Mengen}{S\dsep X\dsep Y}
  \DeltaPrem{\makecell[l]{Dualer, infimal durch
    \(\Semi{\powerset(S),\cap}\)\\gelesener Relationsoperator}}{\InducedGeq}
}
\begin{tabproofwide}
  \proofstepwidestar[1]{X\in\powerset(S)}{\rA}
  \proofstepwidestar[2]{Y\in\powerset(S)}{\rA}
  \proofstepwidestar[]{\Semi{\powerset(S),\cap}}{%
    \FormulaRefAuto{PowerSetIntersectionSemilattice}[thm]}
  \proofstepwide[1,2]{X\InducedGeq Y}{\leftrightarrow}{X\cap Y=X}{%
    \FormulaRefAuto{MeetOrderEvaluation}[thm]{3,1,2}}
  \proofstepwide[1,2]{}{\leftrightarrow}{X\subseteq Y}{%
    \FormulaRefAuto{A\subseteq B \eqvdash A \cap B = A}}
\end{tabproofwide}

Wird nur die infimale Orientierung betrachtet, benennt man ihre Relation
üblicherweise wieder mit \(\leq\). In dieser lokalen Konvention erzeugen
sowohl die Vereinigung als Supremumsoperation als auch der Durchschnitt als
Infimumsoperation genau die Inklusionsordnung.

\chapter{Verbände}

\section{Axiomatische Definition}

Zwei beliebige Halbverbandsoperationen auf derselben Menge müssen nicht
dieselbe Ordnung erzeugen. Erst die beiden Absorptionsgesetze verbinden sie zu
einem Verband. Die folgende Axiomentafel führt alle acht Halbverbandseigenschaften
und beide Absorptionseigenschaften ausdrücklich auf.

\FormulaAxiomDeltaK[Absorption des Infimums]{%
  x,y\in A\vdash x\MeetOp\bigl(x\JoinOp{}y\bigr)=x%
}{LatticeMeetAbsorptionLaw}{%
  \DeltaRow{Mengen}{A\dsep x\dsep y}
  \DeltaRow{Funktionen}{\MeetOp\dsep \JoinOp}
}
\par\smallskip

\FormulaAxiomDeltaK[Absorption des Supremums]{%
  x,y\in A\vdash x\JoinOp\bigl(x\MeetOp{}y\bigr)=x%
}{LatticeJoinAbsorptionLaw}{%
  \DeltaRow{Mengen}{A\dsep x\dsep y}
  \DeltaRow{Funktionen}{\MeetOp\dsep \JoinOp}
}

\FormulaDefDeltaK[Begriff des algebraischen Verbandes]{%
  \LatticeAlg{A,\MeetOp,\JoinOp}%
}{AlgebraicLattice}{%
  \DeltaRow{Mengen}{A\dsep x\dsep y\dsep z}
  \DeltaRow{Funktionen}{\MeetOp\dsep \JoinOp}
  \DeltaRow{\textbf{Axiome für \(\MeetOp\)}}{}
  \DeltaRow{Abgeschlossenheit}{%
    x,y\in A\vdash x\MeetOp{}y\in A}
    [\FormulaRefAuto{SemilatticeClosureLaw}[ax]]
  \DeltaRow{Assoziativität}{%
    \begin{aligned}[t]
      &x,y,z\in A\vdash{}\\[-2pt]
      &(x\MeetOp{}y)\MeetOp{}z=x\MeetOp(y\MeetOp{}z)
    \end{aligned}}
    [\FormulaRefAuto{SemilatticeAssociativityLaw}[ax]]
  \DeltaRow{Kommutativität}{%
    x,y\in A\vdash x\MeetOp{}y=y\MeetOp{}x}
    [\FormulaRefAuto{SemilatticeCommutativityLaw}[ax]]
  \DeltaRow{Idempotenz}{%
    x\in A\vdash x\MeetOp{}x=x}
    [\FormulaRefAuto{SemilatticeIdempotenceLaw}[ax]]
  \DeltaRow{\textbf{Axiome für \(\JoinOp\)}}{}
  \DeltaRow{Abgeschlossenheit}{%
    x,y\in A\vdash x\JoinOp{}y\in A}
    [\FormulaRefAuto{SemilatticeClosureLaw}[ax]]
  \DeltaRow{Assoziativität}{%
    \begin{aligned}[t]
      &x,y,z\in A\vdash{}\\[-2pt]
      &(x\JoinOp{}y)\JoinOp{}z=x\JoinOp(y\JoinOp{}z)
    \end{aligned}}
    [\FormulaRefAuto{SemilatticeAssociativityLaw}[ax]]
  \DeltaRow{Kommutativität}{%
    x,y\in A\vdash x\JoinOp{}y=y\JoinOp{}x}
    [\FormulaRefAuto{SemilatticeCommutativityLaw}[ax]]
  \DeltaRow{Idempotenz}{%
    x\in A\vdash x\JoinOp{}x=x}
    [\FormulaRefAuto{SemilatticeIdempotenceLaw}[ax]]
  \DeltaRow{\textbf{Verträglichkeitsaxiome}}{}
  \DeltaRow{\(\MeetOp\)-Absorption}{%
    x,y\in A\vdash x\MeetOp(x\JoinOp{}y)=x}
    [\FormulaRefAuto{LatticeMeetAbsorptionLaw}[ax]]
  \DeltaRow{\(\JoinOp\)-Absorption}{%
    x,y\in A\vdash x\JoinOp(x\MeetOp{}y)=x}
    [\FormulaRefAuto{LatticeJoinAbsorptionLaw}[ax]]
  \DeltaRow{\textbf{Neues Symbol}}{}
  \DeltaPrem{Verbände}{\LatticeAlg{A,\MeetOp,\JoinOp}}
}

\FormulaAxiomDeltaK[Zugriff auf die infimale Halbverbandsstruktur eines Verbandes]{%
  \LatticeAlg{A,\MeetOp,\JoinOp}\vdash \Semi{A,\MeetOp}%
}{LatticeMeetSemilatticeAxiom}{%
  \DeltaRow{Mengen}{A}
  \DeltaRow{Funktionen}{\MeetOp\dsep\JoinOp}
}
\par\smallskip

\FormulaAxiomDeltaK[Zugriff auf die supremale Halbverbandsstruktur eines Verbandes]{%
  \LatticeAlg{A,\MeetOp,\JoinOp}\vdash \Semi{A,\JoinOp}%
}{LatticeJoinSemilatticeAxiom}{%
  \DeltaRow{Mengen}{A}
  \DeltaRow{Funktionen}{\MeetOp\dsep\JoinOp}
}
\par\smallskip

\FormulaAxiomDeltaK[Zugriff auf das Absorptionsaxiom für \(\MeetOp\)]{%
  \LatticeAlg{A,\MeetOp,\JoinOp}\dsep x,y\in A
  \vdash x\MeetOp\bigl(x\JoinOp{}y\bigr)=x%
}{LatticeMeetAbsorptionAxiom}{%
  \DeltaRow{Mengen}{A\dsep x\dsep y}
  \DeltaRow{Funktionen}{\MeetOp\dsep \JoinOp}
}
\par\smallskip

\FormulaAxiomDeltaK[Zugriff auf das Absorptionsaxiom für \(\JoinOp\)]{%
  \LatticeAlg{A,\MeetOp,\JoinOp}\dsep x,y\in A
  \vdash x\JoinOp\bigl(x\MeetOp{}y\bigr)=x%
}{LatticeJoinAbsorptionAxiom}{%
  \DeltaRow{Mengen}{A\dsep x\dsep y}
  \DeltaRow{Funktionen}{\MeetOp\dsep \JoinOp}
}

\section{Homomorphismen und Isomorphismen von Verbänden}

Bei Verbänden müssen beide Operationen erhalten werden. Die Bezeichnung
führt daher für Quelle und Ziel jeweils den Träger, den Infimumsoperator und
den Supremumsoperator mit.

\FormulaDefDeltaK[Begriff des Verbandshomomorphismus]{%
  \LatticeHom{f;A,\MeetOp_A,\JoinOp_A;B,\MeetOp_B,\JoinOp_B}%
}{LatticeHomomorphism}{%
  \DeltaRow{Mengen}{A\dsep B\dsep x\dsep y}
  \DeltaRow{Funktionen}{%
    f\dsep\MeetOp_A\dsep\JoinOp_A\dsep\MeetOp_B\dsep\JoinOp_B}
  \DeltaPrem{Quellverband}{\LatticeAlg{A,\MeetOp_A,\JoinOp_A}}
  \DeltaPrem{Zielverband}{\LatticeAlg{B,\MeetOp_B,\JoinOp_B}}
  \DeltaPrem{Abbildung}{f\colon A\to B}
  \DeltaRow{Erhaltung des Infimums}{%
    x,y\in A\vdash
    f(x\MeetOp_A{}y)=f(x)\MeetOp_B f(y)}
  \DeltaRow{Erhaltung des Supremums}{%
    x,y\in A\vdash
    f(x\JoinOp_A{}y)=f(x)\JoinOp_B f(y)}
  \DeltaRow{\textbf{Neues Symbol}}{}
  \DeltaPrem{Verbandshomomorphismen}{%
    \LatticeHom{f;A,\MeetOp_A,\JoinOp_A;B,\MeetOp_B,\JoinOp_B}}
}
\par\smallskip

\FormulaAxiomDeltaK[Zugriff auf den Quellverband eines Verbandshomomorphismus]{%
  \LatticeHom{f;A,\MeetOp_A,\JoinOp_A;B,\MeetOp_B,\JoinOp_B}
  \vdash\LatticeAlg{A,\MeetOp_A,\JoinOp_A}%
}{LatticeHomSourceAxiom}{%
  \DeltaRow{Mengen}{A\dsep B}
  \DeltaRow{Funktionen}{%
    f\dsep\MeetOp_A\dsep\JoinOp_A\dsep\MeetOp_B\dsep\JoinOp_B}
}
\par\smallskip

\FormulaAxiomDeltaK[Zugriff auf den Zielverband eines Verbandshomomorphismus]{%
  \LatticeHom{f;A,\MeetOp_A,\JoinOp_A;B,\MeetOp_B,\JoinOp_B}
  \vdash\LatticeAlg{B,\MeetOp_B,\JoinOp_B}%
}{LatticeHomTargetAxiom}{%
  \DeltaRow{Mengen}{A\dsep B}
  \DeltaRow{Funktionen}{%
    f\dsep\MeetOp_A\dsep\JoinOp_A\dsep\MeetOp_B\dsep\JoinOp_B}
}
\par\smallskip

\FormulaAxiomDeltaK[Zugriff auf die Abbildung eines Verbandshomomorphismus]{%
  \LatticeHom{f;A,\MeetOp_A,\JoinOp_A;B,\MeetOp_B,\JoinOp_B}
  \vdash f\colon A\to B%
}{LatticeHomMapAxiom}{%
  \DeltaRow{Mengen}{A\dsep B}
  \DeltaRow{Funktionen}{%
    f\dsep\MeetOp_A\dsep\JoinOp_A\dsep\MeetOp_B\dsep\JoinOp_B}
}
\par\smallskip

\FormulaAxiomDeltaK[Erhaltung des Infimums durch einen Verbandshomomorphismus]{%
  \begin{aligned}[t]
    &\LatticeHom{f;A,\MeetOp_A,\JoinOp_A;B,\MeetOp_B,\JoinOp_B}
      \dsep x,y\in A\vdash{}\\
    &f(x\MeetOp_A{}y)=f(x)\MeetOp_B f(y)
  \end{aligned}
}{LatticeHomMeetEvaluationAxiom}{%
  \DeltaRow{Mengen}{A\dsep B\dsep x\dsep y}
  \DeltaRow{Funktionen}{%
    f\dsep\MeetOp_A\dsep\JoinOp_A\dsep\MeetOp_B\dsep\JoinOp_B}
}
\par\smallskip

\FormulaAxiomDeltaK[Erhaltung des Supremums durch einen Verbandshomomorphismus]{%
  \begin{aligned}[t]
    &\LatticeHom{f;A,\MeetOp_A,\JoinOp_A;B,\MeetOp_B,\JoinOp_B}
      \dsep x,y\in A\vdash{}\\
    &f(x\JoinOp_A{}y)=f(x)\JoinOp_B f(y)
  \end{aligned}
}{LatticeHomJoinEvaluationAxiom}{%
  \DeltaRow{Mengen}{A\dsep B\dsep x\dsep y}
  \DeltaRow{Funktionen}{%
    f\dsep\MeetOp_A\dsep\JoinOp_A\dsep\MeetOp_B\dsep\JoinOp_B}
}

\FormulaDefDeltaK[Begriff des Verbandsisomorphismus]{%
  \LatticeIso{f;A,\MeetOp_A,\JoinOp_A;B,\MeetOp_B,\JoinOp_B}%
}{LatticeIsomorphism}{%
  \DeltaRow{Mengen}{A\dsep B}
  \DeltaRow{Funktionen}{%
    f\dsep\MeetOp_A\dsep\JoinOp_A\dsep\MeetOp_B\dsep\JoinOp_B}
  \DeltaPrem{Verbandshomomorphismus}{%
    \LatticeHom{f;A,\MeetOp_A,\JoinOp_A;B,\MeetOp_B,\JoinOp_B}}
  \DeltaPrem{Bijektivität}{f\colon A\bij B}
  \DeltaRow{\textbf{Neues Symbol}}{}
  \DeltaPrem{Verbandsisomorphismen}{%
    \LatticeIso{f;A,\MeetOp_A,\JoinOp_A;B,\MeetOp_B,\JoinOp_B}}
}
\par\smallskip

\FormulaAxiomDeltaK[Ein Verbandsisomorphismus ist ein Homomorphismus]{%
  \LatticeIso{f;A,\MeetOp_A,\JoinOp_A;B,\MeetOp_B,\JoinOp_B}
  \vdash
  \LatticeHom{f;A,\MeetOp_A,\JoinOp_A;B,\MeetOp_B,\JoinOp_B}%
}{LatticeIsoHomAxiom}{%
  \DeltaRow{Mengen}{A\dsep B}
  \DeltaRow{Funktionen}{%
    f\dsep\MeetOp_A\dsep\JoinOp_A\dsep\MeetOp_B\dsep\JoinOp_B}
}
\par\smallskip

\FormulaAxiomDeltaK[Ein Verbandsisomorphismus ist bijektiv]{%
  \LatticeIso{f;A,\MeetOp_A,\JoinOp_A;B,\MeetOp_B,\JoinOp_B}
  \vdash f\colon A\bij B%
}{LatticeIsoBijectiveAxiom}{%
  \DeltaRow{Mengen}{A\dsep B}
  \DeltaRow{Funktionen}{%
    f\dsep\MeetOp_A\dsep\JoinOp_A\dsep\MeetOp_B\dsep\JoinOp_B}
}

\FormulaThmDeltaK[Ein Verbandshomomorphismus ist infimal ein Halbverbandshomomorphismus]{%
  \begin{aligned}[t]
    &\LatticeHom{f;A,\MeetOp_A,\JoinOp_A;B,\MeetOp_B,\JoinOp_B}
      \vdash{}\\
    &\SemilatticeHom{f;A,\MeetOp_A;B,\MeetOp_B}
  \end{aligned}
}{LatticeHomMeetSemilatticeHom}{%
  \DeltaRow{Mengen}{A\dsep B\dsep x\dsep y}
  \DeltaRow{Funktionen}{%
    f\dsep\MeetOp_A\dsep\JoinOp_A\dsep\MeetOp_B\dsep\JoinOp_B}
}
\begin{tabproofwide}
  \proofstepwidestar[1]{%
    \LatticeHom{f;A,\MeetOp_A,\JoinOp_A;B,\MeetOp_B,\JoinOp_B}}{\rA}
  \proofstepwidestar[1]{\LatticeAlg{A,\MeetOp_A,\JoinOp_A}}{%
    \FormulaRefAuto{LatticeHomSourceAxiom}[ax]{1}}
  \proofstepwidestar[1]{\LatticeAlg{B,\MeetOp_B,\JoinOp_B}}{%
    \FormulaRefAuto{LatticeHomTargetAxiom}[ax]{1}}
  \proofstepwidestar[1]{\Semi{A,\MeetOp_A}}{%
    \FormulaRefAuto{LatticeMeetSemilatticeAxiom}[ax]{2}}
  \proofstepwidestar[1]{\Semi{B,\MeetOp_B}}{%
    \FormulaRefAuto{LatticeMeetSemilatticeAxiom}[ax]{3}}
  \proofstepwidestar[1]{f\colon A\to B}{%
    \FormulaRefAuto{LatticeHomMapAxiom}[ax]{1}}
  \proofstepwidestar[7]{x\in A}{\rA}
  \proofstepwidestar[8]{y\in A}{\rA}
  \proofstepwidestar[1,7,8]{%
    f(x\MeetOp_A{}y)=f(x)\MeetOp_B f(y)}{%
    \FormulaRefAuto{LatticeHomMeetEvaluationAxiom}[ax]{1,7,8}}
  \proofstepwidestar[1]{\SemilatticeHom{f;A,\MeetOp_A;B,\MeetOp_B}}{%
    \FormulaRefAuto{SemilatticeHomomorphism}[def]{4,5,6,9}}
\end{tabproofwide}

\FormulaThmDeltaK[Ein Verbandshomomorphismus ist supremal ein Halbverbandshomomorphismus]{%
  \begin{aligned}[t]
    &\LatticeHom{f;A,\MeetOp_A,\JoinOp_A;B,\MeetOp_B,\JoinOp_B}
      \vdash{}\\
    &\SemilatticeHom{f;A,\JoinOp_A;B,\JoinOp_B}
  \end{aligned}
}{LatticeHomJoinSemilatticeHom}{%
  \DeltaRow{Mengen}{A\dsep B\dsep x\dsep y}
  \DeltaRow{Funktionen}{%
    f\dsep\MeetOp_A\dsep\JoinOp_A\dsep\MeetOp_B\dsep\JoinOp_B}
}
\begin{tabproofwide}
  \proofstepwidestar[1]{%
    \LatticeHom{f;A,\MeetOp_A,\JoinOp_A;B,\MeetOp_B,\JoinOp_B}}{\rA}
  \proofstepwidestar[1]{\LatticeAlg{A,\MeetOp_A,\JoinOp_A}}{%
    \FormulaRefAuto{LatticeHomSourceAxiom}[ax]{1}}
  \proofstepwidestar[1]{\LatticeAlg{B,\MeetOp_B,\JoinOp_B}}{%
    \FormulaRefAuto{LatticeHomTargetAxiom}[ax]{1}}
  \proofstepwidestar[1]{\Semi{A,\JoinOp_A}}{%
    \FormulaRefAuto{LatticeJoinSemilatticeAxiom}[ax]{2}}
  \proofstepwidestar[1]{\Semi{B,\JoinOp_B}}{%
    \FormulaRefAuto{LatticeJoinSemilatticeAxiom}[ax]{3}}
  \proofstepwidestar[1]{f\colon A\to B}{%
    \FormulaRefAuto{LatticeHomMapAxiom}[ax]{1}}
  \proofstepwidestar[7]{x\in A}{\rA}
  \proofstepwidestar[8]{y\in A}{\rA}
  \proofstepwidestar[1,7,8]{%
    f(x\JoinOp_A{}y)=f(x)\JoinOp_B f(y)}{%
    \FormulaRefAuto{LatticeHomJoinEvaluationAxiom}[ax]{1,7,8}}
  \proofstepwidestar[1]{\SemilatticeHom{f;A,\JoinOp_A;B,\JoinOp_B}}{%
    \FormulaRefAuto{SemilatticeHomomorphism}[def]{4,5,6,9}}
\end{tabproofwide}

\FormulaThmDeltaK[Die Identität ist ein Verbandshomomorphismus]{%
  \LatticeAlg{A,\MeetOp_A,\JoinOp_A}\vdash
  \LatticeHom{\Id_A;A,\MeetOp_A,\JoinOp_A;A,\MeetOp_A,\JoinOp_A}%
}{LatticeIdentityHomomorphism}{%
  \DeltaRow{Mengen}{A\dsep x\dsep y}
  \DeltaRow{Funktionen}{\Id_A\dsep\MeetOp_A\dsep\JoinOp_A}
}
\begin{tabproofwide}
  \proofstepwidestar[1]{\LatticeAlg{A,\MeetOp_A,\JoinOp_A}}{\rA}
  \proofstepwidestar[1]{\Semi{A,\MeetOp_A}}{%
    \FormulaRefAuto{LatticeMeetSemilatticeAxiom}[ax]{1}}
  \proofstepwidestar[1]{\Semi{A,\JoinOp_A}}{%
    \FormulaRefAuto{LatticeJoinSemilatticeAxiom}[ax]{1}}
  \proofstepwidestar[1]{%
    \SemilatticeHom{\Id_A;A,\MeetOp_A;A,\MeetOp_A}}{%
    \FormulaRefAuto{SemilatticeIdentityHomomorphism}[thm]{2}}
  \proofstepwidestar[1]{%
    \SemilatticeHom{\Id_A;A,\JoinOp_A;A,\JoinOp_A}}{%
    \FormulaRefAuto{SemilatticeIdentityHomomorphism}[thm]{3}}
  \proofstepwidestar[]{\Id_A\colon A\to A}{\FormulaRefAuto{IdA}}
  \proofstepwidestar[7]{x\in A}{\rA}
  \proofstepwidestar[8]{y\in A}{\rA}
  \proofstepwidestar[1,7,8]{%
    \Id_A(x\MeetOp_A{}y)=\Id_A(x)\MeetOp_A\Id_A(y)}{%
    \FormulaRefAuto{SemilatticeHomEvaluationAxiom}[ax]{4,7,8}}
  \proofstepwidestar[1,7,8]{%
    \Id_A(x\JoinOp_A{}y)=\Id_A(x)\JoinOp_A\Id_A(y)}{%
    \FormulaRefAuto{SemilatticeHomEvaluationAxiom}[ax]{5,7,8}}
  \proofstepwidestar[1]{%
    \LatticeHom{\Id_A;A,\MeetOp_A,\JoinOp_A;A,\MeetOp_A,\JoinOp_A}}{%
    \FormulaRefAuto{LatticeHomomorphism}[def]{1,1,6,9,10}}
\end{tabproofwide}

\FormulaThmDeltaK[Komposition von Verbandshomomorphismen]{%
  \begin{aligned}[t]
    &\LatticeHom{f;A,\MeetOp_A,\JoinOp_A;B,\MeetOp_B,\JoinOp_B}
      \dsep
      \LatticeHom{g;B,\MeetOp_B,\JoinOp_B;C,\MeetOp_C,\JoinOp_C}
      \vdash{}\\
    &\LatticeHom{g\circ f;A,\MeetOp_A,\JoinOp_A;C,\MeetOp_C,\JoinOp_C}
  \end{aligned}
}{LatticeHomComposition}{%
  \DeltaRow{Mengen}{A\dsep B\dsep C\dsep x\dsep y}
  \DeltaRow{Funktionen}{%
    f\dsep g\dsep
    \MeetOp_A\dsep\JoinOp_A\dsep
    \MeetOp_B\dsep\JoinOp_B\dsep
    \MeetOp_C\dsep\JoinOp_C}
}
\begin{tabproofwide}
  \proofstepwidestar[1]{%
    \LatticeHom{f;A,\MeetOp_A,\JoinOp_A;B,\MeetOp_B,\JoinOp_B}}{\rA}
  \proofstepwidestar[2]{%
    \LatticeHom{g;B,\MeetOp_B,\JoinOp_B;C,\MeetOp_C,\JoinOp_C}}{\rA}
  \proofstepwidestar[1]{\SemilatticeHom{f;A,\MeetOp_A;B,\MeetOp_B}}{%
    \FormulaRefAuto{LatticeHomMeetSemilatticeHom}[thm]{1}}
  \proofstepwidestar[2]{\SemilatticeHom{g;B,\MeetOp_B;C,\MeetOp_C}}{%
    \FormulaRefAuto{LatticeHomMeetSemilatticeHom}[thm]{2}}
  \proofstepwidestar[1,2]{%
    \SemilatticeHom{g\circ f;A,\MeetOp_A;C,\MeetOp_C}}{%
    \FormulaRefAuto{SemilatticeHomComposition}[thm]{3,4}}
  \proofstepwidestar[1]{\SemilatticeHom{f;A,\JoinOp_A;B,\JoinOp_B}}{%
    \FormulaRefAuto{LatticeHomJoinSemilatticeHom}[thm]{1}}
  \proofstepwidestar[2]{\SemilatticeHom{g;B,\JoinOp_B;C,\JoinOp_C}}{%
    \FormulaRefAuto{LatticeHomJoinSemilatticeHom}[thm]{2}}
  \proofstepwidestar[1,2]{%
    \SemilatticeHom{g\circ f;A,\JoinOp_A;C,\JoinOp_C}}{%
    \FormulaRefAuto{SemilatticeHomComposition}[thm]{6,7}}
  \proofstepwidestar[1]{\LatticeAlg{A,\MeetOp_A,\JoinOp_A}}{%
    \FormulaRefAuto{LatticeHomSourceAxiom}[ax]{1}}
  \proofstepwidestar[2]{\LatticeAlg{C,\MeetOp_C,\JoinOp_C}}{%
    \FormulaRefAuto{LatticeHomTargetAxiom}[ax]{2}}
  \proofstepwidestar[1,2]{g\circ f\colon A\to C}{%
    \FormulaRefAuto{SemilatticeHomMapAxiom}[ax]{5}}
  \proofstepwidestar[12]{x\in A}{\rA}
  \proofstepwidestar[13]{y\in A}{\rA}
  \proofstepwidestar[1,2,12,13]{%
    (g\circ f)(x\MeetOp_A{}y)
      =(g\circ f)(x)\MeetOp_C(g\circ f)(y)}{%
    \FormulaRefAuto{SemilatticeHomEvaluationAxiom}[ax]{5,12,13}}
  \proofstepwidestar[1,2,12,13]{%
    (g\circ f)(x\JoinOp_A{}y)
      =(g\circ f)(x)\JoinOp_C(g\circ f)(y)}{%
    \FormulaRefAuto{SemilatticeHomEvaluationAxiom}[ax]{8,12,13}}
  \proofstepwidestar[1,2]{%
    \LatticeHom{g\circ f;A,\MeetOp_A,\JoinOp_A;C,\MeetOp_C,\JoinOp_C}}{%
    \FormulaRefAuto{LatticeHomomorphism}[def]{9,10,11,14,15}}
\end{tabproofwide}

\FormulaThmDeltaK[Die Identität ist ein Verbandsisomorphismus]{%
  \LatticeAlg{A,\MeetOp_A,\JoinOp_A}\vdash
  \LatticeIso{\Id_A;A,\MeetOp_A,\JoinOp_A;A,\MeetOp_A,\JoinOp_A}%
}{LatticeIdentityIsomorphism}{%
  \DeltaRow{Mengen}{A}
  \DeltaRow{Funktionen}{\Id_A\dsep\MeetOp_A\dsep\JoinOp_A}
}
\begin{tabproof}
  \proofstep{1}{\LatticeAlg{A,\MeetOp_A,\JoinOp_A}}{\rA}
  \proofstep{1}{%
    \LatticeHom{\Id_A;A,\MeetOp_A,\JoinOp_A;A,\MeetOp_A,\JoinOp_A}}{%
    \FormulaRefAuto{LatticeIdentityHomomorphism}[thm]{1}}
  \proofstep{}{\Id_A\colon A\bij A}{%
    \FormulaRefAuto{\Id_A\colon A\bij A}}
  \proofstep{1}{%
    \LatticeIso{\Id_A;A,\MeetOp_A,\JoinOp_A;A,\MeetOp_A,\JoinOp_A}}{%
    \FormulaRefAuto{LatticeIsomorphism}[def]{2,3}}
\end{tabproof}

\FormulaThmDeltaK[Komposition von Verbandsisomorphismen]{%
  \begin{aligned}[t]
    &\LatticeIso{f;A,\MeetOp_A,\JoinOp_A;B,\MeetOp_B,\JoinOp_B}
      \dsep
      \LatticeIso{g;B,\MeetOp_B,\JoinOp_B;C,\MeetOp_C,\JoinOp_C}
      \vdash{}\\
    &\LatticeIso{g\circ f;A,\MeetOp_A,\JoinOp_A;C,\MeetOp_C,\JoinOp_C}
  \end{aligned}
}{LatticeIsoComposition}{%
  \DeltaRow{Mengen}{A\dsep B\dsep C}
  \DeltaRow{Funktionen}{%
    f\dsep g\dsep
    \MeetOp_A\dsep\JoinOp_A\dsep
    \MeetOp_B\dsep\JoinOp_B\dsep
    \MeetOp_C\dsep\JoinOp_C}
}
\begin{tabproofwide}
  \proofstepwidestar[1]{%
    \LatticeIso{f;A,\MeetOp_A,\JoinOp_A;B,\MeetOp_B,\JoinOp_B}}{\rA}
  \proofstepwidestar[2]{%
    \LatticeIso{g;B,\MeetOp_B,\JoinOp_B;C,\MeetOp_C,\JoinOp_C}}{\rA}
  \proofstepwidestar[1]{%
    \LatticeHom{f;A,\MeetOp_A,\JoinOp_A;B,\MeetOp_B,\JoinOp_B}}{%
    \FormulaRefAuto{LatticeIsoHomAxiom}[ax]{1}}
  \proofstepwidestar[2]{%
    \LatticeHom{g;B,\MeetOp_B,\JoinOp_B;C,\MeetOp_C,\JoinOp_C}}{%
    \FormulaRefAuto{LatticeIsoHomAxiom}[ax]{2}}
  \proofstepwidestar[1,2]{%
    \LatticeHom{g\circ f;A,\MeetOp_A,\JoinOp_A;C,\MeetOp_C,\JoinOp_C}}{%
    \FormulaRefAuto{LatticeHomComposition}[thm]{3,4}}
  \proofstepwidestar[1]{f\colon A\bij B}{%
    \FormulaRefAuto{LatticeIsoBijectiveAxiom}[ax]{1}}
  \proofstepwidestar[2]{g\colon B\bij C}{%
    \FormulaRefAuto{LatticeIsoBijectiveAxiom}[ax]{2}}
  \proofstepwidestar[1,2]{g\circ f\colon A\bij C}{%
    \FormulaRefAuto{BijectiveFunctionComposition}[thm]{6,7}}
  \proofstepwidestar[1,2]{%
    \LatticeIso{g\circ f;A,\MeetOp_A,\JoinOp_A;C,\MeetOp_C,\JoinOp_C}}{%
    \FormulaRefAuto{LatticeIsomorphism}[def]{5,8}}
\end{tabproofwide}

\section{Die gemeinsame Verbandsordnung}

\FormulaThmDeltaK[Übereinstimmung der beiden Verbandsordnungen]{%
  \begin{aligned}[t]
    &\LatticeAlg{A,\MeetOp,\JoinOp}\dsep x,y\in A\vdash{}\\
    &x\MeetOp{}y=x\leftrightarrow x\JoinOp{}y=y
  \end{aligned}
}{LatticeOrdersCoincide}{%
  \DeltaRow{Mengen}{A\dsep x\dsep y}
  \DeltaRow{Funktionen}{\MeetOp\dsep \JoinOp}
}
\begin{tabproofsplitwide}
  \proofpartwide{\(\vdash\)}
    \proofstepwidestar[1]{\LatticeAlg{A,\MeetOp,\JoinOp}}{\rA}
    \proofstepwidestar[2]{x\in A}{\rA}
    \proofstepwidestar[3]{y\in A}{\rA}
    \proofstepwidestar[4]{x\MeetOp{}y=x}{\rA}
    \proofstepwidestar[1]{\Semi{A,\MeetOp}}{%
      \FormulaRefAuto{LatticeMeetSemilatticeAxiom}[ax]{1}}
    \proofstepwidestar[1]{\Semi{A,\JoinOp}}{%
      \FormulaRefAuto{LatticeJoinSemilatticeAxiom}[ax]{1}}
    \proofstepwidestar[1,2,3]{y\MeetOp{}x=x\MeetOp{}y}{%
      \FormulaRefAuto{SemilatticeCommutativityAxiom}[ax]{5,3,2}}
    \proofstepwidestar[1,2,3,4]{y\MeetOp{}x=x}{\rIE{4,7}}
    \proofstepwidestar[1,2,3]{y\JoinOp\bigl(y\MeetOp{}x\bigr)=y}{%
      \FormulaRefAuto{LatticeJoinAbsorptionAxiom}[ax]{1,3,2}}
    \proofstepwidestar[1,2,3,4]{y\JoinOp{}x=y}{\rIE{8,9}}
    \proofstepwidestar[1,2,3]{y\JoinOp{}x=x\JoinOp{}y}{%
      \FormulaRefAuto{SemilatticeCommutativityAxiom}[ax]{6,3,2}}
    \proofstepwidestar[1,2,3,4]{x\JoinOp{}y=y}{%
      \FormulaRefAuto{a=b,\,a=c\vdash b=c}{11,10}}
  \closeproofpartwide

  \proofpartwide{\(\dashv\)}
    \proofstepwidestar[1]{\LatticeAlg{A,\MeetOp,\JoinOp}}{\rA}
    \proofstepwidestar[2]{x\in A}{\rA}
    \proofstepwidestar[3]{y\in A}{\rA}
    \proofstepwidestar[4]{x\JoinOp{}y=y}{\rA}
    \proofstepwidestar[1]{\Semi{A,\MeetOp}}{%
      \FormulaRefAuto{LatticeMeetSemilatticeAxiom}[ax]{1}}
    \proofstepwidestar[1]{\Semi{A,\JoinOp}}{%
      \FormulaRefAuto{LatticeJoinSemilatticeAxiom}[ax]{1}}
    \proofstepwidestar[1,2,3]{x\MeetOp\bigl(x\JoinOp{}y\bigr)=x}{%
      \FormulaRefAuto{LatticeMeetAbsorptionAxiom}[ax]{1,2,3}}
    \proofstepwidestar[1,2,3,4]{x\MeetOp{}y=x}{\rIE{4,7}}
  \closeproofpartwide
\end{tabproofsplitwide}

Bezeichnen wir die dadurch gemeinsame Relation mit \(\leq\), so gilt in
jedem Verband die charakteristische Dreifachäquivalenz
\[
  x\InducedLeq y
  \quad\Longleftrightarrow\quad
  x\MeetOp{}y=x
  \quad\Longleftrightarrow\quad
  x\JoinOp{}y=y.
\]

\section{Beispiele von Verbänden}

\subsection{Das Potenzmengengitter}

Die beiden Halbverbandsbeispiele aus Kapitel~2 werden nun durch die getrennt
zu beweisenden Absorptionsgesetze zu einem Verband verbunden.

\FormulaThmDeltaK[Absorption des Durchschnitts über der Vereinigung]{%
  X,Y\in\powerset(S)\vdash X\cap(X\cup Y)=X
}{PowerSetMeetAbsorption}{%
  \DeltaRow{Mengen}{S\dsep X\dsep Y}
}
\begin{tabproof}
  \proofstep{1}{X\in\powerset(S)}{\rA}
  \proofstep{2}{Y\in\powerset(S)}{\rA}
  \proofstep{}{X\subseteq X\cup Y}{\FormulaRefAuto{A \subseteq A \cup B}}
  \proofstep{}{X\cap(X\cup Y)=X}{%
    \FormulaRefAuto{A \subseteq B \eqvdash A \cap B = A}{3}}
\end{tabproof}

\FormulaThmDeltaK[Absorption der Vereinigung über dem Durchschnitt]{%
  X,Y\in\powerset(S)\vdash X\cup(X\cap Y)=X
}{PowerSetJoinAbsorption}{%
  \DeltaRow{Mengen}{S\dsep X\dsep Y}
}
\begin{tabproof}
  \proofstep{1}{X\in\powerset(S)}{\rA}
  \proofstep{2}{Y\in\powerset(S)}{\rA}
  \proofstep{}{X\cap Y\subseteq X}{\FormulaRefAuto{A \cap B \subseteq A}}
  \proofstep{}{X\subseteq X}{\FormulaRefAuto{A\subseteq A}}
  \proofstep{}{X\cup(X\cap Y)\subseteq X}{%
    \FormulaRefAuto{A \subseteq C,\, B \subseteq C \vdash A \cup B \subseteq C}{4,3}}
  \proofstep{}{X\subseteq X\cup(X\cap Y)}{%
    \FormulaRefAuto{A \subseteq A \cup B}}
  \proofstep{}{X\cup(X\cap Y)=X}{%
    \FormulaRefAuto{A \subseteq B \land B \subseteq A \eqvdash A = B}{5,6}}
\end{tabproof}

\FormulaThmDeltaK[Potenzmengen bilden Verbände]{%
  \LatticeAlg{\powerset(S),\cap,\cup}
}{PowerSetLattice}{%
  \DeltaRow{Mengen}{S\dsep X\dsep Y\dsep Z}
  \DeltaRow{Funktionen}{\cap\dsep\cup}
}
\begin{tabproofsplitwide}
  \proofpartwide{Axiome}
    \proofstepwidestar[1]{X\in\powerset(S)}{\rA}
    \proofstepwidestar[2]{Y\in\powerset(S)}{\rA}
    \proofstepwidestar[3]{Z\in\powerset(S)}{\rA}
    \proofstepwidestar[1,2]{X\cap Y\in\powerset(S)}{%
      \FormulaRefAuto{PowerSetIntersectionClosure}[thm]{1,2}}
    \proofstepwidestar[1,2,3]{(X\cap Y)\cap Z=X\cap(Y\cap Z)}{%
      \FormulaRefAuto{(A \cap B) \cap C = A \cap (B \cap C)}}
    \proofstepwidestar[1,2]{X\cap Y=Y\cap X}{%
      \FormulaRefAuto{A \cap B = B \cap A}}
    \proofstepwidestar[1]{X\cap X=X}{%
      \FormulaRefAuto{a=b\vdash b=a}{\FormulaRefAuto{A = A \cap A}}}
    \proofstepwidestar[1,2]{X\cup Y\in\powerset(S)}{%
      \FormulaRefAuto{PowerSetUnionClosure}[thm]{1,2}}
    \proofstepwidestar[1,2,3]{(X\cup Y)\cup Z=X\cup(Y\cup Z)}{%
      \FormulaRefAuto{(A \cup B) \cup C = A \cup (B \cup C)}}
    \proofstepwidestar[1,2]{X\cup Y=Y\cup X}{%
      \FormulaRefAuto{A \cup B = B \cup A}}
    \proofstepwidestar[1]{X\cup X=X}{%
      \FormulaRefAuto{A \cup A = A}}
    \proofstepwidestar[1,2]{X\cap(X\cup Y)=X}{%
      \FormulaRefAuto{PowerSetMeetAbsorption}[thm]{1,2}}
    \proofstepwidestar[1,2]{X\cup(X\cap Y)=X}{%
      \FormulaRefAuto{PowerSetJoinAbsorption}[thm]{1,2}}
    \proofstepwidestar[]{\LatticeAlg{\powerset(S),\cap,\cup}}{%
      \FormulaRefAuto{AlgebraicLattice}[def]{4,5,6,7,8,9,10,11,12,13}}
  \closeproofpartwide
\end{tabproofsplitwide}

Nach den beiden bereits bewiesenen Ordnungscharakterisierungen ist die
gemeinsame Verbandsordnung genau die Inklusion:
\[
  X\subseteq Y
  \quad\Longleftrightarrow\quad
  X\cap Y=X
  \quad\Longleftrightarrow\quad
  X\cup Y=Y.
\]


\chapter{Supremums-Halbverbände mit Null und Eins}

\section{Supremums-Halbverbände mit Null}

Ein Supremums-Halbverband mit Null besitzt ein kleinstes Element und ist
damit nach unten beschränkt. Ein größtes Element folgt daraus im Allgemeinen
nicht; beispielsweise besitzt der Supremums-Halbverband
\((\mathbb{N},\max,0)\) kein größtes Element. Deshalb bezeichnet
\(\ZeroJoinSemi{A,\JoinOp,0}\) in diesem Band präzise einen
Supremums-Halbverband mit Null und nicht schon einen beschränkten
Supremums-Halbverband.

Im gesamten Kapitel bezeichnet \(\InducedLeq\) jeweils die von der im
Kontext genannten Halbverbandsoperation supremal induzierten Relationsoperator.
Operation und Träger werden im Metablock oder im jeweiligen
Strukturprädikat festgehalten und daher nicht an das Relationssymbol
geschrieben.

Zusätzlich zu den Halbverbandsaxiomen werden die Trägerzugehörigkeit des
Nullelements und sein Neutralitätsaxiom als eigene Strukturaxiome festgehalten.

\FormulaAxiomDeltaK[Halbverbandsaxiom einer Nullelementstruktur]{%
  \Semi{A,\JoinOp}%
}{ZeroJoinBaseSemilatticeLaw}{%
  \DeltaRow{Mengen}{A\dsep 0}
  \DeltaRow{Funktionen}{\JoinOp}
}
\par\smallskip

\FormulaAxiomDeltaK[Trägeraxiom eines Nullelements]{%
  0\in A%
}{ZeroJoinCarrierLaw}{%
  \DeltaRow{Mengen}{A\dsep 0}
  \DeltaRow{Funktionen}{\JoinOp}
}
\par\smallskip

\FormulaAxiomDeltaK[Neutralitätsaxiom eines Nullelements]{%
  x\in A\vdash 0\JoinOp{}x=x%
}{ZeroJoinIdentityLaw}{%
  \DeltaRow{Mengen}{A\dsep x\dsep 0}
  \DeltaRow{Funktionen}{\JoinOp}
}

\FormulaDefDeltaK[Nach unten beschränkter Supremums-Halbverband mit Null]{%
  \ZeroJoinSemi{A,\JoinOp,0}%
}{ZeroJoinSemilattice}{%
  \DeltaRow{Mengen}{A\dsep x\dsep 0}
  \DeltaRow{Funktionen}{\JoinOp}
  \DeltaRow{\textbf{Axiome}}{}
  \DeltaRow{Halbverband}
           {\Semi{A,\JoinOp}}
           [\FormulaRefAuto{ZeroJoinBaseSemilatticeLaw}[ax]]
  \DeltaRow{Nullelement}
           {0\in A}
           [\FormulaRefAuto{ZeroJoinCarrierLaw}[ax]]
  \DeltaRow{Neutralitätsaxiom}
           {x\in A\vdash 0\JoinOp{}x=x}
           [\FormulaRefAuto{ZeroJoinIdentityLaw}[ax]]
  \DeltaRow{\textbf{Neues Symbol}}{}
  \DeltaPrem{\makecell[l]{Supremums-Halbverbände\\mit Nullelement}}{%
    \ZeroJoinSemi{A,\JoinOp,0}}
}

\newpage
\FormulaAxiomDeltaK[Zugriff auf das Halbverbandsaxiom der Nullelementstruktur]{%
  \ZeroJoinSemi{A,\JoinOp,0}\vdash \Semi{A,\JoinOp}%
}{ZeroJoinBaseSemilatticeAxiom}{%
  \DeltaRow{Mengen}{A\dsep 0}
  \DeltaRow{Funktionen}{\JoinOp}
}
\par\smallskip

\FormulaAxiomDeltaK[Zugriff auf das Nullelement im Träger]{%
  \ZeroJoinSemi{A,\JoinOp,0}\vdash 0\in A%
}{ZeroJoinCarrierAxiom}{%
  \DeltaRow{Mengen}{A\dsep 0}
  \DeltaRow{Funktionen}{\JoinOp}
}

\FormulaAxiomDeltaK[Zugriff auf das Neutralitätsaxiom des Nullelements]{%
  \ZeroJoinSemi{A,\JoinOp,0}\dsep x\in A\vdash 0\JoinOp{}x=x%
}{ZeroJoinIdentityAxiom}{%
  \DeltaRow{Mengen}{A\dsep x\dsep 0}
  \DeltaRow{Funktionen}{\JoinOp}
}
\leavevmode\par\medskip

\FormulaThmDeltaK[Das Nullelement ist das kleinste Element]{%
  \ZeroJoinSemi{A,\JoinOp,0}\dsep x\in A
  \vdash 0\InducedLeq x
}{ZeroJoinLeast}{%
  \DeltaRow{Mengen}{A\dsep x\dsep 0}
  \DeltaRow{Funktionen}{\JoinOp}
}
\begin{tabproof}
  \proofstep{1}{\ZeroJoinSemi{A,\JoinOp,0}}{\rA}
  \proofstep{2}{x\in A}{\rA}
  \proofstep{1}{0\in A}{\FormulaRefAuto{ZeroJoinCarrierAxiom}[ax]{1}}
  \proofstep{1,2}{0\JoinOp{}x=x}{\FormulaRefAuto{ZeroJoinIdentityAxiom}[ax]{1,2}}
  \proofstep{1,2}{%
    0\InducedLeq x\leftrightarrow 0\JoinOp{}x=x}{%
    \FormulaRefAuto{JoinOrderEvaluationAxiom}[ax]{%
      \FormulaRefAuto{ZeroJoinBaseSemilatticeAxiom}[ax]{1},3,2}}
  \proofstep{1,2}{0\InducedLeq x}{%
    \FormulaRefAuto{P\leftrightarrow Q, Q\vdash P}{5,4}}
\end{tabproof}

\begin{remark}
Der vorherige Satz rechtfertigt auch die Bezeichnung
\emph{nach unten beschränkter Supremums-Halbverband}. Das Wort
\emph{beschränkt} wird im Folgenden nur verwendet, wenn zusätzlich ein
größtes Element ausgezeichnet ist.
\end{remark}

\section{Beschränkte Supremums-Halbverbände}

Ein beschränkter Supremums-Halbverband besitzt neben seiner Null eine
ausgezeichnete Eins. Algebraisch wird die Eins dadurch charakterisiert, dass
sie jedes Element unter dem Supremum absorbiert.
Die bereits definierte Nullstruktur wird dabei unverändert wiederverwendet
und deshalb nicht ein zweites Mal als eigenes Axiom registriert.

\FormulaAxiomDeltaK[Trägeraxiom der Eins]{%
  1\in A%
}{BoundedJoinTopCarrierLaw}{%
  \DeltaRow{Mengen}{A\dsep 0\dsep 1}
  \DeltaRow{Funktionen}{\JoinOp}
}
\par\smallskip

\FormulaAxiomDeltaK[Absorptionsaxiom der Eins]{%
  x\in A\vdash x\JoinOp{}1=1%
}{BoundedJoinTopAbsorptionLaw}{%
  \DeltaRow{Mengen}{A\dsep x\dsep 0\dsep 1}
  \DeltaRow{Funktionen}{\JoinOp}
}

\FormulaDefDeltaK[Beschränkter Supremums-Halbverband mit Null und Eins]{%
  \BoundedJoinSemi{A,\JoinOp,0,1}%
}{BoundedJoinSemilattice}{%
  \DeltaRow{Mengen}{A\dsep x\dsep 0\dsep 1}
  \DeltaRow{Funktionen}{\JoinOp}
  \DeltaRow{\textbf{Axiome}}{}
  \DeltaRow{Supremums-Halbverband mit Null}
           {\ZeroJoinSemi{A,\JoinOp,0}}
           [\FormulaRefAuto{ZeroJoinSemilattice}[def]]
  \DeltaRow{Eins im Träger}
           {1\in A}
           [\FormulaRefAuto{BoundedJoinTopCarrierLaw}[ax]]
  \DeltaRow{Absorption durch Eins}
           {x\in A\vdash x\JoinOp{}1=1}
           [\FormulaRefAuto{BoundedJoinTopAbsorptionLaw}[ax]]
  \DeltaRow{\textbf{Neues Symbol}}{}
  \DeltaPrem{\makecell[l]{Beschränkte\\Supremums-Halbverbände\\
                          mit Null und Eins}}{%
    \BoundedJoinSemi{A,\JoinOp,0,1}}
}

\FormulaAxiomDeltaK[Zugriff auf die Nullstruktur eines beschränkten Supremums-Halbverbandes]{%
  \BoundedJoinSemi{A,\JoinOp,0,1}
  \vdash \ZeroJoinSemi{A,\JoinOp,0}%
}{BoundedJoinZeroSemilatticeAxiom}{%
  \DeltaRow{Mengen}{A\dsep 0\dsep 1}
  \DeltaRow{Funktionen}{\JoinOp}
}
\par\smallskip

\FormulaAxiomDeltaK[Zugriff auf die Eins im Träger]{%
  \BoundedJoinSemi{A,\JoinOp,0,1}\vdash 1\in A%
}{BoundedJoinTopCarrierAxiom}{%
  \DeltaRow{Mengen}{A\dsep 0\dsep 1}
  \DeltaRow{Funktionen}{\JoinOp}
}
\par\smallskip

\FormulaAxiomDeltaK[Zugriff auf das Absorptionsaxiom der Eins]{%
  \BoundedJoinSemi{A,\JoinOp,0,1}\dsep x\in A
  \vdash x\JoinOp{}1=1%
}{BoundedJoinTopAbsorptionAxiom}{%
  \DeltaRow{Mengen}{A\dsep x\dsep 0\dsep 1}
  \DeltaRow{Funktionen}{\JoinOp}
}
\leavevmode\par\medskip

\subsection{Spezielle Definitionen}

\FormulaDefDeltaK[\(\mathsf{SupHV}_0\)-Homomorphismus]{%
  \ZeroJoinSemiHom{f;\mathcal A_0;\mathcal B_0}%
}{ZeroJoinSemilatticeHomomorphism}{%
  \DeltaRow{Mengen}{A\dsep B\dsep 0_A\dsep 0_B}
  \DeltaRow{Funktionen}{f\dsep\JoinOp\dsep\mathbin{\diamond}}
  \DeltaRow{Strukturkürzel}{%
    \begin{aligned}[t]
      \mathcal A_0&\coloneqq(A,\JoinOp,0_A),\\[-2pt]
      \mathcal B_0&\coloneqq(B,\mathbin{\diamond},0_B)
    \end{aligned}}
  \DeltaPrem{Quellstruktur}{\ZeroJoinSemi{A,\JoinOp,0_A}}
  \DeltaPrem{Zielstruktur}{\ZeroJoinSemi{B,\mathbin{\diamond},0_B}}
  \DeltaPrem{Halbverbandshomomorphismus}{%
    \SemilatticeHom{f;A,\JoinOp;B,\mathbin{\diamond}}}
  \DeltaRow{Erhaltung der Null}{f(0_A)=0_B}
  \DeltaRow{\textbf{Neues Symbol}}{}
  \DeltaPrem{\(\mathsf{SupHV}_0\)-Homomorphismen}{%
    \ZeroJoinSemiHom{f;\mathcal A_0;\mathcal B_0}}
}

\FormulaAxiomDeltaK[Quellstruktur eines \(\mathsf{SupHV}_0\)-Homomorphismus]{%
  \ZeroJoinSemiHom{f;A,\JoinOp,0_A;B,\mathbin{\diamond},0_B}
  \vdash \ZeroJoinSemi{A,\JoinOp,0_A}%
}{ZeroJoinSemilatticeHomSourceAxiom}{%
  \DeltaRow{Mengen}{A\dsep B\dsep 0_A\dsep 0_B}
  \DeltaRow{Funktionen}{f\dsep\JoinOp\dsep\mathbin{\diamond}}
}

\FormulaAxiomDeltaK[Zielstruktur eines \(\mathsf{SupHV}_0\)-Homomorphismus]{%
  \ZeroJoinSemiHom{f;A,\JoinOp,0_A;B,\mathbin{\diamond},0_B}
  \vdash \ZeroJoinSemi{B,\mathbin{\diamond},0_B}%
}{ZeroJoinSemilatticeHomTargetAxiom}{%
  \DeltaRow{Mengen}{A\dsep B\dsep 0_A\dsep 0_B}
  \DeltaRow{Funktionen}{f\dsep\JoinOp\dsep\mathbin{\diamond}}
}

\FormulaAxiomDeltaK[Halbverbandshomomorphismusanteil eines
  \(\mathsf{SupHV}_0\)-Homomorphismus]{%
  \ZeroJoinSemiHom{f;A,\JoinOp,0_A;B,\mathbin{\diamond},0_B}
  \vdash \SemilatticeHom{f;A,\JoinOp;B,\mathbin{\diamond}}%
}{ZeroJoinSemilatticeHomBaseHomAxiom}{%
  \DeltaRow{Mengen}{A\dsep B\dsep 0_A\dsep 0_B}
  \DeltaRow{Funktionen}{f\dsep\JoinOp\dsep\mathbin{\diamond}}
}

\FormulaAxiomDeltaK[Nullerhaltung eines \(\mathsf{SupHV}_0\)-Homomorphismus]{%
  \ZeroJoinSemiHom{f;A,\JoinOp,0_A;B,\mathbin{\diamond},0_B}
  \vdash f(0_A)=0_B%
}{ZeroJoinSemilatticeHomZeroPreservationAxiom}{%
  \DeltaRow{Mengen}{A\dsep B\dsep 0_A\dsep 0_B}
  \DeltaRow{Funktionen}{f\dsep\JoinOp\dsep\mathbin{\diamond}}
}

\FormulaDefDeltaK[\(\mathsf{SupHV}_0\)-Isomorphismus]{%
  \ZeroJoinSemiIso{f;\mathcal A_0;\mathcal B_0}%
}{ZeroJoinSemilatticeIsomorphism}{%
  \DeltaRow{Mengen}{A\dsep B\dsep 0_A\dsep 0_B}
  \DeltaRow{Funktionen}{f\dsep\JoinOp\dsep\mathbin{\diamond}}
  \DeltaRow{Strukturkürzel}{%
    \begin{aligned}[t]
      \mathcal A_0&\coloneqq(A,\JoinOp,0_A),\\[-2pt]
      \mathcal B_0&\coloneqq(B,\mathbin{\diamond},0_B)
    \end{aligned}}
  \DeltaPrem{Homomorphismus}{%
    \ZeroJoinSemiHom{f;\mathcal A_0;\mathcal B_0}}
  \DeltaPrem{Bijektivität}{f\colon A\bij B}
  \DeltaRow{\textbf{Neues Symbol}}{}
  \DeltaPrem{\(\mathsf{SupHV}_0\)-Isomorphismen}{%
    \ZeroJoinSemiIso{f;\mathcal A_0;\mathcal B_0}}
}

\FormulaAxiomDeltaK[Homomorphismusanteil eines
  \(\mathsf{SupHV}_0\)-Isomorphismus]{%
  \ZeroJoinSemiIso{f;A,\JoinOp,0_A;B,\mathbin{\diamond},0_B}
  \vdash
  \ZeroJoinSemiHom{f;A,\JoinOp,0_A;B,\mathbin{\diamond},0_B}%
}{ZeroJoinSemilatticeIsoHomAxiom}{%
  \DeltaRow{Mengen}{A\dsep B\dsep 0_A\dsep 0_B}
  \DeltaRow{Funktionen}{f\dsep\JoinOp\dsep\mathbin{\diamond}}
}

\FormulaAxiomDeltaK[Bijektivität eines
  \(\mathsf{SupHV}_0\)-Isomorphismus]{%
  \ZeroJoinSemiIso{f;A,\JoinOp,0_A;B,\mathbin{\diamond},0_B}
  \vdash f\colon A\bij B%
}{ZeroJoinSemilatticeIsoBijectiveAxiom}{%
  \DeltaRow{Mengen}{A\dsep B\dsep 0_A\dsep 0_B}
  \DeltaRow{Funktionen}{f\dsep\JoinOp\dsep\mathbin{\diamond}}
}

\FormulaDefDeltaK[\(\mathsf{SupHV}_{0,1}\)-Homomorphismus]{%
  \BoundedJoinSemiHom{f;\mathcal A_{01};\mathcal B_{01}}%
}{BoundedJoinSemilatticeHomomorphism}{%
  \DeltaRow{Mengen}{A\dsep B\dsep 0_A\dsep 0_B\dsep 1_A\dsep 1_B}
  \DeltaRow{Funktionen}{f\dsep\JoinOp\dsep\mathbin{\diamond}}
  \DeltaRow{Strukturkürzel}{%
    \begin{aligned}[t]
      \mathcal A_0&\coloneqq(A,\JoinOp,0_A),
        &\mathcal A_{01}&\coloneqq(A,\JoinOp,0_A,1_A),\\[-2pt]
      \mathcal B_0&\coloneqq(B,\mathbin{\diamond},0_B),
        &\mathcal B_{01}&\coloneqq(B,\mathbin{\diamond},0_B,1_B)
    \end{aligned}}
  \DeltaPrem{Quellstruktur}{\BoundedJoinSemi{A,\JoinOp,0_A,1_A}}
  \DeltaPrem{Zielstruktur}{%
    \BoundedJoinSemi{B,\mathbin{\diamond},0_B,1_B}}
  \DeltaPrem{Homomorphismus mit Null}{%
    \ZeroJoinSemiHom{f;\mathcal A_0;\mathcal B_0}}
  \DeltaRow{Erhaltung der Eins}{f(1_A)=1_B}
  \DeltaRow{\textbf{Neues Symbol}}{}
  \DeltaPrem{\(\mathsf{SupHV}_{0,1}\)-Homomorphismen}{%
    \BoundedJoinSemiHom{f;\mathcal A_{01};\mathcal B_{01}}}
}

\FormulaAxiomDeltaK[Quellstruktur eines
  \(\mathsf{SupHV}_{0,1}\)-Homomorphismus]{%
  \BoundedJoinSemiHom{
    f;A,\JoinOp,0_A,1_A;B,\mathbin{\diamond},0_B,1_B}
  \vdash \BoundedJoinSemi{A,\JoinOp,0_A,1_A}%
}{BoundedJoinSemilatticeHomSourceAxiom}{%
  \DeltaRow{Mengen}{A\dsep B\dsep 0_A\dsep 0_B\dsep 1_A\dsep 1_B}
  \DeltaRow{Funktionen}{f\dsep\JoinOp\dsep\mathbin{\diamond}}
}

\FormulaAxiomDeltaK[Zielstruktur eines
  \(\mathsf{SupHV}_{0,1}\)-Homomorphismus]{%
  \BoundedJoinSemiHom{
    f;A,\JoinOp,0_A,1_A;B,\mathbin{\diamond},0_B,1_B}
  \vdash \BoundedJoinSemi{B,\mathbin{\diamond},0_B,1_B}%
}{BoundedJoinSemilatticeHomTargetAxiom}{%
  \DeltaRow{Mengen}{A\dsep B\dsep 0_A\dsep 0_B\dsep 1_A\dsep 1_B}
  \DeltaRow{Funktionen}{f\dsep\JoinOp\dsep\mathbin{\diamond}}
}

\FormulaAxiomDeltaK[Basishomomorphismus eines
  \(\mathsf{SupHV}_{0,1}\)-Homomorphismus]{%
  \BoundedJoinSemiHom{
    f;A,\JoinOp,0_A,1_A;B,\mathbin{\diamond},0_B,1_B}
  \vdash
  \ZeroJoinSemiHom{f;A,\JoinOp,0_A;B,\mathbin{\diamond},0_B}%
}{BoundedJoinSemilatticeHomBaseHomAxiom}{%
  \DeltaRow{Mengen}{A\dsep B\dsep 0_A\dsep 0_B\dsep 1_A\dsep 1_B}
  \DeltaRow{Funktionen}{f\dsep\JoinOp\dsep\mathbin{\diamond}}
}

\FormulaAxiomDeltaK[Einserhaltung eines
  \(\mathsf{SupHV}_{0,1}\)-Homomorphismus]{%
  \BoundedJoinSemiHom{
    f;A,\JoinOp,0_A,1_A;B,\mathbin{\diamond},0_B,1_B}
  \vdash f(1_A)=1_B%
}{BoundedJoinSemilatticeHomOnePreservationAxiom}{%
  \DeltaRow{Mengen}{A\dsep B\dsep 0_A\dsep 0_B\dsep 1_A\dsep 1_B}
  \DeltaRow{Funktionen}{f\dsep\JoinOp\dsep\mathbin{\diamond}}
}

\FormulaDefDeltaK[\(\mathsf{SupHV}_{0,1}\)-Isomorphismus]{%
  \BoundedJoinSemiIso{f;\mathcal A_{01};\mathcal B_{01}}%
}{BoundedJoinSemilatticeIsomorphism}{%
  \DeltaRow{Mengen}{A\dsep B\dsep 0_A\dsep 0_B\dsep 1_A\dsep 1_B}
  \DeltaRow{Funktionen}{f\dsep\JoinOp\dsep\mathbin{\diamond}}
  \DeltaRow{Strukturkürzel}{%
    \begin{aligned}[t]
      \mathcal A_{01}&\coloneqq(A,\JoinOp,0_A,1_A),\\[-2pt]
      \mathcal B_{01}&\coloneqq(B,\mathbin{\diamond},0_B,1_B)
    \end{aligned}}
  \DeltaPrem{Homomorphismus}{%
    \BoundedJoinSemiHom{f;\mathcal A_{01};\mathcal B_{01}}}
  \DeltaPrem{Bijektivität}{f\colon A\bij B}
  \DeltaRow{\textbf{Neues Symbol}}{}
  \DeltaPrem{\(\mathsf{SupHV}_{0,1}\)-Isomorphismen}{%
    \BoundedJoinSemiIso{f;\mathcal A_{01};\mathcal B_{01}}}
}

\FormulaAxiomDeltaK[Homomorphismusanteil eines
  \(\mathsf{SupHV}_{0,1}\)-Isomorphismus]{%
  \BoundedJoinSemiIso{
    f;A,\JoinOp,0_A,1_A;B,\mathbin{\diamond},0_B,1_B}
  \vdash
  \BoundedJoinSemiHom{
    f;A,\JoinOp,0_A,1_A;B,\mathbin{\diamond},0_B,1_B}%
}{BoundedJoinSemilatticeIsoHomAxiom}{%
  \DeltaRow{Mengen}{A\dsep B\dsep 0_A\dsep 0_B\dsep 1_A\dsep 1_B}
  \DeltaRow{Funktionen}{f\dsep\JoinOp\dsep\mathbin{\diamond}}
}

\FormulaAxiomDeltaK[Bijektivität eines
  \(\mathsf{SupHV}_{0,1}\)-Isomorphismus]{%
  \BoundedJoinSemiIso{
    f;A,\JoinOp,0_A,1_A;B,\mathbin{\diamond},0_B,1_B}
  \vdash f\colon A\bij B%
}{BoundedJoinSemilatticeIsoBijectiveAxiom}{%
  \DeltaRow{Mengen}{A\dsep B\dsep 0_A\dsep 0_B\dsep 1_A\dsep 1_B}
  \DeltaRow{Funktionen}{f\dsep\JoinOp\dsep\mathbin{\diamond}}
}

\FormulaThmDeltaK[Die Eins ist das größte Element]{%
  \BoundedJoinSemi{A,\JoinOp,0,1}\dsep x\in A
  \vdash x\InducedLeq 1%
}{BoundedJoinTopGreatest}{%
  \DeltaRow{Mengen}{A\dsep x\dsep 0\dsep 1}
  \DeltaRow{Funktionen}{\JoinOp}
}
\begin{tabproof}
  \proofstep{1}{\BoundedJoinSemi{A,\JoinOp,0,1}}{\rA}
  \proofstep{2}{x\in A}{\rA}
  \proofstep{1}{\ZeroJoinSemi{A,\JoinOp,0}}{%
    \FormulaRefAuto{BoundedJoinZeroSemilatticeAxiom}[ax]{1}}
  \proofstep{1}{\Semi{A,\JoinOp}}{%
    \FormulaRefAuto{ZeroJoinBaseSemilatticeAxiom}[ax]{3}}
  \proofstep{1}{1\in A}{%
    \FormulaRefAuto{BoundedJoinTopCarrierAxiom}[ax]{1}}
  \proofstep{1,2}{x\JoinOp{}1=1}{%
    \FormulaRefAuto{BoundedJoinTopAbsorptionAxiom}[ax]{1,2}}
  \proofstep{1,2}{%
    x\InducedLeq 1\leftrightarrow x\JoinOp{}1=1}{%
    \FormulaRefAuto{JoinOrderEvaluationAxiom}[ax]{4,2,5}}
  \proofstep{1,2}{x\InducedLeq 1}{%
    \FormulaRefAuto{P\leftrightarrow Q, Q\vdash P}{7,6}}
\end{tabproof}

\FormulaThmDeltaK[Ein größtes Element absorbiert beim Supremum]{%
  \begin{aligned}[t]
    &\ZeroJoinSemi{A,\JoinOp,0}\dsep e\in A\dsep
      \forall x\in A\,(x\InducedLeq e)\dsep x\in A\vdash{}\\
    &x\JoinOp{}e=e
  \end{aligned}
}{ZeroJoinGreatestAbsorbs}{%
  \DeltaRow{Mengen}{A\dsep e\dsep x\dsep 0}
  \DeltaRow{Funktionen}{\JoinOp}
}
\begin{tabproof}
  \proofstep{1}{\ZeroJoinSemi{A,\JoinOp,0}}{\rA}
  \proofstep{2}{e\in A}{\rA}
  \proofstep{3}{\forall x\in A\,(x\InducedLeq e)}{\rA}
  \proofstep{4}{x\in A}{\rA}
  \proofstep{1}{\Semi{A,\JoinOp}}{%
    \FormulaRefAuto{ZeroJoinBaseSemilatticeAxiom}[ax]{1}}
  \proofstep{3,4}{x\InducedLeq e}{\rRE{4,\rUE{3}}}
  \proofstep{1,3,4}{x\JoinOp{}e=e}{%
    \FormulaRefAuto{JoinOrderEvaluation}[thm]{5,6}}
\end{tabproof}

\section{Endliche Supremums-Halbverbände mit Null}

\FormulaDefDeltaK[\(\JoinOp\)-irreduzibles Element]{%
  \JoinIrr{A,\JoinOp,0}{j}%
}{JoinIrreducible}{%
  \DeltaRow{Mengen}{A\dsep 0\dsep j\dsep x\dsep y}
  \DeltaRow{Funktionen}{\JoinOp}
  \DeltaPrem{Halbverband mit Nullelement}{\ZeroJoinSemi{A,\JoinOp,0}}
  \DeltaRow{Von Null verschieden}{j\in A\dsep j\neq 0}
  \DeltaRow{Irreduzibilität}{%
    x,y\in A\dsep x\JoinOp{}y=j\vdash x=j\lor y=j}
  \DeltaRow{\textbf{Neues Symbol}}{}
  \DeltaPrem{\(\JoinOp\)-Irreduzible}{\JoinIrr{A,\JoinOp,0}{j}}
}
\leavevmode\par\medskip

\FormulaDefDeltaK[Menge der \(\JoinOp\)-Irreduziblen]{%
  \JoinIrrSet{A,\JoinOp,0}
  \coloneqq\{j\in A\mid\JoinIrr{A,\JoinOp,0}{j}\}
}{JoinIrreducibleSet}{%
  \DeltaRow{Mengen}{A\dsep j\dsep 0}
  \DeltaRow{Funktionen}{\JoinOp}
  \DeltaRow{\textbf{Neues Symbol}}{}
  \DeltaPrem{\(\JoinOp\)-Irreduzible}{\JoinIrrSet{A,\JoinOp,0}}
}

Für gemeinsame untere Schranken eines Elementpaares verwenden wir ohne
neues Symbol den in Band~11 definierten Spezialfall
\[
  \LB_{A,\InducedLeq}(\{x,y\}).
\]
Die allgemeine Charakterisierung dieser Menge und ihre Einbettung in den
Träger sind bereits in Band~11 bewiesen.

\FormulaThmDeltaK[Endlichkeit der gemeinsamen unteren Schranken]{%
  \Finite A\vdash\Finite{\LB_{A,\InducedLeq}(\{x,y\})}
}{PairLowerBoundSetFinite}{%
  \DeltaRow{Mengen}{A\dsep x\dsep y}
  \DeltaRow{Funktionen}{\JoinOp}
}
\begin{tabproof}
  \proofstep{1}{\Finite A}{\rA}
  \proofstep{}{\LB_{A,\InducedLeq}(\{x,y\})\subseteq A}{%
    \FormulaRefAuto{LowerBoundSetSubsetCarrier}[thm]}
  \proofstep{1}{\Finite{\LB_{A,\InducedLeq}(\{x,y\})}}{%
    \FormulaRefAuto{FiniteSubset}[thm]{2,1}}
\end{tabproof}

\FormulaThmDeltaK[Das Nullelement ist eine gemeinsame untere Schranke]{%
  \ZeroJoinSemi{A,\JoinOp,0}\dsep x,y\in A
  \vdash 0\in\LB_{A,\InducedLeq}(\{x,y\})
}{PairLowerBoundSetContainsZero}{%
  \DeltaRow{Mengen}{A\dsep 0\dsep x\dsep y}
  \DeltaRow{Funktionen}{\JoinOp}
}
\begin{tabproof}
  \proofstep{1}{\ZeroJoinSemi{A,\JoinOp,0}}{\rA}
  \proofstep{2}{x\in A}{\rA}
  \proofstep{3}{y\in A}{\rA}
  \proofstep{1}{0\in A}{\FormulaRefAuto{ZeroJoinCarrierAxiom}[ax]{1}}
  \proofstep{1,2}{0\InducedLeq x}{%
    \FormulaRefAuto{ZeroJoinLeast}[thm]{1,2}}
  \proofstep{1,3}{0\InducedLeq y}{%
    \FormulaRefAuto{ZeroJoinLeast}[thm]{1,3}}
  \proofstep{1,2,3}{%
    0\in A\land0\InducedLeq x\land
    0\InducedLeq y}{\rAI{\rAI{4,5},6}}
  \proofstep{1,2,3}{0\in\LB_{A,\InducedLeq}(\{x,y\})}{%
    \FormulaRefAuto{P\leftrightarrow Q, Q\vdash P}{%
      \FormulaRefAuto{PairLowerBoundSetCharacterization}[thm],7}}
\end{tabproof}

\FormulaThmDeltaK[Abschluss der gemeinsamen unteren Schranken]{%
  \begin{aligned}[t]
    &\ZeroJoinSemi{A,\JoinOp,0}\dsep x,y\in A\dsep
      d,e\in\LB_{A,\InducedLeq}(\{x,y\})\vdash{}\\
    &d\JoinOp{}e\in\LB_{A,\InducedLeq}(\{x,y\})
  \end{aligned}
}{PairLowerBoundSetJoinClosed}{%
  \DeltaRow{Mengen}{A\dsep 0\dsep d\dsep e\dsep x\dsep y}
  \DeltaRow{Funktionen}{\JoinOp}
}
\begin{tabproofwide}
  \proofstepwidestar[1]{\ZeroJoinSemi{A,\JoinOp,0}}{\rA}
  \proofstepwidestar[2]{x\in A}{\rA}
  \proofstepwidestar[3]{y\in A}{\rA}
  \proofstepwidestar[4]{d\in\LB_{A,\InducedLeq}(\{x,y\})}{\rA}
  \proofstepwidestar[5]{e\in\LB_{A,\InducedLeq}(\{x,y\})}{\rA}
  \proofstepwidestar[1]{\Semi{A,\JoinOp}}{%
    \FormulaRefAuto{ZeroJoinBaseSemilatticeAxiom}[ax]{1}}
  \proofstepwidestar[4]{%
    d\in A\land d\InducedLeq x\land
    d\InducedLeq y}{%
    \FormulaRefAuto{P\leftrightarrow Q, P\vdash Q}{%
      \FormulaRefAuto{PairLowerBoundSetCharacterization}[thm],4}}
  \proofstepwidestar[5]{%
    e\in A\land e\InducedLeq x\land
    e\InducedLeq y}{%
    \FormulaRefAuto{P\leftrightarrow Q, P\vdash Q}{%
      \FormulaRefAuto{PairLowerBoundSetCharacterization}[thm],5}}
  \proofstepwidestar[4]{d\in A}{\rAEa{\rAEa{7}}}
  \proofstepwidestar[5]{e\in A}{\rAEa{\rAEa{8}}}
  \proofstepwidestar[4]{d\InducedLeq x}{%
    \FormulaRefAuto{P\land Q\land R\vdash Q}{7}}
  \proofstepwidestar[5]{e\InducedLeq x}{%
    \FormulaRefAuto{P\land Q\land R\vdash Q}{8}}
  \proofstepwidestar[4]{d\InducedLeq y}{%
    \FormulaRefAuto{P\land Q\land R\vdash R}{7}}
  \proofstepwidestar[5]{e\InducedLeq y}{%
    \FormulaRefAuto{P\land Q\land R\vdash R}{8}}
  \proofstepwidestar[1,4,5]{d\JoinOp{}e\in A}{%
    \FormulaRefAuto{SemilatticeClosureAxiom}[ax]{6,9,10}}
  \proofstepwidestar[1,2,4,5]{%
    d\JoinOp{}e\InducedLeq x}{%
    \FormulaRefAuto{JoinLeastUpper}[thm]{6,9,10,2,11,12}}
  \proofstepwidestar[1,3,4,5]{%
    d\JoinOp{}e\InducedLeq y}{%
    \FormulaRefAuto{JoinLeastUpper}[thm]{6,9,10,3,13,14}}
  \proofstepwidestar[1,2,3,4,5]{%
    d\JoinOp{}e\in\LB_{A,\InducedLeq}(\{x,y\})}{%
    \FormulaRefAuto{P\leftrightarrow Q, Q\vdash P}{%
      \FormulaRefAuto{PairLowerBoundSetCharacterization}[thm],
      \rAI{\rAI{15,16},17}}}
\end{tabproofwide}

\FormulaThmDeltaK[Ein maximales Element einer abgeschlossenen Teilmenge ist größtes Element]{%
  \begin{aligned}[t]
    &\Semi{A,\JoinOp}\dsep D\subseteq A\dsep m\in D\dsep
      \forall q\in D\,\bigl(m\InducedLeq q\to q=m\bigr)\dsep{}\\
    &\forall d,e\in D\,(d\JoinOp{}e\in D)
      \vdash\forall d\in D\,d\InducedLeq m
  \end{aligned}
}{JoinClosedMaximalIsGreatest}{%
  \DeltaRow{Mengen}{A\dsep D\dsep d\dsep e\dsep m\dsep q}
  \DeltaRow{Funktionen}{\JoinOp}
}
\begin{tabproofwide}
  \proofstepwidestar[1]{\Semi{A,\JoinOp}}{\rA}
  \proofstepwidestar[2]{D\subseteq A}{\rA}
  \proofstepwidestar[3]{m\in D}{\rA}
  \proofstepwidestar[4]{%
    \forall q\in D\,\bigl(m\InducedLeq q\to q=m\bigr)}{\rA}
  \proofstepwidestar[5]{\forall d,e\in D\,(d\JoinOp{}e\in D)}{\rA}
  \proofstepwidestar[6]{d\in D}{\rA}
  \proofstepwidestar[2,3]{m\in A}{%
    \FormulaRefAuto{A\subseteq B,\,x\in A\vdash x\in B}{2,3}}
  \proofstepwidestar[2,6]{d\in A}{%
    \FormulaRefAuto{A\subseteq B,\,x\in A\vdash x\in B}{2,6}}
  \proofstepwidestar[3,5,6]{m\JoinOp{}d\in D}{%
    \rRE{6,\rRE{3,\rUE{\rUE{5}}}}}
  \proofstepwidestar[1,2,3,6]{%
    m\InducedLeq m\JoinOp{}d}{%
    \FormulaRefAuto{JoinUpperLeft}[thm]{1,7,8}}
  \proofstepwidestar[3,4,5,6]{m\JoinOp{}d=m}{%
    \rRE{\rRE{9,\rUE{4}},10}}
  \proofstepwidestar[1,2,3,6]{%
    d\InducedLeq m\JoinOp{}d}{%
    \FormulaRefAuto{JoinUpperRight}[thm]{1,7,8}}
  \proofstepwidestar[1,2,3,4,5,6]{%
    d\InducedLeq m}{\rIE{11,12}}
  \proofstepwidestar[1,2,3,4,5]{%
    \forall d\in D\,d\InducedLeq m}{%
    \rUI{\rRI{6,13}}}
\end{tabproofwide}

\FormulaThmDeltaK[Größtes Element eines endlichen Supremums-Halbverbandes mit Null]{%
  \begin{aligned}[t]
    &\Finite A\dsep\ZeroJoinSemi{A,\JoinOp,0}\vdash{}\\
    &\exists e\in A\,\forall x\in A\,(x\InducedLeq e)
  \end{aligned}
}{FiniteZeroJoinSemilatticeHasGreatest}{%
  \DeltaRow{Mengen}{A\dsep d\dsep e\dsep m\dsep q\dsep x\dsep 0}
  \DeltaRow{Funktionen}{\JoinOp}
}
\begin{tabproofwide}
  \proofstepwidestar[1]{\Finite A}{\rA}
  \proofstepwidestar[2]{\ZeroJoinSemi{A,\JoinOp,0}}{\rA}
  \proofstepwidestar[2]{\Semi{A,\JoinOp}}{%
    \FormulaRefAuto{ZeroJoinBaseSemilatticeAxiom}[ax]{2}}
  \proofstepwidestar[2]{\PartOrd{A,\InducedLeq}}{%
    \FormulaRefAuto{JoinSemilatticeInducesOrder}[thm]{3}}
  \proofstepwidestar[2]{0\in A}{%
    \FormulaRefAuto{ZeroJoinCarrierAxiom}[ax]{2}}
  \proofstepwidestar[2]{A\neq\varnothing}{%
    \FormulaRefAuto{a\in S\vdash S\neq\varnothing}{5}}
  \proofstepwidestar[]{A\subseteq A}{\FormulaRefAuto{A\subseteq A}}
  \proofstepwidestar[1,2]{%
    \exists m\in A\,\forall q\in A\,
      (m\InducedLeq q\to q=m)}{%
    \FormulaRefAuto{FinitePosetMaximalElement}[thm]{4,1,6,7}}
  \proofstepwidestar[9]{%
    m\in A\land
    \forall q\in A\,(m\InducedLeq q\to q=m)}{\rA}
  \proofstepwidestar[10]{d\in A}{\rA}
  \proofstepwidestar[11]{e\in A}{\rA}
  \proofstepwidestar[2,10,11]{d\JoinOp{}e\in A}{%
    \FormulaRefAuto{SemilatticeClosureAxiom}[ax]{3,10,11}}
  \proofstepwidestar[2,9]{%
    \forall x\in A\,(x\InducedLeq m)}{%
    \FormulaRefAuto{JoinClosedMaximalIsGreatest}[thm]{%
      3,7,\rAEa{9},\rAEb{9},12}}
  \proofstepwidestar[2,9]{%
    \exists e\in A\,\forall x\in A\,(x\InducedLeq e)}{%
    \rEI{\rAI{\rAEa{9},13}}}
  \proofstepwidestar[1,2]{%
    \exists e\in A\,\forall x\in A\,(x\InducedLeq e)}{%
    \rEE{8,9,14}}
\end{tabproofwide}

\FormulaThmDeltaK[Jeder endliche Supremums-Halbverband mit Null ist beschränkt]{%
  \Finite A\dsep\ZeroJoinSemi{A,\JoinOp,0}
  \vdash\exists e\,\BoundedJoinSemi{A,\JoinOp,0,e}%
}{FiniteZeroJoinSemilatticeIsBounded}{%
  \DeltaRow{Mengen}{A\dsep e\dsep x\dsep 0}
  \DeltaRow{Funktionen}{\JoinOp}
}
\begin{tabproofwide}
  \proofstepwidestar[1]{\Finite A}{\rA}
  \proofstepwidestar[2]{\ZeroJoinSemi{A,\JoinOp,0}}{\rA}
  \proofstepwidestar[1,2]{%
    \exists e\in A\,\forall x\in A\,(x\InducedLeq e)}{%
    \FormulaRefAuto{FiniteZeroJoinSemilatticeHasGreatest}[thm]{1,2}}
  \proofstepwidestar[4]{%
    e\in A\land\forall x\in A\,(x\InducedLeq e)}{\rA}
  \proofstepwidestar[4]{e\in A}{\rAEa{4}}
  \proofstepwidestar[4]{%
    \forall x\in A\,(x\InducedLeq e)}{\rAEb{4}}
  \proofstepwidestar[7]{x\in A}{\rA}
  \proofstepwidestar[2,4,7]{x\JoinOp{}e=e}{%
    \FormulaRefAuto{ZeroJoinGreatestAbsorbs}[thm]{2,5,6,7}}
  \proofstepwidestar[2,4]{%
    \BoundedJoinSemi{A,\JoinOp,0,e}}{%
    \FormulaRefAuto{BoundedJoinSemilattice}[def]{2,5,8}}
  \proofstepwidestar[2,4]{%
    \exists e\,\BoundedJoinSemi{A,\JoinOp,0,e}}{\rEI{9}}
  \proofstepwidestar[1,2]{%
    \exists e\,\BoundedJoinSemi{A,\JoinOp,0,e}}{\rEE{3,4,10}}
\end{tabproofwide}

Damit fällt für endliche Supremums-Halbverbände mit Null die Unterscheidung
zwischen \emph{nach unten beschränkt} und \emph{beschränkt} weg. Außerhalb
des endlichen Falls bleibt sie wesentlich.

\FormulaThmDeltaK[Paarinfima in endlichen Halbverbänden mit Nullelement]{%
  \begin{aligned}[t]
    &\Finite A\dsep\ZeroJoinSemi{A,\JoinOp,0}\dsep x,y\in A\vdash{}\\
    &\exists m\in A\,\Bigl(
      m\in\LB_{A,\InducedLeq}(\{x,y\})\ \land{}\\[-2pt]
    &\hspace{23mm}\bigl(\forall d\in\LB_{A,\InducedLeq}(\{x,y\})\,
        d\InducedLeq m\bigr)\ \land
      \PairInf{A}{\InducedLeq}{x}{y}{m}\Bigr)
  \end{aligned}
}{FiniteZeroJoinSemilatticePairInfima}{%
  \DeltaRow{Mengen}{A\dsep d\dsep e\dsep m\dsep x\dsep y\dsep 0}
  \DeltaRow{Funktionen}{\JoinOp}
}
\begin{tabproofwide}
  \proofstepwidestar[1]{\Finite A}{\rA}
  \proofstepwidestar[2]{\ZeroJoinSemi{A,\JoinOp,0}}{\rA}
  \proofstepwidestar[3]{x\in A}{\rA}
  \proofstepwidestar[4]{y\in A}{\rA}
  \proofstepwidestar[2]{\Semi{A,\JoinOp}}{%
    \FormulaRefAuto{ZeroJoinBaseSemilatticeAxiom}[ax]{2}}
  \proofstepwidestar[2]{\PartOrd{A,\InducedLeq}}{%
    \FormulaRefAuto{JoinSemilatticeInducesOrder}[thm]{5}}
  \proofstepwidestar[3,4]{\{x,y\}\subseteq A}{%
    \FormulaRefAuto{a \in A,\, b \in A \vdash \{a,b\} \subseteq A}{3,4}}
  \proofstepwidestar[2,3,4]{%
    \LB_{A,\InducedLeq}(\{x,y\})\subseteq A}{%
    \FormulaRefAuto{LowerBoundSetSubsetCarrier}[thm]{6,7}}
  \proofstepwidestar[1]{%
    \Finite{\LB_{A,\InducedLeq}(\{x,y\})}}{%
    \FormulaRefAuto{PairLowerBoundSetFinite}[thm]{1}}
  \proofstepwidestar[2,3,4]{%
    0\in\LB_{A,\InducedLeq}(\{x,y\})}{%
    \FormulaRefAuto{PairLowerBoundSetContainsZero}[thm]{2,3,4}}
  \proofstepwidestar[2,3,4]{%
    \LB_{A,\InducedLeq}(\{x,y\})\neq\varnothing}{%
    \FormulaRefAuto{a\in S\vdash S\neq\varnothing}{10}}
  \proofstepwidestar[1,2,3,4]{%
    \begin{aligned}[t]
      &\exists m\in\LB_{A,\InducedLeq}(\{x,y\})\,{}\\[-2pt]
      &\quad\forall q\in\LB_{A,\InducedLeq}(\{x,y\})\,
        (m\InducedLeq q\to q=m)
    \end{aligned}}{%
    \FormulaRefAuto{FinitePosetMaximalElement}[thm]{6,9,11,8}}
  \proofstepwidestar[13]{%
    \begin{aligned}[t]
      &m\in\LB_{A,\InducedLeq}(\{x,y\})\land{}\\[-2pt]
      &\forall q\in\LB_{A,\InducedLeq}(\{x,y\})\,
        (m\InducedLeq q\to q=m)
    \end{aligned}}{\rA}
  \proofstepwidestar[13]{%
    m\in\LB_{A,\InducedLeq}(\{x,y\})}{\rAEa{13}}
  \proofstepwidestar[13]{%
    \forall q\in\LB_{A,\InducedLeq}(\{x,y\})\,
      (m\InducedLeq q\to q=m)}{\rAEb{13}}
  \proofstepwidestar[16]{%
    d\in\LB_{A,\InducedLeq}(\{x,y\})}{\rA}
  \proofstepwidestar[17]{%
    e\in\LB_{A,\InducedLeq}(\{x,y\})}{\rA}
  \proofstepwidestar[2,3,4,16,17]{%
    d\JoinOp{}e\in\LB_{A,\InducedLeq}(\{x,y\})}{%
    \FormulaRefAuto{PairLowerBoundSetJoinClosed}[thm]{%
      2,3,4,16,17}}
  \proofstepwidestar[2,3,4,13]{%
    \forall d\in\LB_{A,\InducedLeq}(\{x,y\})\,
      (d\InducedLeq m)}{%
    \FormulaRefAuto{JoinClosedMaximalIsGreatest}[thm]{%
      5,8,14,15,18}}
  \proofstepwidestar[2,3,4,13]{m\in A}{%
    \FormulaRefAuto{A\subseteq B,\,x\in A\vdash x\in B}{8,14}}
  \proofstepwidestar[2,3,4,13]{%
    \PairInf{A}{\InducedLeq}{x}{y}{m}}{%
    \FormulaRefAuto{InfimumCandidateEqualsInfimum}[thm]{%
      6,7,14,19}}
  \proofstepwidestar[2,3,4,13]{%
    \begin{aligned}[t]
      &m\in\LB_{A,\InducedLeq}(\{x,y\})\land{}\\[-2pt]
      &\bigl(\forall d\in\LB_{A,\InducedLeq}(\{x,y\})\,
          d\InducedLeq m\bigr)\land
        \PairInf{A}{\InducedLeq}{x}{y}{m}
    \end{aligned}}{%
    \rAI{\rAI{14,19},21}}
  \proofstepwidestar[2,3,4,13]{%
    \begin{aligned}[t]
      &\exists m\in A\,\Bigl(
        m\in\LB_{A,\InducedLeq}(\{x,y\})\land{}\\[-2pt]
      &\quad\bigl(\forall d\in\LB_{A,\InducedLeq}(\{x,y\})\,
          d\InducedLeq m\bigr)\land
        \PairInf{A}{\InducedLeq}{x}{y}{m}\Bigr)
    \end{aligned}}{\rEI{\rAI{20,22}}}
  \proofstepwidestar[1,2,3,4]{%
    \begin{aligned}[t]
      &\exists m\in A\,\Bigl(
        m\in\LB_{A,\InducedLeq}(\{x,y\})\land{}\\[-2pt]
      &\quad\bigl(\forall d\in\LB_{A,\InducedLeq}(\{x,y\})\,
          d\InducedLeq m\bigr)\land
        \PairInf{A}{\InducedLeq}{x}{y}{m}\Bigr)
    \end{aligned}}{\rEE{12,13,23}}
\end{tabproofwide}

\FormulaDefDeltaK[Trennende untere Elemente]{%
  \operatorname{Sep}_{A,\JoinOp}(x,y)
  \coloneqq
  \bigl\{d\in A\mid
    d\InducedLeq x\land
    d\not\InducedLeq y
  \bigr\}
}{JoinSeparatorSet}{%
  \DeltaRow{Mengen}{A\dsep d\dsep x\dsep y}
  \DeltaRow{Funktionen}{\JoinOp}
  \DeltaRow{Relationsoperatoren}{\InducedLeq}
  \DeltaRow{\textbf{Neues Symbol}}{}
  \DeltaPrem{Trennende untere Elemente}{%
    \operatorname{Sep}_{A,\JoinOp}(x,y)}
}

\FormulaThmDeltaK[Charakterisierung der trennenden unteren Elemente]{%
  \begin{aligned}[t]
    d\in\operatorname{Sep}_{A,\JoinOp}(x,y)
    \quad\eqvdash\quad{}&d\in A\land
      d\InducedLeq x\land{}\\
    &d\not\InducedLeq y
  \end{aligned}
}{JoinSeparatorSetCharacterization}{%
  \DeltaRow{Mengen}{A\dsep d\dsep x\dsep y}
  \DeltaRow{Funktionen}{\JoinOp}
}
\begin{tabproofwide}
  \proofstepwide{d\in\operatorname{Sep}_{A,\JoinOp}(x,y)}{\leftrightarrow}{%
    d\in A\land d\InducedLeq x\land
    d\not\InducedLeq y}{%
    \FormulaRefAuto{JoinSeparatorSet}[def]}
\end{tabproofwide}

\FormulaThmDeltaK[Trennende untere Elemente liegen im Träger]{%
  \operatorname{Sep}_{A,\JoinOp}(x,y)\subseteq A
}{JoinSeparatorSetSubset}{%
  \DeltaRow{Mengen}{A\dsep d\dsep x\dsep y}
  \DeltaRow{Funktionen}{\JoinOp}
}
\begin{tabproof}
  \proofstep{1}{d\in\operatorname{Sep}_{A,\JoinOp}(x,y)}{\rA}
  \proofstep{1}{%
    d\in A\land d\InducedLeq x\land
    d\not\InducedLeq y}{%
    \FormulaRefAuto{P\leftrightarrow Q, P\vdash Q}{%
      \FormulaRefAuto{JoinSeparatorSetCharacterization}[thm],1}}
  \proofstep{1}{d\in A}{\rAEa{\rAEa{2}}}
  \proofstep{}{\operatorname{Sep}_{A,\JoinOp}(x,y)\subseteq A}{%
    \FormulaRefAuto{A \subseteq B \coloneq \forall x\,(x\in A \rightarrow x\in B)}[def]{%
      \rUI{\rRI{1,3}}}}
\end{tabproof}

\FormulaThmDeltaK[Das linke Element trennt sich selbst vom rechten]{%
  \begin{aligned}[t]
    &\ZeroJoinSemi{A,\JoinOp,0}\dsep x,y\in A\dsep
      x\not\InducedLeq y\vdash{}\\
    &x\in\operatorname{Sep}_{A,\JoinOp}(x,y)
  \end{aligned}
}{JoinSeparatorSetContainsLeft}{%
  \DeltaRow{Mengen}{A\dsep 0\dsep x\dsep y}
  \DeltaRow{Funktionen}{\JoinOp}
}
\begin{tabproofwide}
  \proofstepwidestar[1]{\ZeroJoinSemi{A,\JoinOp,0}}{\rA}
  \proofstepwidestar[2]{x\in A}{\rA}
  \proofstepwidestar[3]{y\in A}{\rA}
  \proofstepwidestar[4]{x\not\InducedLeq y}{\rA}
  \proofstepwidestar[1]{\Semi{A,\JoinOp}}{%
    \FormulaRefAuto{ZeroJoinBaseSemilatticeAxiom}[ax]{1}}
  \proofstepwidestar[1,2]{x\InducedLeq x}{%
    \FormulaRefAuto{JoinOrderReflexive}[thm]{5,2}}
  \proofstepwidestar[1,2,3,4]{%
    x\in A\land x\InducedLeq x\land
    x\not\InducedLeq y}{\rAI{\rAI{2,6},4}}
  \proofstepwidestar[1,2,3,4]{%
    x\in\operatorname{Sep}_{A,\JoinOp}(x,y)}{%
    \FormulaRefAuto{P\leftrightarrow Q, Q\vdash P}{%
      \FormulaRefAuto{JoinSeparatorSetCharacterization}[thm],7}}
\end{tabproofwide}

\FormulaThmDeltaK[Endlichkeit der trennenden unteren Elemente]{%
  \Finite A\vdash\Finite{\operatorname{Sep}_{A,\JoinOp}(x,y)}
}{JoinSeparatorSetFinite}{%
  \DeltaRow{Mengen}{A\dsep x\dsep y}
  \DeltaRow{Funktionen}{\JoinOp}
}
\begin{tabproof}
  \proofstep{1}{\Finite A}{\rA}
  \proofstep{}{\operatorname{Sep}_{A,\JoinOp}(x,y)\subseteq A}{%
    \FormulaRefAuto{JoinSeparatorSetSubset}[thm]}
  \proofstep{1}{\Finite{\operatorname{Sep}_{A,\JoinOp}(x,y)}}{%
    \FormulaRefAuto{FiniteSubset}[thm]{2,1}}
\end{tabproof}

\FormulaThmDeltaK[Echte untere Trenner liegen unter dem rechten Element]{%
  \begin{aligned}[t]
    &\Semi{A,\JoinOp}\dsep x,y,u\in A\dsep
      j\in\operatorname{Sep}_{A,\JoinOp}(x,y)\dsep{}\\
    &\forall d\in\operatorname{Sep}_{A,\JoinOp}(x,y)\,
      \bigl(d\InducedLeq j\to d=j\bigr)\dsep{}\\
    &u\InducedLeq j\dsep u\neq j
      \vdash u\InducedLeq y
  \end{aligned}
}{StrictLowerSeparatorBelowRight}{%
  \DeltaRow{Mengen}{A\dsep d\dsep j\dsep u\dsep x\dsep y}
  \DeltaRow{Funktionen}{\JoinOp}
}
\begin{tabproofwide}
  \proofstepwidestar[1]{\Semi{A,\JoinOp}}{\rA}
  \proofstepwidestar[2]{x\in A}{\rA}
  \proofstepwidestar[3]{y\in A}{\rA}
  \proofstepwidestar[4]{u\in A}{\rA}
  \proofstepwidestar[5]{j\in\operatorname{Sep}_{A,\JoinOp}(x,y)}{\rA}
  \proofstepwidestar[6]{%
    \forall d\in\operatorname{Sep}_{A,\JoinOp}(x,y)\,
      (d\InducedLeq j\to d=j)}{\rA}
  \proofstepwidestar[7]{u\InducedLeq j}{\rA}
  \proofstepwidestar[8]{u\neq j}{\rA}
  \proofstepwidestar[5]{%
    j\in A\land j\InducedLeq x\land j\not\InducedLeq y}{%
    \FormulaRefAuto{P\leftrightarrow Q, P\vdash Q}{%
      \FormulaRefAuto{JoinSeparatorSetCharacterization}[thm],5}}
  \proofstepwidestar[5]{j\InducedLeq x}{%
    \FormulaRefAuto{P\land Q\land R\vdash Q}{9}}
  \proofstepwidestar[11]{u\not\InducedLeq y}{\rA}
  \proofstepwidestar[1,5,7]{u\InducedLeq x}{%
    \FormulaRefAuto{JoinOrderTransitive}[thm]{1,7,10}}
  \proofstepwidestar[1,4,5,7,11]{%
    u\in A\land u\InducedLeq x\land u\not\InducedLeq y}{%
    \rAI{\rAI{4,12},11}}
  \proofstepwidestar[1,4,5,7,11]{%
    u\in\operatorname{Sep}_{A,\JoinOp}(x,y)}{%
    \FormulaRefAuto{P\leftrightarrow Q, Q\vdash P}{%
      \FormulaRefAuto{JoinSeparatorSetCharacterization}[thm],13}}
  \proofstepwidestar[1,4,5,6,7,11]{u=j}{%
    \rRE{7,\rRE{14,\rUE{6}}}}
  \proofstepwidestar[1,4,5,6,7,8,11]{\bot}{\rBI{15,8}}
  \proofstepwidestar[1,2,3,4,5,6,7,8]{%
    u\InducedLeq y}{\rCI{11,16}}
\end{tabproofwide}

\FormulaThmDeltaK[Ein minimaler Trenner ist \(\JoinOp\)-irreduzibel]{%
  \begin{aligned}[t]
    &\ZeroJoinSemi{A,\JoinOp,0}\dsep x,y\in A\dsep
      j\in\operatorname{Sep}_{A,\JoinOp}(x,y)\dsep{}\\
    &\forall d\in\operatorname{Sep}_{A,\JoinOp}(x,y)\,
      \bigl(d\InducedLeq j\to d=j\bigr)
      \vdash\JoinIrr{A,\JoinOp,0}{j}
  \end{aligned}
}{MinimalSeparatorJoinIrreducible}{%
  \DeltaRow{Mengen}{A\dsep 0\dsep d\dsep j\dsep u\dsep v\dsep x\dsep y}
  \DeltaRow{Funktionen}{\JoinOp}
}
\begin{tabproofsplitwide}
  \proofpartwideindR[Ein minimaler Trenner ist nicht Null]{%
    \begin{aligned}[t]
      &\ZeroJoinSemi{A,\JoinOp,0}\dsep x,y\in A\dsep
        j\in\operatorname{Sep}_{A,\JoinOp}(x,y)\vdash{}\\
      &j\neq0
    \end{aligned}
  }{\ZeroJoinSemi{A,\JoinOp,0}\dsep x,y\in A\dsep
    j\in\operatorname{Sep}_{A,\JoinOp}(x,y)\vdash j\neq0}%
    [MinimalSeparatorNonzero]
    \proofstepwidestar[1]{\ZeroJoinSemi{A,\JoinOp,0}}{\rA}
    \proofstepwidestar[2]{x\in A}{\rA}
    \proofstepwidestar[3]{y\in A}{\rA}
    \proofstepwidestar[4]{j\in\operatorname{Sep}_{A,\JoinOp}(x,y)}{\rA}
    \proofstepwidestar[4]{%
      j\in A\land j\InducedLeq x\land j\not\InducedLeq y}{%
      \FormulaRefAuto{P\leftrightarrow Q, P\vdash Q}{%
        \FormulaRefAuto{JoinSeparatorSetCharacterization}[thm],4}}
    \proofstepwidestar[4]{j\not\InducedLeq y}{%
      \FormulaRefAuto{P\land Q\land R\vdash R}{5}}
    \proofstepwidestar[1,3]{0\InducedLeq y}{%
      \FormulaRefAuto{ZeroJoinLeast}[thm]{1,3}}
    \proofstepwidestar[8]{j=0}{\rA}
    \proofstepwidestar[8]{0=j}{\FormulaRefAuto{a=b\vdash b=a}{8}}
    \proofstepwidestar[1,3,8]{j\InducedLeq y}{\rIE{9,7}}
    \proofstepwidestar[1,3,4,8]{\bot}{\rBI{6,10}}
    \proofstepwidestar[1,3,4]{j\neq0}{\rCI{8,11}}
  \closeproofpartwideindR

  \proofpartwideindR[Jede Zerlegung enthält den minimalen Trenner]{%
    \begin{aligned}[t]
      &\ZeroJoinSemi{A,\JoinOp,0}\dsep x,y\in A\dsep
        j\in\operatorname{Sep}_{A,\JoinOp}(x,y)\dsep{}\\[-2pt]
      &\forall d\in\operatorname{Sep}_{A,\JoinOp}(x,y)\,
        (d\InducedLeq j\to d=j)\vdash{}\\[-2pt]
      &u,v\in A\dsep u\JoinOp{}v=j\vdash u=j\lor v=j
    \end{aligned}
  }{\begin{aligned}[t]
      &\ZeroJoinSemi{A,\JoinOp,0}\dsep x,y\in A\dsep
        j\in\operatorname{Sep}_{A,\JoinOp}(x,y)\dsep{}\\[-2pt]
      &\forall d\in\operatorname{Sep}_{A,\JoinOp}(x,y)\,
        (d\InducedLeq j\to d=j)\vdash{}\\[-2pt]
      &u,v\in A\dsep u\JoinOp{}v=j\vdash u=j\lor v=j
    \end{aligned}}[MinimalSeparatorDecomposition]
    \proofstepwidestar[1]{\ZeroJoinSemi{A,\JoinOp,0}}{\rA}
    \proofstepwidestar[2]{x\in A}{\rA}
    \proofstepwidestar[3]{y\in A}{\rA}
    \proofstepwidestar[4]{j\in\operatorname{Sep}_{A,\JoinOp}(x,y)}{\rA}
    \proofstepwidestar[5]{%
      \forall d\in\operatorname{Sep}_{A,\JoinOp}(x,y)\,
        (d\InducedLeq j\to d=j)}{\rA}
    \proofstepwidestar[1]{\Semi{A,\JoinOp}}{%
      \FormulaRefAuto{ZeroJoinBaseSemilatticeAxiom}[ax]{1}}
    \proofstepwidestar[4]{%
      j\in A\land j\InducedLeq x\land j\not\InducedLeq y}{%
      \FormulaRefAuto{P\leftrightarrow Q, P\vdash Q}{%
        \FormulaRefAuto{JoinSeparatorSetCharacterization}[thm],4}}
    \proofstepwidestar[4]{j\not\InducedLeq y}{%
      \FormulaRefAuto{P\land Q\land R\vdash R}{7}}
    \proofstepwidestar[9]{u\in A}{\rA}
    \proofstepwidestar[10]{v\in A}{\rA}
    \proofstepwidestar[11]{u\JoinOp{}v=j}{\rA}
    \proofstepwidestar[6,9,10]{u\InducedLeq u\JoinOp{}v}{%
      \FormulaRefAuto{JoinUpperLeft}[thm]{6,9,10}}
    \proofstepwidestar[6,9,10,11]{u\InducedLeq j}{\rIE{11,12}}
    \proofstepwidestar[6,9,10]{v\InducedLeq u\JoinOp{}v}{%
      \FormulaRefAuto{JoinUpperRight}[thm]{6,9,10}}
    \proofstepwidestar[6,9,10,11]{v\InducedLeq j}{\rIE{11,14}}
    \proofstepwidestar[]{u=j\lor u\neq j}{\FormulaRefAuto{P\lor\neg P}}
    \proofstepwidestar[17]{u=j}{\rA}
    \proofstepwidestar[17]{u=j\lor v=j}{\rOIa{17}}
    \proofstepwidestar[19]{u\neq j}{\rA}
    \proofstepwidestar[1,2,3,4,5,9,13,19]{u\InducedLeq y}{%
      \FormulaRefAuto{StrictLowerSeparatorBelowRight}[thm]{%
        6,2,3,9,4,5,13,19}}
    \proofstepwidestar[]{v=j\lor v\neq j}{\FormulaRefAuto{P\lor\neg P}}
    \proofstepwidestar[22]{v=j}{\rA}
    \proofstepwidestar[22]{u=j\lor v=j}{\rOIb{22}}
    \proofstepwidestar[24]{v\neq j}{\rA}
    \proofstepwidestar[1,2,3,4,5,10,15,24]{v\InducedLeq y}{%
      \FormulaRefAuto{StrictLowerSeparatorBelowRight}[thm]{%
        6,2,3,10,4,5,15,24}}
    \proofstepwidestar[1,2,3,9,10,19,24]{u\JoinOp{}v\InducedLeq y}{%
      \FormulaRefAuto{JoinLeastUpper}[thm]{6,9,10,3,20,25}}
    \proofstepwidestar[1,2,3,4,9,10,11,19,24]{j\InducedLeq y}{%
      \rIE{11,26}}
    \proofstepwidestar[1,2,3,4,9,10,11,19,24]{\bot}{\rBI{8,27}}
    \proofstepwidestar[1,2,3,4,9,10,11,19,24]{u=j\lor v=j}{%
      \FormulaRefAuto{\bot\vdash P}{28}}
    \proofstepwidestar[1,2,3,4,5,9,10,11,19]{u=j\lor v=j}{%
      \rOE{21,22,23,24,29}}
    \proofstepwidestar[1,2,3,4,5,9,10,11]{u=j\lor v=j}{%
      \rOE{16,17,18,19,30}}
  \closeproofpartwideindR

  \proofpartwide{Schluss}
    \proofstepwidestar[1]{\ZeroJoinSemi{A,\JoinOp,0}}{\rA}
    \proofstepwidestar[2]{x\in A}{\rA}
    \proofstepwidestar[3]{y\in A}{\rA}
    \proofstepwidestar[4]{j\in\operatorname{Sep}_{A,\JoinOp}(x,y)}{\rA}
    \proofstepwidestar[5]{%
      \forall d\in\operatorname{Sep}_{A,\JoinOp}(x,y)\,
        (d\InducedLeq j\to d=j)}{\rA}
    \proofstepwidestar[4]{%
      j\in A\land j\InducedLeq x\land j\not\InducedLeq y}{%
      \FormulaRefAuto{P\leftrightarrow Q, P\vdash Q}{%
        \FormulaRefAuto{JoinSeparatorSetCharacterization}[thm],4}}
    \proofstepwidestar[4]{j\in A}{\rAEa{\rAEa{6}}}
    \proofstepwidestar[1,2,3,4]{j\neq0}{%
      \FormulaRefAuto{MinimalSeparatorNonzero}[thm]{1,2,3,4}}
    \proofstepwidestar[9]{u,v\in A}{\rA}
    \proofstepwidestar[10]{u\JoinOp{}v=j}{\rA}
    \proofstepwidestar[1,2,3,4,5,9,10]{u=j\lor v=j}{%
      \FormulaRefAuto{MinimalSeparatorDecomposition}[thm]{%
        1,2,3,4,5,9,10}}
    \proofstepwidestar[1,2,3,4,5]{\JoinIrr{A,\JoinOp,0}{j}}{%
      \FormulaRefAuto{JoinIrreducible}[def]{7,8,11}\qedhere}
  \closeproofpartwide
\end{tabproofsplitwide}

\FormulaThmDeltaK[\(\JoinOp\)-Irreduzible trennen Elemente]{%
  \begin{aligned}[t]
    &\Finite A\dsep\ZeroJoinSemi{A,\JoinOp,0}\dsep x,y\in A\dsep
      x\not\InducedLeq y\vdash{}\\
    &\exists j\in\JoinIrrSet{A,\JoinOp,0}\,
      \bigl(j\InducedLeq x\land
            j\not\InducedLeq y\bigr)
  \end{aligned}
}{JoinIrreduciblesSeparate}{%
  \DeltaRow{Mengen}{A\dsep d\dsep j\dsep x\dsep y\dsep 0}
  \DeltaRow{Funktionen}{\JoinOp}
}
\begin{tabproofwide}
  \proofstepwidestar[1]{\Finite A}{\rA}
  \proofstepwidestar[2]{\ZeroJoinSemi{A,\JoinOp,0}}{\rA}
  \proofstepwidestar[3]{x\in A}{\rA}
  \proofstepwidestar[4]{y\in A}{\rA}
  \proofstepwidestar[5]{x\not\InducedLeq y}{\rA}
  \proofstepwidestar[2]{\Semi{A,\JoinOp}}{%
    \FormulaRefAuto{ZeroJoinBaseSemilatticeAxiom}[ax]{2}}
  \proofstepwidestar[2]{\PartOrd{A,\InducedLeq}}{%
    \FormulaRefAuto{JoinSemilatticeInducesOrder}[thm]{6}}
  \proofstepwidestar[]{%
    \operatorname{Sep}_{A,\JoinOp}(x,y)\subseteq A}{%
    \FormulaRefAuto{JoinSeparatorSetSubset}[thm]}
  \proofstepwidestar[1]{%
    \Finite{\operatorname{Sep}_{A,\JoinOp}(x,y)}}{%
    \FormulaRefAuto{JoinSeparatorSetFinite}[thm]{1}}
  \proofstepwidestar[2,3,4,5]{%
    x\in\operatorname{Sep}_{A,\JoinOp}(x,y)}{%
    \FormulaRefAuto{JoinSeparatorSetContainsLeft}[thm]{2,3,4,5}}
  \proofstepwidestar[2,3,4,5]{%
    \operatorname{Sep}_{A,\JoinOp}(x,y)\neq\varnothing}{%
    \FormulaRefAuto{a\in S\vdash S\neq\varnothing}{10}}
  \proofstepwidestar[1,2,3,4,5]{%
    \begin{aligned}[t]
      &\exists j\in\operatorname{Sep}_{A,\JoinOp}(x,y)\,{}\\[-2pt]
      &\quad\forall d\in\operatorname{Sep}_{A,\JoinOp}(x,y)\,
        (d\InducedLeq j\to d=j)
    \end{aligned}}{%
    \ProofRefStack{%
      \FormulaRefAuto{FinitePosetMinimalElement}[thm]\\
      (7,9,11,8)}}
  \proofstepwidestar[13]{%
    \begin{aligned}[t]
      &j\in\operatorname{Sep}_{A,\JoinOp}(x,y)\land{}\\[-2pt]
      &\forall d\in\operatorname{Sep}_{A,\JoinOp}(x,y)\,
        (d\InducedLeq j\to d=j)
    \end{aligned}}{\rA}
  \proofstepwidestar[13]{%
    j\in\operatorname{Sep}_{A,\JoinOp}(x,y)}{\rAEa{13}}
  \proofstepwidestar[13]{%
    \forall d\in\operatorname{Sep}_{A,\JoinOp}(x,y)\,
      (d\InducedLeq j\to d=j)}{\rAEb{13}}
  \proofstepwidestar[2,3,4,13]{\JoinIrr{A,\JoinOp,0}{j}}{%
    \FormulaRefAuto{MinimalSeparatorJoinIrreducible}[thm]{%
      2,3,4,14,15}}
  \proofstepwidestar[13]{j\in A}{%
    \FormulaRefAuto{A\subseteq B,\,x\in A\vdash x\in B}{8,14}}
  \proofstepwidestar[2,3,4,13]{j\in\JoinIrrSet{A,\JoinOp,0}}{%
    \FormulaRefAuto{JoinIrreducibleSet}[def]{17,16}}
  \proofstepwidestar[13]{%
    j\in A\land j\InducedLeq x\land j\not\InducedLeq y}{%
    \FormulaRefAuto{P\leftrightarrow Q, P\vdash Q}{%
      \FormulaRefAuto{JoinSeparatorSetCharacterization}[thm],14}}
  \proofstepwidestar[13]{j\InducedLeq x}{%
    \FormulaRefAuto{P\land Q\land R\vdash Q}{19}}
  \proofstepwidestar[13]{j\not\InducedLeq y}{%
    \FormulaRefAuto{P\land Q\land R\vdash R}{19}}
  \proofstepwidestar[2,3,4,13]{%
    \exists j\in\JoinIrrSet{A,\JoinOp,0}\,
      (j\InducedLeq x\land j\not\InducedLeq y)}{%
    \rEI{\rAI{18,\rAI{20,21}}}}
  \proofstepwidestar[1,2,3,4,5]{%
    \exists j\in\JoinIrrSet{A,\JoinOp,0}\,
      (j\InducedLeq x\land j\not\InducedLeq y)}{%
    \rEE{12,13,22}}
\end{tabproofwide}


\FormulaThmDeltaK[Ein nichttrivialer endlicher Halbverband besitzt Irreduzible]{%
  \begin{aligned}[t]
    &\Finite A\dsep\ZeroJoinSemi{A,\JoinOp,0}\dsep
      \exists x\in A\,(x\neq0)\vdash{}\\
    &\JoinIrrSet{A,\JoinOp,0}\neq\varnothing
  \end{aligned}
}{NontrivialFiniteZeroJoinHasIrreducible}{%
  \DeltaRow{Mengen}{A\dsep j\dsep x\dsep 0}
  \DeltaRow{Funktionen}{\JoinOp}
}
\begin{tabproofwide}
  \proofstepwidestar[1]{\Finite A}{\rA}
  \proofstepwidestar[2]{\ZeroJoinSemi{A,\JoinOp,0}}{\rA}
  \proofstepwidestar[3]{\exists x\in A\,(x\neq0)}{\rA}
  \proofstepwidestar[4]{x\in A\land x\neq0}{\rA}
  \proofstepwidestar[4]{x\in A}{\rAEa{4}}
  \proofstepwidestar[4]{x\neq0}{\rAEb{4}}
  \proofstepwidestar[2]{\Semi{A,\JoinOp}}{%
    \FormulaRefAuto{ZeroJoinBaseSemilatticeAxiom}[ax]{2}}
  \proofstepwidestar[8]{x\InducedLeq 0}{\rA}
  \proofstepwidestar[2,4]{0\InducedLeq x}{%
    \FormulaRefAuto{ZeroJoinLeast}[thm]{2,5}}
  \proofstepwidestar[2,4,8]{x=0}{%
    \FormulaRefAuto{JoinOrderAntisymmetric}[thm]{7,8,9}}
  \proofstepwidestar[2,4,8]{\bot}{\rBI{10,6}}
  \proofstepwidestar[2,4]{x\not\InducedLeq 0}{\rCI{8,11}}
  \proofstepwidestar[1,2,4]{%
    \exists j\in\JoinIrrSet{A,\JoinOp,0}\,\bigl(
      j\InducedLeq x\land
      j\not\InducedLeq 0\bigr)}{%
    \FormulaRefAuto{JoinIrreduciblesSeparate}[thm]{%
      1,2,5,\FormulaRefAuto{ZeroJoinCarrierAxiom}[ax]{2},12}}
  \proofstepwidestar[14]{%
    j\in\JoinIrrSet{A,\JoinOp,0}\land\bigl(
      j\InducedLeq x\land
      j\not\InducedLeq 0\bigr)}{\rA}
  \proofstepwidestar[14]{j\in\JoinIrrSet{A,\JoinOp,0}}{\rAEa{14}}
  \proofstepwidestar[14]{\exists j\,j\in\JoinIrrSet{A,\JoinOp,0}}{\rEI{15}}
  \proofstepwidestar[1,2,4]{%
    \exists j\,j\in\JoinIrrSet{A,\JoinOp,0}}{\rEE{13,14,16}}
  \proofstepwidestar[1,2,4]{\JoinIrrSet{A,\JoinOp,0}\neq\varnothing}{%
    \FormulaRefAuto{A\neq\varnothing \eqvdash \exists x\,(x \in A)}{17}}
  \proofstepwidestar[1,2,3]{\JoinIrrSet{A,\JoinOp,0}\neq\varnothing}{%
    \rEE{3,4,18}}
\end{tabproofwide}


\section{Echte untere Elemente von Irreduziblen}

\FormulaDefDeltaK[Echte untere Elemente eines Irreduziblen]{%
  D_L(P)\coloneqq
  \{Q\in L\mid Q\InducedLeq P\land Q\neq P\}
}{JoinIrreducibleProperLowerSet}{%
  \DeltaRow{Mengen}{L\dsep P\dsep Q}
  \DeltaRow{Funktionen}{\JoinOp}
  \DeltaRow{\textbf{Neues Symbol}}{}
  \DeltaPrem{Echte untere Elemente}{D_L(P)}
}
\leavevmode\par\medskip

\FormulaThmDeltaK[Echte untere Elemente liegen im Träger]{%
  Q\in D_L(P)\vdash Q\in L
}{JoinIrreducibleProperLowerCarrier}{%
  \DeltaRow{Mengen}{L\dsep P\dsep Q}
  \DeltaRow{Funktionen}{\JoinOp}
}
\begin{tabproof}
  \proofstep{1}{Q\in D_L(P)}{\rA}
  \proofstep{1}{Q\in L\land Q\InducedLeq P\land Q\neq P}{%
    \FormulaRefAuto{JoinIrreducibleProperLowerSet}[def]{1}}
  \proofstep{1}{Q\in L}{\rAEa{\rAEa{2}}}
\end{tabproof}

\FormulaThmDeltaK[Die echten unteren Elemente bilden eine endliche nichtleere Teilmenge]{%
  \begin{aligned}[t]
    &\Finite L\dsep\ZeroJoinSemi{L,\JoinOp,\varnothing}\dsep
      \JoinIrr{L,\JoinOp,\varnothing}{P}\vdash{}\\
    &\Finite{D_L(P)}\land D_L(P)\neq\varnothing
      \land D_L(P)\subseteq L
  \end{aligned}
}{JoinIrreducibleProperLowerFinite}{%
  \DeltaRow{Mengen}{L\dsep P\dsep Q}
  \DeltaRow{Funktionen}{\JoinOp}
}
\begin{tabproofwide}
  \proofstepwidestar[1]{\Finite L}{\rA}
  \proofstepwidestar[2]{\ZeroJoinSemi{L,\JoinOp,\varnothing}}{\rA}
  \proofstepwidestar[3]{\JoinIrr{L,\JoinOp,\varnothing}{P}}{\rA}
  \proofstepwidestar[3]{P\in L}{\FormulaRefAuto{JoinIrreducible}[def]{3}}
  \proofstepwidestar[3]{P\neq\varnothing}{\FormulaRefAuto{JoinIrreducible}[def]{3}}
  \proofstepwidestar[2]{\varnothing\in L}{%
    \FormulaRefAuto{ZeroJoinCarrierAxiom}[ax]{2}}
  \proofstepwidestar[2,4]{\varnothing\InducedLeq P}{%
    \FormulaRefAuto{ZeroJoinLeast}[thm]{2,4}}
  \proofstepwidestar[2,3]{\varnothing\in D_L(P)}{%
    \FormulaRefAuto{JoinIrreducibleProperLowerSet}[def]{6,7,5}}
  \proofstepwidestar[2,3]{D_L(P)\neq\varnothing}{%
    \FormulaRefAuto{a\in S\vdash S\neq\varnothing}{8}}
  \proofstepwidestar[]{D_L(P)\subseteq L}{%
    \FormulaRefAuto{JoinIrreducibleProperLowerSet}[def]}
  \proofstepwidestar[1]{\Finite{D_L(P)}}{%
    \FormulaRefAuto{FiniteSubset}[thm]{10,1}}
  \proofstepwidestar[1,2,3]{\Finite{D_L(P)}\land D_L(P)\neq\varnothing
    \land D_L(P)\subseteq L}{\rAI{11,\rAI{9,10}}}
\end{tabproofwide}

\FormulaThmDeltaK[Die echten unteren Elemente sind unter dem Supremum abgeschlossen]{%
  \ZeroJoinSemi{L,\JoinOp,\varnothing}\dsep
  \JoinIrr{L,\JoinOp,\varnothing}{P}\dsep Q,S\in D_L(P)
  \vdash Q\JoinOp{}S\in D_L(P)
}{JoinIrreducibleProperLowerJoinClosed}{%
  \DeltaRow{Mengen}{L\dsep P\dsep Q\dsep S}
  \DeltaRow{Funktionen}{\JoinOp}
}
\begin{tabproofsplitwide}
  \proofpartwideindR[Das Supremum liegt im Träger und unter \(P\)]{%
    \begin{aligned}[t]
      &\ZeroJoinSemi{L,\JoinOp,\varnothing}\dsep
        \JoinIrr{L,\JoinOp,\varnothing}{P}\dsep
        Q,S\in D_L(P)\vdash{}\\
      &Q\JoinOp{}S\in L\land
        Q\JoinOp{}S\InducedLeq P
    \end{aligned}
  }{%
    \ZeroJoinSemi{L,\JoinOp,\varnothing}\dsep
    \JoinIrr{L,\JoinOp,\varnothing}{P}\dsep Q,S\in D_L(P)\vdash
    Q\JoinOp{}S\in L\land
    Q\JoinOp{}S\InducedLeq P
  }[JoinIrreducibleProperLowerJoinBound]
    \proofstepwidestar[1]{\ZeroJoinSemi{L,\JoinOp,\varnothing}}{\rA}
    \proofstepwidestar[2]{\JoinIrr{L,\JoinOp,\varnothing}{P}}{\rA}
    \proofstepwidestar[3]{Q,S\in D_L(P)}{\rA}
    \proofstepwidestar[3]{Q\in D_L(P)}{\rAEa{3}}
    \proofstepwidestar[3]{S\in D_L(P)}{\rAEb{3}}
    \proofstepwidestar[3]{Q\in L}{%
      \FormulaRefAuto{JoinIrreducibleProperLowerCarrier}[thm]{4}}
    \proofstepwidestar[3]{S\in L}{%
      \FormulaRefAuto{JoinIrreducibleProperLowerCarrier}[thm]{5}}
    \proofstepwidestar[2]{P\in L}{%
      \FormulaRefAuto{JoinIrreducible}[def]{2}}
    \proofstepwidestar[3]{%
      Q\in L\land Q\InducedLeq P\land Q\neq P}{%
      \FormulaRefAuto{JoinIrreducibleProperLowerSet}[def]{4}}
    \proofstepwidestar[3]{Q\InducedLeq P}{\rAEa{\rAEb{9}}}
    \proofstepwidestar[3]{%
      S\in L\land S\InducedLeq P\land S\neq P}{%
      \FormulaRefAuto{JoinIrreducibleProperLowerSet}[def]{5}}
    \proofstepwidestar[3]{S\InducedLeq P}{\rAEa{\rAEb{11}}}
    \proofstepwidestar[1]{\Semi{L,\JoinOp}}{%
      \FormulaRefAuto{ZeroJoinBaseSemilatticeAxiom}[ax]{1}}
    \proofstepwidestar[1,3]{Q\JoinOp{}S\in L}{%
      \FormulaRefAuto{SemilatticeClosureAxiom}[ax]{13,6,7}}
    \proofstepwidestar[1,2,3]{Q\JoinOp{}S\InducedLeq P}{%
      \FormulaRefAuto{JoinLeastUpper}[thm]{13,6,7,8,10,12}}
    \proofstepwidestar[1,2,3]{%
      Q\JoinOp{}S\in L\land
      Q\JoinOp{}S\InducedLeq P}{\rAI{14,15}}
  \closeproofpartwideindR

  \proofpartwideindR[Das Supremum bleibt von \(P\) verschieden]{%
    \ZeroJoinSemi{L,\JoinOp,\varnothing}\dsep
    \JoinIrr{L,\JoinOp,\varnothing}{P}\dsep Q,S\in D_L(P)
    \vdash Q\JoinOp{}S\neq P
  }{%
    \ZeroJoinSemi{L,\JoinOp,\varnothing}\dsep
    \JoinIrr{L,\JoinOp,\varnothing}{P}\dsep Q,S\in D_L(P)
    \vdash Q\JoinOp{}S\neq P
  }[JoinIrreducibleProperLowerJoinProper]
    \proofstepwidestar[1]{\ZeroJoinSemi{L,\JoinOp,\varnothing}}{\rA}
    \proofstepwidestar[2]{\JoinIrr{L,\JoinOp,\varnothing}{P}}{\rA}
    \proofstepwidestar[3]{Q,S\in D_L(P)}{\rA}
    \proofstepwidestar[3]{Q\in D_L(P)}{\rAEa{3}}
    \proofstepwidestar[3]{S\in D_L(P)}{\rAEb{3}}
    \proofstepwidestar[3]{Q\in L}{%
      \FormulaRefAuto{JoinIrreducibleProperLowerCarrier}[thm]{4}}
    \proofstepwidestar[3]{S\in L}{%
      \FormulaRefAuto{JoinIrreducibleProperLowerCarrier}[thm]{5}}
    \proofstepwidestar[3]{%
      Q\in L\land Q\InducedLeq P\land Q\neq P}{%
      \FormulaRefAuto{JoinIrreducibleProperLowerSet}[def]{4}}
    \proofstepwidestar[3]{Q\neq P}{\rAEb{\rAEb{8}}}
    \proofstepwidestar[3]{%
      S\in L\land S\InducedLeq P\land S\neq P}{%
      \FormulaRefAuto{JoinIrreducibleProperLowerSet}[def]{5}}
    \proofstepwidestar[3]{S\neq P}{\rAEb{\rAEb{10}}}
    \proofstepwidestar[12]{Q\JoinOp{}S=P}{\rA}
    \proofstepwidestar[2,3,12]{Q=P\lor S=P}{%
      \FormulaRefAuto{JoinIrreducible}[def]{2,6,7,12}}
    \proofstepwidestar[14]{Q=P}{\rA}
    \proofstepwidestar[3,14]{\bot}{\rBI{14,9}}
    \proofstepwidestar[16]{S=P}{\rA}
    \proofstepwidestar[3,16]{\bot}{\rBI{16,11}}
    \proofstepwidestar[2,3,12]{\bot}{\rOE{13,14,15,16,17}}
    \proofstepwidestar[2,3]{Q\JoinOp{}S\neq P}{\rCI{12,18}}
  \closeproofpartwideindR

  \proofpartwide{Schluss}
    \proofstepwidestar[1]{\ZeroJoinSemi{L,\JoinOp,\varnothing}}{\rA}
    \proofstepwidestar[2]{\JoinIrr{L,\JoinOp,\varnothing}{P}}{\rA}
    \proofstepwidestar[3]{Q,S\in D_L(P)}{\rA}
    \proofstepwidestar[1,2,3]{%
      Q\JoinOp{}S\in L\land
      Q\JoinOp{}S\InducedLeq P}{%
      \FormulaRefAuto{JoinIrreducibleProperLowerJoinBound}[thm]{1,2,3}}
    \proofstepwidestar[1,2,3]{Q\JoinOp{}S\in L}{\rAEa{4}}
    \proofstepwidestar[1,2,3]{%
      Q\JoinOp{}S\InducedLeq P}{\rAEb{4}}
    \proofstepwidestar[1,2,3]{Q\JoinOp{}S\neq P}{%
      \FormulaRefAuto{JoinIrreducibleProperLowerJoinProper}[thm]{1,2,3}}
    \proofstepwidestar[1,2,3]{Q\JoinOp{}S\in D_L(P)}{%
      \FormulaRefAuto{JoinIrreducibleProperLowerSet}[def]{5,6,7}}
  \closeproofpartwide
\end{tabproofsplitwide}

\FormulaThmDeltaK[Größtes echtes unteres Element eines Irreduziblen]{%
  \Finite L\dsep\ZeroJoinSemi{L,\JoinOp,\varnothing}\dsep
  \JoinIrr{L,\JoinOp,\varnothing}{P}\vdash
  \exists P_\ast\in D_L(P)\,\forall Q\in D_L(P)\,
    Q\InducedLeq P_\ast
}{JoinIrreducibleLargestProperLower}{%
  \DeltaRow{Mengen}{L\dsep P\dsep P_\ast\dsep Q\dsep S}
  \DeltaRow{Funktionen}{\JoinOp}
}
\begin{tabproofsplitwide}
  \proofpartwideindR[Ein maximales echtes unteres Element existiert]{%
    \begin{aligned}[t]
      &\Finite L\dsep\ZeroJoinSemi{L,\JoinOp,\varnothing}\dsep
        \JoinIrr{L,\JoinOp,\varnothing}{P}\vdash{}\\
      &\exists P_\ast\in D_L(P)\,\forall Q\in D_L(P)\,
        (P_\ast\InducedLeq Q\to Q=P_\ast)
    \end{aligned}
  }{%
    \Finite L\dsep\ZeroJoinSemi{L,\JoinOp,\varnothing}\dsep
    \JoinIrr{L,\JoinOp,\varnothing}{P}\vdash
    \exists P_\ast\in D_L(P)\,\forall Q\in D_L(P)\,
    (P_\ast\InducedLeq Q\to Q=P_\ast)
  }[JoinIrreducibleProperLowerMaximalExists]
    \proofstepwidestar[1]{\Finite L}{\rA}
    \proofstepwidestar[2]{\ZeroJoinSemi{L,\JoinOp,\varnothing}}{\rA}
    \proofstepwidestar[3]{\JoinIrr{L,\JoinOp,\varnothing}{P}}{\rA}
    \proofstepwidestar[1,2,3]{%
      \Finite{D_L(P)}\land D_L(P)\neq\varnothing\land D_L(P)\subseteq L}{%
      \FormulaRefAuto{JoinIrreducibleProperLowerFinite}[thm]{1,2,3}}
    \proofstepwidestar[1,2,3]{\Finite{D_L(P)}}{\rAEa{4}}
    \proofstepwidestar[1,2,3]{D_L(P)\neq\varnothing}{\rAEa{\rAEb{4}}}
    \proofstepwidestar[1,2,3]{D_L(P)\subseteq L}{\rAEb{\rAEb{4}}}
    \proofstepwidestar[2]{\Semi{L,\JoinOp}}{%
      \FormulaRefAuto{ZeroJoinBaseSemilatticeAxiom}[ax]{2}}
    \proofstepwidestar[2]{\PartOrd{L,\InducedLeq}}{%
      \FormulaRefAuto{JoinSemilatticeInducesOrder}[thm]{8}}
    \proofstepwidestar[1,2,3]{%
      \exists P_\ast\in D_L(P)\,\forall Q\in D_L(P)\,
      (P_\ast\InducedLeq Q\to Q=P_\ast)}{%
      \FormulaRefAuto{FinitePosetMaximalElement}[thm]{9,5,6,7}}
  \closeproofpartwideindR

  \proofpartwideindR[Ein maximales echtes unteres Element ist das größte]{%
    \begin{aligned}[t]
      &\ZeroJoinSemi{L,\JoinOp,\varnothing}\dsep
        \JoinIrr{L,\JoinOp,\varnothing}{P}\dsep P_\ast\in D_L(P)\dsep{}\\
      &\forall S\in D_L(P)
        (P_\ast\InducedLeq S\to S=P_\ast)
        \dsep Q\in D_L(P)\vdash
        Q\InducedLeq P_\ast
    \end{aligned}
  }{%
    \ZeroJoinSemi{L,\JoinOp,\varnothing}\dsep
    \JoinIrr{L,\JoinOp,\varnothing}{P}\dsep P_\ast\in D_L(P)\dsep
    \forall S\in D_L(P)
    (P_\ast\InducedLeq S\to S=P_\ast)
    \dsep Q\in D_L(P)\vdash Q\InducedLeq P_\ast
  }[JoinIrreducibleProperLowerMaximalDominates]
    \proofstepwidestar[1]{\ZeroJoinSemi{L,\JoinOp,\varnothing}}{\rA}
    \proofstepwidestar[2]{\JoinIrr{L,\JoinOp,\varnothing}{P}}{\rA}
    \proofstepwidestar[3]{P_\ast\in D_L(P)}{\rA}
    \proofstepwidestar[4]{%
      \forall S\in D_L(P)
      (P_\ast\InducedLeq S\to S=P_\ast)}{\rA}
    \proofstepwidestar[5]{Q\in D_L(P)}{\rA}
    \proofstepwidestar[1]{\Semi{L,\JoinOp}}{%
      \FormulaRefAuto{ZeroJoinBaseSemilatticeAxiom}[ax]{1}}
    \proofstepwidestar[3]{P_\ast\in L}{%
      \FormulaRefAuto{JoinIrreducibleProperLowerCarrier}[thm]{3}}
    \proofstepwidestar[5]{Q\in L}{%
      \FormulaRefAuto{JoinIrreducibleProperLowerCarrier}[thm]{5}}
    \proofstepwidestar[1,2,3,5]{P_\ast\JoinOp{}Q\in D_L(P)}{%
      \FormulaRefAuto{JoinIrreducibleProperLowerJoinClosed}[thm]{1,2,3,5}}
    \proofstepwidestar[1,3,5]{%
      P_\ast\InducedLeq P_\ast\JoinOp{}Q}{%
      \FormulaRefAuto{JoinUpperLeft}[thm]{6,7,8}}
    \proofstepwidestar[1,2,3,4,5]{P_\ast\JoinOp{}Q=P_\ast}{%
      \rRE{\rRE{9,\rUE{4}},10}}
    \proofstepwidestar[1,3,5]{%
      Q\InducedLeq P_\ast\JoinOp{}Q}{%
      \FormulaRefAuto{JoinUpperRight}[thm]{6,7,8}}
    \proofstepwidestar[1,2,3,4,5]{%
      Q\InducedLeq P_\ast}{\rIE{11,12}}
  \closeproofpartwideindR

  \proofpartwide{Schluss}
    \proofstepwidestar[1]{\Finite L}{\rA}
    \proofstepwidestar[2]{\ZeroJoinSemi{L,\JoinOp,\varnothing}}{\rA}
    \proofstepwidestar[3]{\JoinIrr{L,\JoinOp,\varnothing}{P}}{\rA}
    \proofstepwidestar[1,2,3]{%
      \exists P_\ast\in D_L(P)\,\forall Q\in D_L(P)\,
      (P_\ast\InducedLeq Q\to Q=P_\ast)}{%
      \FormulaRefAuto{JoinIrreducibleProperLowerMaximalExists}[thm]{1,2,3}}
    \proofstepwidestar[5]{%
      P_\ast\in D_L(P)\land
      \forall S\in D_L(P)
      (P_\ast\InducedLeq S\to S=P_\ast)}{\rA}
    \proofstepwidestar[5]{P_\ast\in D_L(P)}{\rAEa{5}}
    \proofstepwidestar[5]{%
      \forall S\in D_L(P)
      (P_\ast\InducedLeq S\to S=P_\ast)}{\rAEb{5}}
    \proofstepwidestar[8]{Q\in D_L(P)}{\rA}
    \proofstepwidestar[2,3,5,8]{Q\InducedLeq P_\ast}{%
      \FormulaRefAuto{JoinIrreducibleProperLowerMaximalDominates}[thm]{%
        2,3,6,7,8}}
    \proofstepwidestar[1,2,3]{%
      \exists P_\ast\in D_L(P)\,\forall Q\in D_L(P)
      Q\InducedLeq P_\ast}{%
      \rEE{4,5,\rEI{\rAI{6,9}}}}
  \closeproofpartwide
\end{tabproofsplitwide}


\end{document}
